\documentclass[12pt,a4paper]{article}

\usepackage{mathtools}
\usepackage{mathrsfs}
\usepackage{float}
\usepackage{caption}
\usepackage{tikz}
\usepackage{tikz-cd}
\usetikzlibrary{arrows,positioning,calc,decorations.pathmorphing}

%% Macros-v2_198: current project v3.072 III-D and bibliography binding; mathematical macros unchanged.
%% Current project v3.070 source-reference bindings; previous macro bodies unchanged.
% Current II-D reference root for PapersMinimal v3.063; mathematical macros otherwise unchanged.
% Macros-v2_188: re-based current paper identities and A--G family; no change to geometry.
% Macros-v2_184: active Paper III-E v1.200/III-F v1.267 binding for framework v3.031; mathematical macros unchanged.
% Macros-v2_183: updated shared cross-reference bindings for Paper1B/III-E/III-F v3.030; mathematical macros unchanged.
% Macros-v2_182: active Paper II-F/V-A cross-references for framework v3.029; definitions unchanged.
% Paper-series v2.381: new I-B torsion/modular results, bibliography and current source roots.
% Paper-series v2.380: new exceptional bibliography and synchronized paper references.
% Paper-series v2.379 current 31-paper macro/aux registry.
% Paper-series v2.378: current 31-paper AUX/bibliography registry updated for Paper I-B v1.2 literature additions.
% Framework v2.371: current-source reference registry refreshed; old theorem labels retained.
% Framework v2.370: current 30-paper source/AUX root rebinding; native mathematical macros unchanged.
% Framework v2.368: synchronized merged 30-root registry; mathematical notation unchanged.
% Framework v2.358: active 30-paper source registry synchronized; existing mathematical macros unchanged.
% Framework v2.357: selective new source-owned certificates; active reference roots updated for five revised papers.
% Framework v2.356: active 30-root registry synchronized to the 30 papers in this release.
% Framework v2.275: exact positive dense-twin word cubature, finite repeatability and eighth-moment classification, nested sampling frames, joint action and POVM-jet exports; preserves v2.274.
% Framework v2.274: filename and root-source synchronization.
% Framework v2.273: restored missing Problem 3 theorems from v2.268; exact exceptional frames, twin geometry, convolution repeatability and projector/occurrence firewalls; merged on v2.272 without overwriting independent branches.
% Framework v2.272: merged v2.270 covariant native-readout mathematics with v2.271 Problem 4 finite-action Hopf hypercharge/holonomy corrections; no universal bulk-gauge uniqueness claim.
% Independent branches: v2.270 readout/guard exports and v2.270 K16 action-consumed charge closure; all prior scoped assumptions retained.
% Framework v2.269: synchronized 30 active AUX roots; mathematics and bibliography preserved from Macros-v2_161.
% Framework v2.266: Paper II-C converse-derived outcome resources; II-E, III-C, IV-A/B, V-A and VI-B downstream mathematical closures; 30-root registry synchronized.
% Framework v2.265: II-C coordinate-stage constitution and exact first-symbol host inheritance; prior mathematical domains retained.
% Framework v2.264: paper II-B reverse/audit and IV-A observable-minimal exchange; synchronized external references to 30 active roots.
% Framework v2.263: Paper I-C adds bridge-controlled root-carry parity and source-relative old-cap cancellation; other root mathematical bodies retained; active 30-root references synchronized.
% Framework v2.262: integrates Paper II-F radius, Casimir and readout results; active 30-root references synchronized.
% Framework v2.261: merged Paper II-B GP2-level/geometric repairs; preserves v2.260 II-D repairs; active references synchronized.
% Framework v2.260: Paper II-D signed RS, finite-symbol typing, and conditional curvature/Jordan repairs; references synchronized.
% Framework v2.259: conditional native SU(2) atlas and explicit later discharge; v2.258 mathematical bodies retained elsewhere.
% Framework v2.258: preserves parallel II-A/dynamics update and adds complete readout/frequency closures.
% Framework v2.257: exact minimal complete position readout and arbitrary-preparation frequency law.
% Framework v2.256: exact native inertia and sharp bounded ordered-source closure; preserves concurrent v2.255 work.
% Framework v2.255: preserves parallel detector integration and adds exact finite readout closures.
% Framework v2.254: Pasterski bibliography and synchronized 30-paper external-reference registry.
% Framework v2.253: typed raw coordinate defect in I-A and synchronized VI-A import; active references synchronized.
% Framework v2.252: synchronized source registry; optional III-G shell-frame calculation.
% Framework v2.251: forward common-action execution, joint source jets and scoped coupled control.
% Framework v2.251: exact full-heat variations, fermion source jets and explicit physical completion frontier.
% Framework v2.251: certified interacting source jets, pole stability and scoped completion proofs.
% Framework v2.251: synchronized active source registry; mathematical notation unchanged.
% Framework v2.251: synchronized active source registry; mathematical notation unchanged.
% Framework v2.251: synchronized active source registry; mathematical notation unchanged.
% Framework v2.251: synchronized source registry; no mathematical notation changes.
% Framework v2.244: synchronized 30-paper AUX registry and rooted phase-source action-order tests.
% Framework v2.242: signed rooted phase lift, complete transverse return and exact family cubic.
% Framework v2.241: proved rooted conormal geometry, conditional cap realization and scoped action tests.
% Framework v2.239: all 30 active source roots synchronized; no mathematical macro changes.
% Framework v2.238: Paper IV A/B/C split and 30-root reference registry.
% Framework v2.237: Paper VII-D technical split and 28-root reference registry.
% Framework v2.236: Paper VI-C technical-volume reference added; all current root bindings advanced.
% Framework v2.235: 26-root AUX and bibliography bindings synchronized; mathematical macros unchanged.
% BEGIN INLINE SOURCE: Paper00-TeXBuildCompatibility-v1_0.tex
% Rendering-only repairs for inherited long appendices and bitmap text fonts.
% Integer counters through 26 retain the original labels; later values use AA,
% AB, ... . This file changes no action coefficient, state or theorem premise.
\RequirePackage{alphalph}
\makeatletter
\let\BHPrimitiveAlph\@Alph
\newalphalph{\BHExtendedAlph}{\BHPrimitiveAlph}{26}
\let\@Alph\BHExtendedAlph
\AtBeginDocument{%
  \@ifpackageloaded{microtype}{\microtypesetup{expansion=false}}{}%
}
\makeatother

% END INLINE SOURCE: Paper00-TeXBuildCompatibility-v1_0.tex

% Framework v2.164: current reference/bibliography profile; mathematical definitions retained.
% Framework v2.162: finite metric, functional Hilbert and scoped Ward integration.
% Framework v2.161: inherits Macros-v2_57.tex exactly except the current external-reference binding.
% Macros v2.57 / framework v2.156: current reference registry only; mathematical definitions unchanged.
% Framework v2.143 / Macros v2.55: current external-reference bindings; mathematical notation unchanged.
% Framework v2.110 / Macros v2.54: current external-reference registry and bibliography; mathematical notation inherited.
% Framework v2.67 / Macros-v2_53: adds Born-action-cell and joint-EH readout symbols; absolute unit realization is metrological; MATCH_{J/scal}^{EH} remains the sole G24 normalization bridge.
% Framework v2.35 / Macros-v2_52: grouped paper numbering, active AUX registry, and Born-v6_20 bibliography binding synchronized; notation unchanged.
% Framework v2.26: active-source AUX and bibliography bindings synchronized; notation unchanged.
% Framework v2.09: reviewed early-reciprocity integration; scoped physical imports.
% Framework v2.08: Born--Euler mathematics and conditional complete-action potential criterion.
% Framework v2.07: active AUX registry includes Paper I-B v1.2; shared bibliography Born-v6_18.
% Macros-v2_46 / framework v2.04: mathematical atlas and downstream realization correction.
% Macros-v2_46 / framework v2.03: scoped structure integration and current dependency bindings.
% Macros-v2_46 / framework v2.01: shared source bindings; prior revision comments are historical.
% v2.42 / framework v1.99: Synchronized RG/AS owner and bibliography bindings.
% v2.41 / framework v1.99: active manuscript AUX bindings and geometric-closure bibliography synchronized.
% Release-v1.80 centralized source synchronization for scalar-circle/comparison-band integration.
% Release v1.79: Hurwitz hierarchy and kinetic-normalized action response integration; physical mass matching remains conditional.
% Release v1.78: approved M8 closure-batch integration; physical and analytic domains retained.
% Current v1.78 registry: all active TeX roots are synchronized below.
% Earlier release comments are historical scope records, not current status.
% Release-v1.77 synchronization: native same-owner Noether closure, generator--flow/history typing, analytic-certification firewalls, scalaron/K16, LF/QCD, quantum-Ward, and RAR boundary-owner transport are registered; shared mathematical notation is unchanged.
% Release-v1.76 synchronization: registry points to the v1.76 active paper roots.
% Release-v1.76 synchronization: registers the local A+Phi+g Noether/constraint closure and dynamics-frontier propagation; shared mathematical notation is unchanged.
% Release-v1.75 synchronization: registry points to the v1.75 active paper roots.
% Release-v1.75 synchronization: integrates the joint-Dirac Einstein bulk real form and the local odd-scalar heat-coefficient action on explicit analytic first-jet/interior domains; v1.74 Yang--Mills closure is retained. Fermionic reality, pure-metric curvature-squared/exact-nonlocal terms, noncanonical boundaries, physical vacuum selection, and actual solution-history tokens remain open.
% Release-v1.74 synchronization: registry points to the v1.74 active paper roots and retains the v1.73 theorem labels.
% Release-v1.74 synchronization: canonical Macros-v2_34.tex; generated local nonlinear Yang--Mills action, constrained connection-field phase space, source-free Gauss momentum map, and Legendre closure are integrated on their proved domains. The remaining non-gauge ACT_cl sectors and actual solution-history selection remain open.
% Release-v1.71 synchronization: closes the former JOINT_x/JOINT_5
% connection-occurrence gap by canonical structure-group extension of the
% nonflat Born connection, faithful associated representation, and the local
% cotangent--Clifford operator constructor.  Actual fluctuation members,
% actions, vacua, solutions, calibrations, and analytic QFT domains remain
% separately typed.  The active-AUX registry is advanced once.
% Release-v1.71 canonical macro and active-AUX registry.
% ======================================================================
% Macros-v2_39.tex -- canonical publication macro file
% Release v1.71 retains the v1.70 null-ray/QCD boundaries and additionally
% removes the duplicate structural connection actuality premise.  It does not
% claim global bundle topology, select a gauge/Higgs member, or promote the
% finite construction to an interacting continuum theory.
% Release v1.68 constructs the finite light-front total-fraction operators and
% parton number/momentum operator measures, proves their simplex and Hausdorff-limit
% theorems, and separates positive cut measures from scheme-renormalized
% light-ray distributions.  It also separates finite real Schur poles, stable
% norm-resolvent/Riesz pole-residue limits, and continued resonance
% pole/residue limits.  The Paper-VII-C scattering package is claim-indexed:
% local QFT and spectral claims import no S-matrix or infrared premise, while
% asymptotic amplitudes, LSZ, S-matrix poles, completeness, and cross sections
% import only their named massive or massless/long-range clauses.
% Release v1.67 factorizes BND48_H into a same-action Hessian/Schur
% descendant and still-open typed boundary-trace, critical-lift, and
% refinement-limit conditions.  It proves the finite-collar Dirichlet-to-
% Neumann law sqrt(y)tanh(L sqrt(y)), its nonuniform critical limit, and the
% finite-return ultra-infrared turnover.  The Paper-I-A Hardy constant mode is
% an exact geometric near-hit, not yet the physical metric boundary mode.
% The current local f(R) Hessian remains massive and therefore still refutes
% a RAR prediction on the minimal action branch.
% Release v1.66 derives the transverse depth-six normalizer bridge from the
% already generated C6 Reynolds image and Paper-I-A binary-dihedral normalizer.
% The bridge is the unique least 2D3-stable scalar-containing realization of
% the faithful quotient 2D3/C6=C2, has trace defect six and hence K=48, but is
% transverse to the nested active carrier.  Exact K=44 and K=52 contraction,
% semigroup, and energy-midpoint identities are registered as structural
% readouts rather than projector events.  The active E48 carrier gap remains
% unchanged.  A same-action boundary-symbol/Calderon intertwiner is isolated
% as the necessary and still unproved owner bridge to the static DtN sort;
% representation equivalence alone neither derives RAR nor creates a K=48
% action junction, pole, occupation, or cosmological transition.
% Release v1.65 derives the native E_48-centered carrier mass gap from the
% GP1--GP3-generated active spectrum, distinguishes its carrier-free interval
% from a particle located at E_48, and withdraws the unsupported promotion to
% a positive continuum Yang--Mills mass gap.  The latter remains a valid
% conditional transfer theorem on an enlarged continuum branch, but it is not
% an active framework prediction or an unfinished native-gap proof.
% Release v1.64 synchronizes Derivation Batches 14--22.  It registers the
% GP1--GP3 heat-profile orbit, same-Hessian electroweak response and pole
% ownership, K16/K48 action-response corrections, perturbation state and
% non-Gaussianity typing, finite-vacuum/gap/QCD limit repairs, factorized
% position-area ownership, the actual static metric response, and the ordered
% Wilson/string-tension limits.  Strict continuum convergence, positive
% physical gap, surviving higher cumulants, and positive tension remain open.
% Release v1.61 closes TIMEWIDTH16 as an independent direct-production datum.
% JUNC_{2|7} continues to own the finite intrinsic microscopic profile/width,
% while the actual same-action background/production history determines its
% physical-time pullback, support, and duration.  The current Paper-I-A--VII-C AUX
% registry is advanced once; shared mathematical notation is unchanged.
% Release v1.60 derives the final-octet primitive-remnant depletion, canonical
% Hopf-current representative, exact Hopf descent to dOmega/(4pi), and
% closed-FLRW a^{-3} neutral-current shape.  It retypes K=48 as the onset
% of the final one-unit geometric transfer and records cDM/RAR as distinct
% physical projections of that common neutral parent while preserving
% occurrence, Hilbert-owner, stress-matching, and static-kernel premises.
% Release v1.59 transports the v1.58 exceptional-shell/K16 closure through
% Papers III-F and VI-B--VII-C.  It replaces Paper VI-B's obsolete mixed-Hessian
% output typing by the third-variation pair vertex and Bogoliubov occupation
% support, derives the post-seed Higgs-pair feedback vertex and its exact
% alpha--delta kernel factors, propagates the unique K16/K48 action-junction
% split, identifies the QCD delta count with the upstream 2O shell atom only
% as a common-scalar numerical compatibility, and narrows THERM16 to dynamical
% approach/beta selection because the finite Gibbs/KMS target is already
% unique on a supplied simple observable algebra once its Hamiltonian and
% inverse temperature are supplied; direct-sum central weights remain typed.
% Release v1.58 registers the exceptional nearest-shell closure, the derived
% transverse K=16 action event, the unique interior action junction, the
% type-correct ordered pair vertex, the thin-layer squeeze law, and the
% radiation-like/matter-like kinematic split.  REG_0, transfer magnitude,
% physical layer width, and thermalization remain explicitly conditional.
% Release v1.57 registers the active-depth-six obstruction, the native
% K=48 carrier desert, the typed 7|9 owner reduction, and the dynamical
% reheating handoff across Papers III-B--VI-B.
% Release v1.56 registers the least-support n_2 inflation dressing of the
% persistent n_7 Laurent response, the analytic global-convexity and positive
% endpoint-measure proofs, the derived inflation core and K=16 coefficient
% freeze, and the no-particle dark-sector split into static RAR plus the
% conditional homogeneous effective-curvature response at reflected K=48.
% There is no n_6 action, scalaron-DM mass, or second Einstein--Hilbert term.
% Release v1.55 registers the protected-octet resolved/residual reflection
% q_sp+Z_W=8, its j<->8-j and K<->64-K scale covariance, and the derived
% rank-one-exit quadratic grade C_1^2/U=C_2 in Papers III-C, III-E, and VIII.
% The reflection is a typed trace/scale theorem, not a time reversal or a
% physical event-occurrence theorem; reheating transfer remains conditional.
% Release v1.54 registers the exact positive-part Laurent indices, the
% active-domain/terminal-endpoint action split, the no-second-EH firewall, and
% the q_DM=6/K_DM=48 scalaron readout typing across Papers III-D/E/F, VI,
% VIII, and IX.
% Release v1.53 registers the common nonnegative Laurent obstacle projection
% for n_1^+ and n_7^+, discharges the three former exit/freeze assumptions,
% and advances the Paper-III-D/E/F, VIII, and IX consumers.  Shared notation
% and bibliography ownership are unchanged.
% Release v1.51 registers the PQTI derivation batch and its updated
% Paper-III-A/C/D/E/F, VIII, and IX consumers.  Shared notation and
% bibliography ownership are unchanged.
% Release v1.50 advances the Paper-III-B AUX target to v1.40 after the
% rigorous audit repair; shared notation and bibliography ownership are unchanged.
% Release v1.49 advances the Paper-II-E/F AUX target to v3.14 after the
% minimal audit repair; shared notation and bibliography ownership are unchanged.
% Release v1.46 synchronizes Paper II-A's minimal frozen constructor-origin
% theorem and its five dependent AUX targets.  Shared notation is unchanged.
% Release v1.45 updates the Paper-II-B AUX target to v2.23 after the
% GP3-rigidity and appendix-placement repair; shared notation is unchanged.
% Release v1.44 advances the centralized Paper-I-A--VII-C AUX registry after the
% Paper-I-A v5.16 protected-completion scope repair.  Every current
% paper imports this one file; no paper-level opt-in or \externaldocument list
% is permitted.  The bibliography remains centralized here through Born-v6_13.
% Shared bibliography, theorem infrastructure, colors, notation alignment,
% and current Paper-I-A--VII-C external-reference registration.
% This file centralizes common imports, bibliography sorting, colors, theorem
% infrastructure, and the notation shared by the complete series.
% v2.9 (release v1.43) made the Paper-I-A--VII-C AUX registry automatic and
% centralized the bibliography and all external-reference ownership.
% v2.8 release-registry refresh (v1.33) registers the rooted generated-parent
% cascade and the POL/REG/free-law/11+4/Born-group closure revisions.  No
% shared mathematical notation or Paper-II-A fundamental principle is changed.
% v2.8 release-registry refresh (v1.32) registers the universal generated-
% continuation revisions in Papers II-A, III-A, III-C, IV, and VIII.  No
% shared mathematical notation or Paper-II-A fundamental principle is changed.
% v2.8 release-registry refresh (v1.30) registers the CRP4/projective-state
% source revisions and their direct registry-only consumers.  No shared
% mathematical notation or Paper-II-A axiom is changed.
% v2.7 release-registry refresh (v1.27) updates only the current external
% Paper-I-A--VII-C basenames after the placement/assumption synchronization; no
% shared mathematical notation or Paper-II-A axiom is changed.
% v2.7 synchronizes the geometric defect-repair revision: Paper-II-A v3.42,
% Paper-III-C v1.26, Paper-III-D v1.27, Paper-III-E v1.32, and Paper-III-F
% v1.34.  It registers the common trace-defect/spectral-filtration theorem,
% its owner-typed response pullback, and the reduced GRED and TERM interfaces.
% No shared mathematical notation is changed.
% v2.6 synchronizes Paper-III-E v1.31's cone-type repair, minimal-response
% classification, and explicit classical action-role boundary, together with
% Paper-III-F v1.32's downstream response-atlas synchronization.  No shared
% mathematical notation is changed.
% v2.5 synchronizes the Paper-III-D v1.26 faithful-trace-duality proof,
% exact ordered-owner joint locus, and typed alternative-presentation
% bookkeeping repair, together with its Paper-III-E and closure-frontier
% references.  No shared mathematical notation is changed.
% v2.4 synchronizes the Paper-III-C v1.25 orbit/process/owner repair, its
% downstream Papers III-B--F, V, VI, VIII, and IX revisions, and Born-v6_13.
% v2.3 synchronizes the Paper-II-A v3.41 witness/support/principal inheritance
% layer, the Paper-III-A v1.24 joint-realizability/readout firewall, all dependent
% host-rate revisions, and the current Paper-I-A--VII-B source registry.
% v2.2 synchronizes the generated active-realization, determinant-depth,
% dimension-weighted hierarchy, and typed cosmological-event interfaces.
% v2.1 synchronizes the exceptional-support interface revisions: Paper I-A
% v5.9, Paper II-A v3.30, Paper II-E/F v3.7, and Paper IV v4.16.
% v2.0 introduces the stage-safe G0 occurrence notation, synchronizes the
% active Paper-I-A--IV source registry (current rigorous Paper-III closure:
% III-A v1.17, III-B v1.22, III-C--E v1.17, III-F v1.14, IV v4.14), and makes
% external-reference self-skip depend on logical paper role rather than a
% versioned job name.
% ======================================================================

% ======================================================================
% macros.tex — Shared macros for the Born Geometry series
% ======================================================================

% ======================================================================
% 1. PACKAGES, BIBLIOGRAPHY, COLORS, AND UNIT DEFINITIONS
% ======================================================================
% Core packages shared across the complete Paper-I-A--VII-C series.  Loading them here makes
% the macro file a genuine series preamble while remaining compatible with existing
% paper preambles that load the same packages before \input.
\usepackage{amsmath}
\usepackage{amssymb}
\usepackage{amsthm}
\usepackage{mathrsfs}
\usepackage{array}
\usepackage{booktabs}
\usepackage{xcolor}
\makeatletter
\@ifpackageloaded{geometry}{}{\usepackage{geometry}\geometry{margin=1in}}
\makeatother
\usepackage{mdframed}

% Shared support packages required by the macros below and by most papers.
\usepackage{ragged2e}
\usepackage{longtable}
\usepackage{tabularx}
\usepackage{enumitem}
\usepackage[most]{tcolorbox}

% Series-wide numeric bibliography convention using the portable BibTeX path.
% All Papers I-A--VII-C use the ordinary \cite interface.
\usepackage[numbers,sort&compress]{natbib}
\providecommand{\printbibliography}{%
  \bibliographystyle{plainnat}%
  \bibliography{Born-v6_60}%
}

% Shared colors.  \providecolor avoids clobbering a paper-specific override.
\providecolor{postulatecolor}{RGB}{230,242,255}
\providecolor{compare}{RGB}{230,245,230}
\providecolor{definitioncolor}{RGB}{230,242,255}

\providecolor{openriskcolor}{RGB}{255,235,235}
\providecolor{bridgecolor}{RGB}{240,240,255}
\providecolor{gold}{RGB}{255,215,0}

% ======================================================================
% 2. CUSTOM ENVIRONMENTS
% ======================================================================
% Common table column types used in all papers.  Guarded so that a paper may
% define its own versions before loading this file.
\makeatletter
\@ifundefined{NC@find@L}{%
  \newcolumntype{L}[1]{>{\RaggedRight\arraybackslash}p{#1}}%
}{}
\@ifundefined{NC@find@C}{%
  \newcolumntype{C}[1]{>{\Centering\arraybackslash}p{#1}}%
}{}
\@ifundefined{NC@find@Y}{%
  \newcolumntype{Y}{>{\RaggedRight\arraybackslash}X}%
}{}
\makeatother

% ======================================================================
% 3. THEOREM STYLES AND ENVIRONMENTS
% ======================================================================
% Custom underlined theorem style (useful for remarks or notes)
\newtheoremstyle{underlined}
  {3pt}   % Space above
  {3pt}   % Space below
  {}      % Body font (empty means normal/roman text, typical for remarks)
  {}      % Indent amount
  {}      % Theorem head font (can be \itshape, \bfseries, etc.)
  {.}     % Punctuation after theorem head
  { }     % Space after theorem head
  {\underline{\thmname{#1}\thmnumber{ #2}}\thmnote{ (#3)}} % Theorem head formatting

% Standard plain theorem style environments.  Each environment is guarded so
% the shared macro file can be loaded after a paper-specific theorem setup.
\theoremstyle{plain}
\makeatletter
\@ifundefined{theorem}{\newtheorem{theorem}{Theorem}[section]}{}
\@ifundefined{lemma}{\newtheorem{lemma}{Lemma}[section]}{}
\@ifundefined{proposition}{\newtheorem{proposition}{Proposition}[section]}{}
\@ifundefined{corollary}{\newtheorem{corollary}{Corollary}[section]}{}
\makeatother

% Definition style environments (typically non-italicized body)
\theoremstyle{definition}
\makeatletter
\@ifundefined{definition}{\newtheorem{definition}{Definition}[section]}{}
\@ifundefined{postulate}{\newtheorem{postulate}{Postulate}[section]}{}
\@ifundefined{closure}{\newtheorem{closure}{Closure}[section]}{}
\makeatother

% Applying the custom underlined style
\theoremstyle{underlined}
\makeatletter
\@ifundefined{remark}{\newtheorem{remark}{Remark}[section]}{}
\@ifundefined{convention}{\newtheorem{convention}{Convention}[section]}{}
\@ifundefined{principle}{\newtheorem{principle}{Principle}[section]}{}
\makeatother

% Framed proof environment.
% Every \begin{proof} ... \end{proof} is automatically framed/shaded.
% This avoids mdframed/trivlist nesting, which can drop the top border after
% theorem-like environments.
\makeatletter
\renewenvironment{proof}[1][\proofname]{%
  \par\addvspace{6pt}%
  \noindent
  \begin{tcolorbox}[%
    enhanced,
    breakable,
    colback=black!4,
    colframe=black!35,
    boxrule=0.35pt,
    arc=1pt,
    left=6pt,
    right=6pt,
    top=4pt,
    bottom=4pt,
    boxsep=0pt,
    before skip=0pt,
    after skip=6pt,
    parbox=false
  ]%
  \pushQED{\qed}%
  \normalfont\itshape #1\@addpunct{.}\enspace
  \normalfont\ignorespaces
}{%
  \unskip\nobreak\hfill\popQED
  \end{tcolorbox}%
  \par\addvspace{2pt}%
}
\makeatother

\mdfdefinestyle{postulatebox}{
  backgroundcolor=yellow!10,
  linecolor=black,
  linewidth=0.4pt,
  skipabove=\baselineskip,
  skipbelow=\baselineskip,
  nobreak=true
}

% ======================================================================
% PROOF-READABILITY HELPERS
% ======================================================================

\newcommand{\ProofPlan}[1]{%
  \par\smallskip
  \noindent\textbf{Proof plan. }#1%
  \par\smallskip
}

\newlist{proofsteps}{enumerate}{2}
\setlist[proofsteps,1]{
  label=\textbf{(\alph*)},
  ref=\alph*,
  leftmargin=2.2em,
  itemsep=0.45\baselineskip,
  topsep=0.35\baselineskip
}
\setlist[proofsteps,2]{
  label=\textbf{(\roman*)},
  ref=\roman*,
  leftmargin=2.0em,
  itemsep=0.25\baselineskip,
  topsep=0.2\baselineskip
}


% ======================================================================
% 4. NUMBER SYSTEMS AND ALPHABETS
% ======================================================================
\providecommand{\RR}{\mathbb{R}} % Real numbers
\providecommand{\PP}{\mathbb{P}} % Projective space / Probability
\providecommand{\CC}{\mathbb{C}} % Complex numbers
\providecommand{\ZZ}{\mathbb{Z}} % Integers
\providecommand{\NN}{\mathbb{N}} % Natural numbers
\providecommand{\HH}{\mathbb{H}} % Quaternions / Hyperbolic space
\providecommand{\LFRayGeom}{\mathsf{LF\!\mbox{-}RAY\!\mbox{-}GEOM}}
\providecommand{\LFRep}{\mathsf{LF\!\mbox{-}REP}}
\providecommand{\LFState}{\mathsf{LF\!\mbox{-}STATE}}
\providecommand{\LFLim}{\mathsf{LF\!\mbox{-}LIM}}
\providecommand{\LFScat}{\mathsf{LF\!\mbox{-}SCAT}}
% Number systems / projective spaces
\providecommand{\CP}{\mathbb{CP}}
\providecommand{\RP}{\mathbb{RP}}

% ======================================================================
% 5. GROUPS AND LIE THEORY
% ======================================================================
\providecommand{\SO}{\ensuremath{\mathrm{SO}}}     % Special Orthogonal group
\providecommand{\SU}{\ensuremath{\mathrm{SU}}}     % Special Unitary group
\providecommand{\Spin}{\ensuremath{\mathrm{Spin}}} % Spin group
\providecommand{\CS}{\ensuremath{\mathrm{CS}}}     % Chern-Simons (capitalized)
\providecommand{\cs}{\ensuremath{\mathrm{cs}}}     % Chern-Simons (lowercase)

% ======================================================================
% 6. COMMON MANIFOLDS AND TOPOLOGY
% ======================================================================
\providecommand{\Sone}{S^1}                 % Standard circle

% ======================================================================
% 7. MATH OPERATORS (Declared via \DeclareMathOperator)
% ======================================================================
\DeclareMathOperator{\End}{End}   % Endomorphisms
\DeclareMathOperator{\Aut}{Aut}   % Automorphisms
\DeclareMathOperator{\spec}{spec} % Lowercase spectrum
\DeclareMathOperator{\tr}{tr}     % Lowercase trace (if you already have \Tr)
\DeclareMathOperator{\supp}{supp} % Support
\DeclareMathOperator{\TrA}{Tr_{A}} % Partial trace over subsystem A
\DeclareMathOperator{\TrB}{Tr_{B}} % Partial trace over subsystem B
% Lie/algebra operators
\DeclareMathOperator{\Ad}{Ad}
\DeclareMathOperator{\Sym}{Sym}
\DeclareMathOperator{\su}{su}
% ======================================================================
% 8. STANDARD OPERATOR-LIKE SYMBOLS (Safe definitions via \providecommand)
% ======================================================================
\providecommand{\vol}{\operatorname{vol}}   % Volume form
\providecommand{\Vol}{\operatorname{Vol}}   % Total volume
\providecommand{\Ric}{\operatorname{Ric}}   % Ricci tensor
\providecommand{\Rm}{\operatorname{Rm}}     % Riemann tensor
\providecommand{\Tr}{\operatorname{Tr}}     % Uppercase trace
\providecommand{\Spec}{\operatorname{Spec}} % Uppercase spectrum
\providecommand{\diag}{\operatorname{diag}} % Diagonal matrix
\providecommand{\rank}{\operatorname{rank}} % Rank
\providecommand{\scal}{\mathcal R} % Scalar curvature
\providecommand{\Cl}{\operatorname{Cl}}     % Clifford algebra
\providecommand{\diam}{\operatorname{diam}} % Diameter

% ======================================================================
% 9. FUNCTIONS, TRIGONOMETRY, AND CALCULUS
% ======================================================================

% Codifferential naming conventions
% (Avoid \ddag: commonly defined by physics.sty and symbol packages)
\providecommand{\codiff}{\mathrm{d}^{\dagger}} % Codifferential operator

% Quantum mechanics bra-ket notation (needed for basis tables)
\providecommand{\ket}[1]{\lvert #1 \rangle}
\providecommand{\bra}[1]{\langle #1 \rvert}
\providecommand{\braket}[2]{\langle #1 \mid #2 \rangle}

% ======================================================================
% 10. PHYSICAL SCALES AND CONSTANTS (Paper 1c / Born Geometry)
% ======================================================================
\providecommand{\Rp}{\ensuremath{R_{\mathrm P}}}     % matched common EH curvature scale (not defined from observed G)
\providecommand{\hBorn}{\ensuremath{h_{\mathrm B}}} % physical representative 2\pi\hbar of one primitive Born scalar action cell
\providecommand{\UEHjoint}{\ensuremath{U_{J}}} % joint-Dirac EH curvature-coordinate scale 8 kappa_J
\providecommand{\Gread}{\ensuremath{G_{\rm scal}}} % scalar-action coordinate c_op^3/(h_B U); physical matching explicit
\providecommand{\GreadJoint}{\ensuremath{G_{J,\mathrm{raw}}}} % raw ordinary-trace coordinate = G_scal/96
\providecommand{\Rds}{R_{\mathrm{dS}}}    % de Sitter Curvature
\providecommand{\Mpl}{M_{\mathrm{Pl}}}    % Planck mass
\providecommand{\lp}{\ell_{\mathrm{Pl}}}  % Planck length
\providecommand{\Pmax}{P_{\mathrm{max}}}  % Maximum momentum
\providecommand{\Mplbar}{\bar{M}_{\mathrm{Pl}}} % Reduced Planck mass
\providecommand{\eff}{\mathrm{eff}}       % Subscript for effective quantities
\providecommand{\Cosmo}{\ensuremath{\Lambda}}                % Cosmological constant (use \Cosmo in source)
\providecommand{\LCDM}{\ensuremath{\Cosmo\mathrm{CDM}}}  % \LambdaCDM (typesets with \Cosmo)
\providecommand{\Lcut}{\ensuremath{\Lambda_{\mathrm{cut}}}}  % Cutoff scale (brace when used as script)
\providecommand{\Rcs}{\ensuremath{R_{\mathrm{cs}}}}      % Configuration-space radius
\providecommand{\CosmoEff}{\ensuremath{\Cosmo_{\eff}}}   % Effective cosmological constant

% Disambiguation: \Rmom is the momentum-space (S^3) radius. 
% Reserve plain R for the 4D curvature scalar in f(R).
\providecommand{\Rmom}{\ensuremath{R_{\mathrm{mom}}}} 
\providecommand{\Punit}{P_{0}} % Momentum unit set by \Rmom (= \hbar/\Rmom)
\providecommand{\Rvec}{\ensuremath{R_{\mathrm{vec}}}} % invariant vector-space radius in the reduced round sector
% Specific numeric values and expressions
\providecommand{\Vlogcap}{L_{\mathrm{Hopf}}^{(1)}}
\providecommand{\Hopf}{\operatorname{Hopf}}
% Scalar phase (S^1) sector + universal cover notation (Paper I-A)
\providecommand{\tper}{T_\star}      % Time period/scale
\providecommand{\Eper}{E_\star}      % Energy period/scale
\providecommand{\ttilde}{\tilde t}   % Universal cover time
\providecommand{\Etilde}{\tilde E}   % Universal cover energy

% ======================================================================
% 10A. SERIES-WIDE LOGICAL / NOTATIONAL DISAMBIGUATIONS
% ======================================================================
% Use these symbols to avoid collisions among: (i) the Paper--III
% no-insertion/geometric-closure class, (ii) Paper--I-A trimer diagonals,
% (iii) Paper--V-A/V-C family-space diagonals, (iv) scalar-index running,
% and (v) protected shell-clock labels.
\providecommand{\Rprim}{\ensuremath{\mathsf R_{\mathrm{prim}}^{\Sigma}}} % primitive finite structural readouts
\providecommand{\Rloop}{\ensuremath{\mathsf R_{\circ}^{\Sigma}}} % closed-loop readout role
\providecommand{\Rbdry}{\ensuremath{\mathsf R_{\partial}^{\Sigma}}} % boundary-conditioned readout role
\providecommand{\Rgen}{\ensuremath{\mathsf R_{\mathrm{fin}}^{\Sigma}}} % finite generated closure of primitive roles
\providecommand{\PhysOcc}{\operatorname{PhysOcc}}
\providecommand{\Actual}{\operatorname{Actual}}
\providecommand{\GeomLic}{\operatorname{GeomLic}}
\providecommand{\InvDist}{\operatorname{InvDist}}
\providecommand{\Dcl}{\ensuremath{\mathsf{D}_1}}      % unmarked positive Dirac-even closed spectral-readout form class
\providecommand{\Dheat}{\ensuremath{\mathsf D_{\mathrm{heat}}}} % one-parameter heat-semigroup subclass
\providecommand{\DtriI}{\ensuremath{D^{\mathrm{tri}}_1}} % Paper--I-A cyclic-trimer diagonal
\providecommand{\DtriII}{\ensuremath{D^{\mathrm{tri}}_2}}
\providecommand{\DfamI}{\ensuremath{D^{\mathrm{fam}}_1}} % Family-space diagonal selector
\providecommand{\DfamII}{\ensuremath{D^{\mathrm{fam}}_2}}
\providecommand{\alpharun}{\ensuremath{\alpha_{\mathrm{run}}}} % scalar spectral-index running
\providecommand{\Ksh}{\ensuremath{K_{\mathrm{sh}}}} % protected shell-clock coordinate
\providecommand{\ddef}{\ensuremath{\delta_{\mathrm{def}}}} % universal defect scalar after Paper--III realization

% ======================================================================
% 10B. SERIES-WIDE LOGICAL STATUS TAGS
% ======================================================================
% Upright small-caps wrapper for audit/status labels.  It prevents LaTeX from
% trying to use unavailable small-caps italic font shapes inside theorem bodies
% such as TU/lmr/m/scit.
\providecommand{\AuditLabel}[1]{}
\renewcommand{\AuditLabel}[1]{{\normalfont\scshape #1}}

% These tags are bookkeeping labels for consistency audits. They do not
% create new mathematical hypotheses; they record how a datum is being used
% in the proof chain.
\providecommand{\StatusPrimitive}{}
\renewcommand{\StatusPrimitive}{\AuditLabel{Primitive}}
\providecommand{\StatusDerived}{}
\renewcommand{\StatusDerived}{\AuditLabel{Derived}}
\providecommand{\StatusImported}{}
\renewcommand{\StatusImported}{\AuditLabel{Imported}}
\providecommand{\StatusConditional}{}
\renewcommand{\StatusConditional}{\AuditLabel{Conditional}}
\providecommand{\StatusBridge}{}
\renewcommand{\StatusBridge}{\AuditLabel{Bridge-derived}}
\providecommand{\StatusComparison}{}
\renewcommand{\StatusComparison}{\AuditLabel{Comparison-only}}

% ======================================================================
% 10C. SERIES-WIDE CLAIM-WORDING AUDIT LABELS
% ======================================================================
% These are textual audit labels used to keep modal language consistent
% across the series.  They introduce no new mathematical assumptions.
\providecommand{\ClaimForced}{}
\renewcommand{\ClaimForced}{\AuditLabel{forced}}
\providecommand{\ClaimUnique}{}
\renewcommand{\ClaimUnique}{\AuditLabel{unique}}
\providecommand{\ClaimCanonical}{}
\renewcommand{\ClaimCanonical}{\AuditLabel{canonical}}
\providecommand{\ClaimMinimal}{}
\renewcommand{\ClaimMinimal}{\AuditLabel{minimal}}
\providecommand{\ClaimClosed}{}
\renewcommand{\ClaimClosed}{\AuditLabel{closed}}
\providecommand{\ClaimPrediction}{}
\renewcommand{\ClaimPrediction}{\AuditLabel{prediction}}

% ======================================================================
% 10D. SERIES-WIDE THEOREM-CHAIN AUDIT LABELS
% ======================================================================
% These labels mark the direction of load-bearing theorem chains. They are
% audit annotations only; they introduce no new mathematical assumptions.
\providecommand{\ChainInput}{}
\renewcommand{\ChainInput}{\AuditLabel{Chain input}}
\providecommand{\ChainOutput}{}
\renewcommand{\ChainOutput}{\AuditLabel{Chain output}}
\providecommand{\ChainGuard}{}
\renewcommand{\ChainGuard}{\AuditLabel{Chain guard}}
\providecommand{\ChainNoBackflow}{}
\renewcommand{\ChainNoBackflow}{\AuditLabel{No backflow}}

% ======================================================================
% 10E. SERIES-WIDE LOCAL THEOREM-AUDIT LABELS
% ======================================================================
% These labels mark a statement-level dependency audit around individual
% theorem/definition nodes. They are bookkeeping labels only.
\providecommand{\LocalAudit}{}
\renewcommand{\LocalAudit}{\AuditLabel{Local audit}}
\providecommand{\UsesOnly}{}
\renewcommand{\UsesOnly}{\AuditLabel{Uses only}}
\providecommand{\Exports}{}
\renewcommand{\Exports}{\AuditLabel{Exports}}
\providecommand{\DoesNotUse}{}
\renewcommand{\DoesNotUse}{\AuditLabel{Does not use}}

% ======================================================================
% 10F. SERIES-WIDE NUMERICAL-ROW AUDIT LABELS
% ======================================================================
% These labels classify individual rows in numerical, benchmark, status,
% and comparison tables.  They are audit annotations only and do not change
% any theorem, value, formula, or comparison number.
\providecommand{\RowAudit}{}
\renewcommand{\RowAudit}{\AuditLabel{Numerical-row audit}}
\providecommand{\RowTheorem}{}
\renewcommand{\RowTheorem}{\AuditLabel{Theorem-fixed row}}
\providecommand{\RowBridge}{}
\renewcommand{\RowBridge}{\AuditLabel{Bridge-fixed row}}
\providecommand{\RowRepresentative}{}
\renewcommand{\RowRepresentative}{\AuditLabel{Representative row}}
\providecommand{\RowComparison}{}
\renewcommand{\RowComparison}{\AuditLabel{Comparison row}}
\providecommand{\RowCalibrated}{}
\renewcommand{\RowCalibrated}{\AuditLabel{Calibrated row}}
\providecommand{\RowEstimate}{}
\renewcommand{\RowEstimate}{\AuditLabel{Estimate row}}
\providecommand{\RowOpen}{}
\renewcommand{\RowOpen}{\AuditLabel{Open row}}
\providecommand{\RowNoSelector}{}
\renewcommand{\RowNoSelector}{\AuditLabel{No selector}}

% ======================================================================
% 10G. TABLE HELPERS
% ======================================================================
% Table helpers used by Paper V-C.  The unused readout-label preset has been
% removed; only active table helpers remain.
\providecommand{\gpsection}[1]{%
  \addlinespace[0.55em]
  \multicolumn{3}{@{}l}{\textsc{#1}}\\[-0.25em]
  \midrule
}
\providecommand{\mathcell}[1]{%
  \(\begin{gathered}#1\end{gathered}\)%
}

\newcommand{\DerBy}[1]{%
  \begin{tabular}[t]{@{}l@{}}#1\end{tabular}%
}

\newcommand{\TableNotes}[1]{%
  \par\medskip
  \begin{minipage}{0.94\linewidth}
  \footnotesize\emph{Notes.} #1
  \end{minipage}\par
}
% ======================================================================
% 10H. PAPER VI-A / QCD NOTATION
% ======================================================================
% Shared symbols used by Paper VI-A and downstream QCD/Yang--Mills handoffs.
\providecommand{\Pqcd}{\mathfrak P_{\rm VIII}^{\rm QCD}}
\providecommand{\Iqcd}{\mathfrak I_{\rm QCD}^{(7)}}
\providecommand{\muQCD}{\mu_{\rm QCD}^{\rm int}}
\providecommand{\LambdaQCD}{\Lambda_{\rm QCD}}
\providecommand{\LambdaBH}{\Lambda_{\rm BH}^{\rm QCD}}
\providecommand{\HQCD}{\mathcal H_{\rm QCD}^{\rm BH}}
\providecommand{\HQCDfin}{H_{\rm QCD}^{\rm fin}}
\providecommand{\Hraw}{\mathcal H_{\rm raw}}
\providecommand{\HphysQCD}{\mathcal H_{\rm phys}^{\rm QCD}}
\providecommand{\Vcol}{V_C}
\providecommand{\Ps}{P_{\rm sing}}
\providecommand{\PiOct}{\Pi_8}
\providecommand{\MSbar}{\overline{\rm MS}}
\providecommand{\YM}{\mathrm{YM}}
\providecommand{\Ktri}{K_{\triangle}}
\providecommand{\KYM}{K_{\YM}^{\rm BH}}
\providecommand{\aBH}{a_{\rm BH}}
\providecommand{\epsBH}{\varepsilon_{\rm BH}}
\newcommand{\RowConstructive}{\textnormal{\textsc{Constructive}}}
\newcommand{\RowObligation}{\textnormal{\textsc{Bridge obligation}}}
\newcommand{\RowTarget}{\textnormal{\textsc{Target row}}}
\newcommand{\RowValidation}{\textnormal{\textsc{Validation row}}}
\providecommand{\slashed}[1]{\not{#1}}

% ======================================================================
% 10I. PAPER VI-B / CONTINUOUS-QFT NOTATION
% ======================================================================
% Shared symbols used by Paper VI-B and any downstream continuous-QFT or
% Yang--Mills handoff paper.  All definitions are guarded so individual
% papers may override them locally before loading this file if needed.
\providecommand{\A}{\mathfrak A}
\providecommand{\Hphys}{\mathcal H_{\rm phys}}
\providecommand{\QCD}{\mathrm{QCD}}
\providecommand{\BH}{\mathrm{BH}}
\providecommand{\loc}{\mathrm{loc}}

\providecommand{\RBH}{\mathcal R_{\rm BH}}
\providecommand{\DG}{\mathfrak D_G}
\providecommand{\PW}{\mathcal{P\!W}} % Peter--Weyl spectral-symbol subspace
\providecommand{\Toe}{\operatorname{Toe}} % Toeplitz compression/representative
\providecommand{\Hol}{\operatorname{Hol}} % Holonomy

% ======================================================================
% 11. LATE-LOADED PACKAGES (Must be at the end of the preamble)
% ======================================================================
% Load late packages only if a paper has not already loaded them, avoiding
% option clashes in paper-specific preambles.
\makeatletter
\@ifpackageloaded{hyperref}{}{\usepackage{hyperref}}
\makeatother
% --- Series-wide disambiguations (Pass 3) ---
\providecommand{\diracdelta}{\delta_{\rm D}} % Dirac delta distribution (delta reserved for hierarchy)
\providecommand{\Hg}{H_{\mathrm g}} % Hodge harmonic projector (keep H for Hubble)
\providecommand{\Gr}{G_{\mathrm r}} % Green operator (keep G for Newton)
\providecommand{\FD}{\mathrm{D}} % functional/variational derivative symbol (delta reserved)
\providecommand{\Id}{\mathbf{1}} % identity/Kronecker tensor (avoid \delta_{ij})


% BEGIN INLINE SOURCE: Paper00-ActiveExternalReferences-v2_130.tex
% Framework v2.191: current reference binding; existing definitions retained.
% Framework v2.190: current22-root registry with native clock-evolution proofs.
% Framework v2.189: exact current22-root external registry, including retained v2.188 hardening owners.
% Framework v2.187: current registry synchronization only; source release v2.186 preserved.
% Framework v2.186: authoritative current 22-paper registry.
% Framework v2.182: all current 22 logical paper AUX bindings.
% Framework v2.170: current I-B v1.43 external AUX binding.
% Framework v2.167: exact current 22-root registry after CP--BL2 integration.
% Framework v2.166: exact current 22-root compact-release AUX registry.
% Framework v2.165: authoritative current root bindings.
% Framework v2.164: exact current 22 manuscript AUX bindings.
% Framework v2.162: finite metric, functional Hilbert and scoped Ward integration.
% Framework v2.161: current 22 manuscript AUX bindings; source only.
% Framework v2.156 / registry v2.113: synchronized current 22 manuscript roots; source-only registry.
% Framework v2.143 / registry v2.111: active manuscript AUX basenames; no generated AUX or PDF included.
% ======================================================================
% Central external-reference imports, numbers only, no PDF jumping
% ======================================================================

% xr-hyper is loaded centrally so that papers may retain their existing
% \usepackage{xr-hyper} lines without redefining \externaldocument after this
% registry is installed.  If plain xr was loaded earlier, retain it.
\makeatletter
\@ifpackageloaded{xr-hyper}{}{%
  \@ifpackageloaded{xr}{}{\usepackage{xr-hyper}}%
}
\makeatother

% Make external-document registration idempotent.  The registry identifies a
% logical paper by its Paper<number><optional letter> prefix.  Prefix-qualified
% imports remain independent.  Current paper roots contain no local
% \externaldocument calls; this guard also makes accidental duplicates harmless.
\ExplSyntaxOn
\prop_new:N \g_bh_external_refs_prop
\cs_new_eq:NN \BHOriginalExternalDocument \externaldocument

\cs_new_protected:Npn \bh_external_document:nnnn #1#2#3#4
 {
  \tl_set:Nn \l_tmpa_tl {#3}
  \regex_replace_once:nnN
   { \A (?: xref/ )? (Paper[0-9]+[A-G]?) .* \Z }
   { \1 }
   \l_tmpa_tl
  \tl_set:Nx \l_tmpb_tl { #1 | \tl_use:N \l_tmpa_tl }
  \prop_if_in:NVTF \g_bh_external_refs_prop \l_tmpb_tl
   {
    \PackageInfo{BornHopf}
     {Skipping~duplicate~external~refs~for~#3}
   }
   {
    \prop_gput:NVn \g_bh_external_refs_prop \l_tmpb_tl {#3}
    \IfNoValueTF{#4}
     { \BHOriginalExternalDocument[#1][#2]{#3} }
     { \BHOriginalExternalDocument[#1][#2]{#3}[#4] }
   }
 }

\RenewDocumentCommand{\externaldocument}{O{} O{} m o}
 { \bh_external_document:nnnn {#1}{#2}{#3}{#4} }

\NewDocumentCommand{\BHExternalRefs}{m}
 {
  % Do not import the aux file of the logical paper currently being compiled,
  % even when the job and registry use different version suffixes.
  \tl_set:Nx \l_tmpa_tl { \c_sys_jobname_str }
  \regex_replace_once:nnN
   { \A (Paper[0-9]+[A-G]?) .* \Z }
   { \1 }
   \l_tmpa_tl
  \tl_set:Nn \l_tmpb_tl { #1 }
  \regex_replace_once:nnN
   { \A (Paper[0-9]+[A-G]?) .* \Z }
   { \1 }
   \l_tmpb_tl
  \str_if_eq:eeTF { \tl_use:N \l_tmpa_tl } { \tl_use:N \l_tmpb_tl }
   {
    \PackageInfo{BornHopf}
     {Skipping~external~refs~for~#1~because~this~is~the~current~document}
   }
   {
    \IfFileExists{xref/#1.aux}
     { \externaldocument[][nocite]{xref/#1}[#1.pdf] }
     {
      \IfFileExists{#1.aux}
       { \externaldocument[][nocite]{#1} }
       {
        \PackageInfo{BornHopf}
         {External~aux~file~#1.aux~not~found;~no~labels~imported}
       }
     }
   }
 }
\ExplSyntaxOff

% External references are automatic.  This file is the sole owner of both the
% bibliography and the current Paper-I-A--VII-D AUX registry (30 roots including I-C).  Missing auxiliaries
% are informational on the first pass and are populated by the source-only
% two-pass series build.  Basenames deliberately omit temporary download/copy
% suffixes such as ``(1)'' or ``-1''.
  % Paper I-A: round S3 compendium.
  \BHExternalRefs{Paper1A-Geometry-v5_181}

  % Paper I-B: complete exceptional geometry and finite spectral structures.
  \BHExternalRefs{Paper1B-ExceptionalGeometry-v1_25}
  % Paper I-D: technical companion of I-A through I-C.
  \BHExternalRefs{Paper1D-TechnicalMaterial-v1_69}
  % Paper I-C: reciprocal M8 compendium.
  \BHExternalRefs{Paper1C-M8Geometry-v1_146}

\ifdefined\BHMathOnly
  % Papers I-A and I-B import only the mathematical atlas labels.
\else
  % Paper II: five physical-mathematical main papers plus one technical proof volume.
  \BHExternalRefs{Paper2A-GeometricConstitution-v3_227}
  \BHExternalRefs{Paper2B-MomentumCarrier-v2_147}
  \BHExternalRefs{Paper2C-ReciprocalLocalization-v2_176}
  \BHExternalRefs{Paper2D-SymmetrySpectra-v2_145}
  % The former II-F export labels remain in II-E; current II-F is the technical volume.
  \BHExternalRefs{Paper2E-KinematicExports-v3_140}
  % Paper II-F: complete technical proof and calculation owner for II-A--II-E.
  \BHExternalRefs{Paper2F-TechnicalMaterial-v1_71}

  % Paper III: Dynamics and Intrinsic Geometric Constitution (six components).
  \BHExternalRefs{Paper3A-GeneratedDynamics-v1_181}
  \BHExternalRefs{Paper3B-SpectralAction-v1_192}
  \BHExternalRefs{Paper3C-ProtectedHierarchy-v1_196}
  \BHExternalRefs{Paper3D-RunningResponse-v1_168}
  \BHExternalRefs{Paper3E-GravityAction-v1_204}
  \BHExternalRefs{Paper3F-CompleteDynamics-v1_279}
  % Paper III-G: complete preserved technical proof owner.
  \BHExternalRefs{Paper3G-TechnicalMaterial-v1_57}

  % Paper IV-A: generated finite quantum mechanics.
  \BHExternalRefs{Paper4A-FiniteQuantumMechanics-v4_217}
  % Paper IV-B: complete finite quantum processes and physical response.
  \BHExternalRefs{Paper4B-QuantumResponse-v1_55}
  % Paper IV-C: extended proofs and finite technical calculations.
  \BHExternalRefs{Paper4C-TechnicalMaterial-v1_50}

  % Paper V-A: The Standard Model.
  \BHExternalRefs{Paper5A-FiniteMatterGeometry-v3_218}

  % Paper V-B: Finite Spectrum and Flavor Closure.
  \BHExternalRefs{Paper5B-TerminalSpectrum-v5_202}

  % Paper V-C: Finite Interaction Boundaries of the Standard Model.
  \BHExternalRefs{Paper5C-InteractionResponse-v2_184}
  % Paper V-D: preserved technical proof and calculation owner.
  \BHExternalRefs{Paper5D-TechnicalMaterial-v1_58}

  % Paper VI-A: Gravity and Closed Cosmology.
  \BHExternalRefs{Paper6A-GravityCosmology-v2_268}

  % Paper VI-B: Geometric Dark Sector and Cosmological Readouts.
  \BHExternalRefs{Paper6B-CosmologyReadouts-v1_215}

  % Paper VI-C: complete technical proof and comparison owner.
  \BHExternalRefs{Paper6C-TechnicalMaterial-v1_53}

  % Paper VII-A: Finite QCD Kernel and Nonperturbative Closure.
\BHExternalRefs{Paper7A-QCDKernel-v1_183}

  % Paper VII-B: Hadronic, Thermal, and Nuclear QCD Readouts.
  \BHExternalRefs{Paper7B-QCDReadouts-v1_167}

  % Paper VII-C: Continuous Quantum Field Theory.
  \BHExternalRefs{Paper7C-ContinuousQFT-v2_208}

  % Paper VII-D: complete technical material for VII-A through VII-C.
  \BHExternalRefs{Paper7D-TechnicalMaterial-v1_60}

\fi

% END INLINE SOURCE: Paper00-ActiveExternalReferences-v2_130.tex


\hypersetup{hidelinks}


\usepackage{xr-hyper}

\title{\textbf{Born--Hopf Geometry III-E: Spectral Gravity and Conditional Scalaron Branches}\\[0.5em]
\large Geometric action, capped profiles, and source-port requirements}

\author{Peter Kugelmann\thanks{\texttt{pkugelmann@proton.me}}}
\date{Version 1.204 --- 20 November 2026 \\ Framework v3.073}
\setlength{\emergencystretch}{3em}
\begin{document}
\maketitle
\begin{abstract}
Paper III-E develops the spectral geometric gravitational
functional and its normalization, scalar threshold, smooth hierarchy response,
and the capped action used as the current gravitational working branch.
The capped reduction is exact on its specified boundary/source port domain,
while identification with the native microscopic common action remains open. Paper III-F owns full physical source and metric
work; detailed branch comparisons are maintained in Paper~III-G.
\end{abstract}
\tableofcontents
\clearpage

\section{Inherited spectral geometry and action scope}
\label{paper3e:sec:organized-scope}
\subsection{Introduction and Paper-III series allocation}
\label{p3e:sec:series-allocation}
Paper~III-E imports the fixed momentum seed, reciprocal-role interfaces,
completed spectral-action class, protected hierarchy data, and exact reduced
response from Papers~II-B and III-A--III-D.  It owns the Einstein--Hilbert
normalization handoff, the Lorentzian/causal scalar domain, the physical
shell-transport interpretation of the already constructed III-D entropy
coordinate, and any additional gravitational calibration used to place the
abstract response roots on an Einstein--Hilbert curvature scale.  Common-action
backreaction remains the task of Paper~III-F.

\subparagraph{Inherited geometric and finite-role prefix.}
The current upstream chain is
\[
\begin{gathered}
\text{Paper~II-B}
\Longrightarrow
\mathcal H_p\cong\mathcal H_1,\qquad m_p=1,\\
+\text{Paper-II-C intrinsic }E_1\text{-evaluation}\\[-2pt]
\xRightarrow{\text{Paper~II-A/III-A GP2 reciprocal closure}}\\[-2pt]
\text{distinct Born/GP2-generated reciprocal roles},\\
\text{phase interfaces and }\mathcal B^2=1,\\
+\text{generated self-adjoint operator owner}\\[-2pt]
\xRightarrow{\text{Paper~III-B native heat closure}}\\[-2pt]
(\mathrm{CLOOP}_{\rm heat}^{\rm der}),\quad
\text{loop-derived spectral form }\Dcl,\\
+(\mathrm{ND}_{\rm sp})+\text{occurring Hurwitz-active orbit family}\\[-2pt]
\xRightarrow{\text{Paper~III-C}}\\[-2pt]
\mathsf C_{\rm IR}^{\rm der}:\quad
\Cl(0,3)=\mathbb R1\oplus W,\quad \dim_{\mathbb R}W=7,\\
 +({\rm RED})+({\rm UV})+({\rm IR})
\xRightarrow{\text{Paper~III-D}}
(C,U,g_{\rm eff},n_{\rm sec}),\\[-2pt]
0<C<U,\qquad
\scal_{\rm resp}=C+U\quad\text{in the adopted }\kappa=2
\text{ normalization representative}.
\end{gathered}
\]
Here \(\mathsf C_{\rm IR}^{\rm der}\) abbreviates
\((\mathrm{ND}_{\rm sp})\), a G0-occurring Hurwitz-even active orbit family,
and its orbit-invariant terminal identity \(P_{\rm term}=P_0\).  It is not a new
assumption package: Paper~II-C proves that all nonzero protected-coordinate
completions have the same normalized represented even algebra, while finite
functional calculus, equivariant descent, and GP2 co-realization inherit the
active composite's G0 witness without selecting a diagonal.  The historical
\((\mathrm{IR}_{O(3)})\) tag certifies the resulting
parity-even Haar projector.  The subnormalized real scalar-corner readout used
in the Einstein--Hilbert projection is derived from the occurring ambient
protected-octet presentation; \((\mathrm{CORN})\) records that certification.
The remaining protected-coordinate antecedent is only
\[
 (\mathrm{ND}_{\rm sp}):\qquad \varrho_{\rm sp}>0.
\]
The carrier identity \(\mathcal H_p\cong\mathcal H_1\) does not itself force
this nonvanishing, and no absolute coordinate-completion scale is inferred.
The fixed-shell realization family imported from Papers~III-C and III-D is
\[
 \boxed{\mathfrak A_{\rm REG}
 =\{\mathrm{REG}_{\sigma}:\sigma_{\rm fin}\in\mathbb R\}},
 \qquad
 \delta_{\sigma}=\frac1{6\sqrt2}+\sigma_{\rm fin}.
\]
For every frozen member,
\(H_{56}^{(\sigma)}=\alpha_{\rm hier}(\sigma_{\rm fin})^7\).
The notation
\[
 (\mathrm{REG}_0):\qquad \sigma_{\rm fin}=0,
 \qquad
 H_{56}^{\rm REG}=\alpha_{\rm hier}(0)^7
\]
denotes the source member only; neither GP closure nor the finite protected
classification selects it.
These are neutral determinant characters.  Their dual ratios are generated by
Paper~III-C; action on an occurring physical scale requires a proved generated
ancestry morphism into the determinant-dual owner.  The owner port, the
dimensionful reference magnitude, and its unit convention are separate
interfaces.  On the MRESP branch, the ancestry morphism in
\((\mathrm{CRP4})\) is the physical specialization of Paper~III-D's abstract
\((\mathrm{OWN}_{56})\) interface.  The historical
\((\mathrm{OWN}_{\rm resp/det})\) label is only an alias for that summand,
not a second owner-ancestry premise.

\subparagraph{Imported stage-safe G0.}
All physical interfaces below specialize Paper~II-A's G0 certificate.  Each
asserted occurrence requires a contingent token of a geometrically licensed,
claim-relative invariant distinction with result-independent provenance; its
tag alone supplies no witness, and no instance reopens frozen upstream
geometry.  A new role or standard enters only through a declared downstream
frontier.  The clocked-Hopf seed and its least operation doctrine are
choice-free generated descendants of the actual scalar-succession provenance
and require no new G0 occurrence certificate.  The typed terminal endpoint
data \((\mathrm{TERM}_{f})\) and a realized nonlinear response
\((\mathrm{MRESP}_{\min})\) remain G0 boundary/process or occurrence
certificates.
The tags \((\mathrm{CORN})\) and
\((\mathrm{BD}_{\rm EH})\) instead record the derived constituent and
Born-orbit identification theorems below; the former \((\mathrm{MON})\)
one-shell occurrence tag is retired.  The gravitational heat role is likewise
a choice-free generated descendant of the occurring scalar metric/protected
split.  The labels \((\mathrm{GRAV}_{\circ})\) and
\((\mathrm{GRAV\!-\!HEAT})\) certify its co-realization and carry no
independent process witness; its outer trace and exponential law are derived
below.  A non-exponential gravitational profile would be an enlarged process
datum and would require a new G0 root.  The reciprocal metric placement
\((\mathrm{CAL}_{xp})\), the primitive-turn duration
\((\mathrm{CAL}_t):T_t>0\), the response/EH upper-scale magnitude placement
\((\mathrm{CAL}_{\rm resp/EH})\), the dimensionless
\(\gamma_\sigma\), and the numerical realization of the Born action unit
\(\hbar\) are calibration/metrological data: they fix representatives of
existing invariants and do not add new geometric structure.  Newton's
coupling is instead a downstream coordinate \(G_{J,\rm raw}\) of the generated
joint-EH coefficient; comparison with an observed value is validation.  The action-role bridge
\((\mathrm{ACT}_{\rm EH}^{\rm cl})\) is instead a typed relation, not a
numerical calibration.  Its local joint-Dirac realization is derived later in
Paper~III-F.  The raw joint/scalar relation is $U_J^{\rm raw}=96U$. It leaves distinct
scalar and raw-joint Newton coordinates. The complete physical coefficient
and its source normalization must be matched before comparison with measured
Newton data. The intrinsic momentum curvature and Hurwitz cell normalization
are known; their physical EH matching is the evaluated
$(\mathrm{CELL}_{K\to U})$ test.

The $(\mathrm{CRP4})$ tag is an evaluated cross-role certificate,
not a calibration and not a consequence of equal numerical values.  Its
historical response-owner alias carries no second occurrence.
The legacy
\((\mathrm{MATCH}_{56})\) tag now certifies the generated terminal ratio and
adds no calibration.  Within Paper~III-E the claim-specific Wick tags mark reconstruction
claims not proved by the elliptic heat calculus itself; G0 types any physical
use but does not derive those reconstruction statements.  In the
Paper~III-A theorem-status taxonomy they are analytic-certification interfaces,
not native same-owner closure.  Paper~II-A's unmarked-automorphism theorem
cannot create the missing continuation/reconstruction owner map, because its
constructor-origin hypothesis already requires the relevant output sort to be
licensed.  Paper~IV may later
discharge finite certificates on explicitly stated Hamiltonian-generated or
positive-kernel domains, but such later theorems are not imported as unstated
Paper~III premises.  No independent momentum witness is needed; particular
processes, bridges, and readouts remain token-level G0 claims.

\subparagraph{Status synchronization after the full residual reclassification.}
Generated or co-realized structure includes the momentum seed, reciprocal
product, protected \(1\oplus7\) split, generated gravitational heat role,
closed trace, Born-diagonal Einstein--Hilbert copy, causal operation cone,
hierarchy-active state, determinant-depth ratios, minimal-response transport,
and entropy coupling.  Actual response processes, local leaves, terminal or
transition histories, and physical readouts remain G0 occurrences or solution
data.  An occurring complete minimal-response root includes the physical
\((\mathrm{ACT}_{\rm MRESP})\) action-role adoption when that role is claimed;
the fields of this bookkeeping tuple are not generated from one another.
Conditional on the complete root, the displayed response graph and transports
are generated descendants.  Numerical unit placements remain calibrations.

The REG theorem is uniform over the positive affine scheme family.  On an unanchored branch, a numerical point evaluation such as REG0 is a
boundary choice. The occurring octahedral compiler anchors its declared
formal half-determinant prescription; this does not identify a physical
complex Gaussian measure.
The optional rank-one \((\mathrm{FLAG})\) and marked profile refinements are
retained only when a theorem explicitly consumes member-specific projectors.
The classical Einstein--Hilbert action-role adoption, physical POVM,
Wick/reconstruction, and interacting-QFT hypotheses remain separately typed
physical or analytic interfaces; none is added to the structural-root count.


\begin{proposition}[Analytic-reconstruction constructor firewall]
\label{p3e:prop:v169-analytic-reconstruction-firewall}
Let \(D_E\) be any Euclidean heat/spectral owner generated in the present
paper, and let \(Q_L\) denote a Lorentzian real-time or reconstructed
correlation role.  Paper~II-A's same-owner unmarked-automorphism theorem can
prove equivariance or invariance of a constructor output only after a licensed
constructor word
\[
   w:D_E\longrightarrow Q
\]
has already been generated in the complete owner.  It therefore cannot, from
Euclidean covariance or internal automorphism invariance alone, manufacture a
new constructor
\[
   D_E\longrightarrow Q_L,
\]
nor can it supply analyticity, contour control, reflection positivity,
same-field compatibility, thermodynamic/infinite-volume convergence, or the
regularity required by a reconstruction theorem.  Consequently the
claim-specific \((\mathrm{WICK}_\bullet)\) interfaces below are analytic
certifications in the sense of Paper~III-A, not additional physical principles
and not consequences of G0/GP2 same-owner closure.  If a later theorem proves
all required analytic hypotheses on a narrower generated domain, its
reconstruction output is then a theorem-descendant on that domain only.
\end{proposition}

\begin{proof}
The constructor-origin rule of Paper~II-A is antecedent to its naturality and
same-owner equivariance conclusions: those conclusions compare already
licensed outputs of one constructor under unmarked automorphisms.  A
Euclidean-to-Lorentzian reconstruction instead changes the output role and
requires extra analytic relations not present in the Euclidean heat owner.
GP2 forbids inserting that missing relation, while G0 types an actual use but
does not prove the analytic hypotheses.  Hence internal same-owner closure
cannot discharge the reconstruction interface.  When an independent analytic
theorem does establish the missing hypotheses and output map, ordinary
recursive theorem closure applies exactly on its proved domain.
\end{proof}

\begin{definition}[Derived causal domain and remaining continuation/calibration hypotheses]
\label{p3:def:causal-wick-response-hypotheses}
The following tags separate theorem-level upstream structure from genuinely
additional III-E continuation or placement data.
\begin{enumerate}[label=(\roman*)]
\item \emph{Derived operational causal domain \((\mathrm{CAUS})\).}
Paper~I-A supplies the positive primitive Reeb lattice, Paper~II-A transports
it through the actual scalar-succession provenance occurrence and includes
the upstream provenance-fixed representative of the unmarked Hopf class in the vector parent's opposite, and Paper~III-A
composes these arrows with the faithful reciprocal isometry action.  Its
choice-free seed and least constructor closure give
\(\mathscr O_{\rm prop}^{\rm phys}\), containing the stationary history and a
nonstationary clocked Hopf history and closed under the already generated
spatial isometries, scalar-cover translations, restriction, and compatible
sequential composition.  Its time-sharing/closed convex relaxation is a
mathematical causal envelope, not an additional physical history doctrine.
For the explicit clock calibration
\((\mathrm{CAL}_t):T_t>0\),
Theorem~\ref{p3:thm:gp1-speed-calibration} gives
\[
 d_h(x(t),x(s))\le c_{\rm op}|t-s|,
 \qquad c_{\rm op}=\frac{2\pi R_x}{T_t}.
\]
The tag abbreviates the resulting sharp closed operation cone and the
Lorentzian conformal class proved below within this occurrence domain; it is
therefore a theorem-status name, not an independent causal-sign or
speed-ceiling hypothesis.  Homothetic changes of \(h\), elapsed coordinate,
and \(c_{\rm op}\) leave the cone and conformal class unchanged.
\((\mathrm{CAL}_{xp})\) is needed only to attach numerical spatial units;
\((\mathrm{CAL}_t)\) attaches the duration assigned to one already generated
positive primitive turn.  Existence and nonstationarity of the seed, and co-realization
of its least generated doctrine, are derived; this does not assert separate
physical-token occurrence for every member history.  Exhaustion of the cone
boundary by a field or light channel, the numerical witnesses for \(R_x\) and
\(T_t\), and an actual system trajectory are not derived here.

\item \emph{Claim-specific Euclidean reconstruction hierarchy.}
The single legacy label \((\mathrm{WICK})\) is replaced by the
claim-scoped tags
\[
\begin{aligned}
(\mathrm{WICK}_{2{\rm cyc}})&:\ \text{cyclic two-point continuation},\\
(\mathrm{WICK}_{\rm same})&:\ \text{same Euclidean and Lorentzian}\\
&\hspace{1em}\text{field-variable compatibility},\\
(\mathrm{WICK}_{\rm OS,fin})&:\ \text{finite-certificate reflection positivity}\\
&\hspace{1em}\text{and reconstruction},\\
(\mathrm{WICK}_{\rm int,\infty})&:\ \text{interacting infinite-volume reconstruction.}
\end{aligned}
\]
Within Paper~III-E each tag marks exactly the analyticity,
reflection-positivity, contour, compatibility, and limit conditions not proved
here but needed for its named reconstruction conclusion.  By
Proposition~\ref{p3e:prop:v169-analytic-reconstruction-firewall}, these are
analytic-certification hypotheses rather than G0 roots, GP assumptions, or
native same-owner theorem closure.  Elliptic spectral calculus alone implies
none of those reconstructions.  Paper~IV later proves
finite certificates on narrower positive-kernel or Hamiltonian-generated
domains.  Accordingly \((\mathrm{WICK}_{2{\rm cyc}})\),
\((\mathrm{WICK}_{\rm same})\), and finite
\((\mathrm{WICK}_{\rm OS,fin})\) are \texttt{DERIVED-LATER} on precisely
those stated Paper-IV theorem domains.  It now also packages Euclidean
contraction, Lorentzian unitary evolution, and their state/test duals into one
complex-time co-realized diagram, and derives the action-form continuation for
stationary stable quadratic bosonic blocks.  These results do not identify the
spectral cutoff \(\Lambda^{-2}\) with elapsed Euclidean time and do not
continue the complete spectral action.  The interacting infinite-volume tag
is not discharged.  Those later results do not alter the present theorem
chronology.  The free Lorentzian field equation derived from the operational
cone below requires no Wick tag.

\item \emph{Classical Einstein--Hilbert action role, typed cross-owner
matching, and Newton-coordinate readout.}
The heat coefficient proved below is Euclidean and belongs to the
unit-preserving scalar-resolution owner.  The historical claim-scoped action
bridge is
\[
 (\mathrm{ACT}_{\rm EH}^{\rm cl})
 :=(\mathrm{LOR}_{\rm dir})\ \vee\ (\mathrm{CONT}_{\rm cl}).
\]
Paper~III-F later derives the local interior action-role component from the
joint Dirac owner on its analytic clock--ADM first-jet domain.  Hence
\((\mathrm{ACT}_{\rm EH}^{\rm cl})\) is \texttt{DERIVED-LATER} there as an
action-role statement.  It does not identify the two owner coefficients.

For the generated joint-Dirac coefficient define, without using measured
Newton data,
\begin{equation}
\label{p3e:eq:v267-joint-eh-coordinate-readout}
 \boxed{\UEHjoint:=8\kappa_J,\qquad
        \hBorn:=2\pi\hbar,\qquad
        G_{J,\rm raw}:=\frac{c_{\rm op}^{3}}{\hBorn\UEHjoint}.}
\end{equation}
Paper~III-A proves that \(c_{\rm op}^{3}\) is the determinant character of
isotropic propagation on the generated three-dimensional spatial carrier, and
\(\hBorn\) is the physical representative of one primitive Born scalar action
cell.  Equation~\eqref{p3e:eq:v267-joint-eh-coordinate-readout} is therefore a
coupling-coordinate readout of the already generated joint-EH coefficient; it
uses no observed value of \(G\).  Algebraically it is equivalent to
\[
 \hbar\kappa_J=\frac{c_{\rm op}^{3}}{16\pi G_{J,\rm raw}}.
\]
Denote this theorem-status coordinate assignment by
\((\mathrm{GREAD}_{J})\).

The downstream common-heat carrier theorem replaces the former equality
condition by the exact raw transport
\begin{equation}
\label{p3e:eq:v267-cross-owner-eh-match}
 \boxed{
 (\mathrm{TRANSPORT}_{\rm J/scal}^{\rm EH,raw}):\qquad
 U_J^{\rm raw}=96\Lambda_{\rm EH}^{2}.}
\end{equation}
The primitive curvature coordinate is supplied by the momentum-curvature
actualization below rather than by a second owner branch.  We therefore write
\begin{equation}
\label{p3e:eq:conditional-newton-matching}
 \boxed{\begin{gathered}
 \Rp:=U,\qquad G_{\rm scal}:=\frac{c_{\rm op}^{3}}{\hBorn U},\\
 G_{J,\rm raw}:=\frac{c_{\rm op}^{3}}{\hBorn U_J^{\rm raw}}
 =\frac{G_{\rm scal}}{96}.
 \end{gathered}}
\end{equation}
Here $\Rp$ is optional scalar-resolution notation. The same symbol cannot
name both displayed couplings when $U_J^{\rm raw}=96U$.
The Newton coupling of an actual final action is
$G_{\rm phys}=c_{\rm op}^3/(16\pi\hbar\kappa_{\rm phys})$ for its
complete dimensionless EH coefficient $\kappa_{\rm phys}$. If the physical
action is the raw joint action, $\kappa_{\rm phys}=\kappa_J$ and
$G_{\rm phys}=G_{J,\rm raw}$. A different normalized projection must
transport the complete action, every source/gauge/scalar term and mixed jet;
rescaling the gravity trace alone does not establish that projection.
Comparison of $G_{\rm phys}$ with observed Newton data is then validation.

Paper~III-F, in the definition entitled \emph{Physical action at a specified reference}%
\ifcsname r@p3f:def:v3041-adopted-metric-assembly\endcsname
\ (Definition~\ref{p3f:def:v3041-adopted-metric-assembly})%
\fi,
adopts the physical assembly \(M_{\rm rp}+S_J-i_v^*S_J\) on its
declared common metric, source and boundary chart. Here \(i_v\) is a
reference section fixed independently of the force. Its pullback gives
the complete joint work at that reference. The nonmetric variations at
fixed metric and boundary data are those of \(S_J\). With the stationary
flat reference, calibrated sources and single tensor pole specified in
III-F, the marked geometric coefficient \(U/8\) gives
\(G_{\rm tensor}=G_{\rm scal}\). Physical selection of the history,
additional long-range contributions and comparison with a measurement
retain their respective conditions. The following certificate treats the
case in which one constant weights the entire action.
\begin{proposition}[Whole-action Newton normalization certificate]
\label{p3e:prop:v2211-whole-action-newton-certificate}
Suppose an actual generated projection $\mathcal P$ maps the raw joint
action, fields, source marks, density and boundary data to a physical action.
If it is one common constant normalization $z>0$, then
\[
 \kappa_{\rm phys}=z\kappa_J,\qquad
 U_{J,\rm phys}=zU_J^{\rm raw},\qquad
 G_{\rm phys}=G_{J,\rm raw}/z.
\]
Therefore $G_{\rm phys}=G_{\rm scal}$ if and only if $z=1/96$,
\emph{provided the same map transports the complete source and every
retained variation}. This is a necessary/sufficient numerical condition
inside an already certified whole-action projection. The factor-96 raw
relation does not construct that projection. In particular no equality of
the two raw coordinates follows. A nonconstant or source-dependent map
retains its product-rule and mixed-jet terms and uses the actual resulting
$\kappa_{\rm phys}$ instead.
\end{proposition}
\begin{proof}
A common constant multiplies every action variation by the same $z$.
Substitute the resulting EH coefficient into the coupling definition and use
$G_{J,\rm raw}=G_{\rm scal}/96$. Source-dependent normalization is
not a common constant and its derivatives cannot be omitted.
\end{proof}

\item \emph{Derived Born-orbit Einstein--Hilbert certification
\((\mathrm{BD}_{\rm EH})\).}
On the smooth scalar parent, retain both reciprocal curvature functionals as
the finite Born orbit and evaluate their invariant average.  The generated
Born-diagonal embedding transports one scalar metric to its reciprocal copy.
Naturality of scalar curvature and its volume density then makes the pullback
of the two half-weighted orbit terms one full Einstein--Hilbert functional.
Theorem~\ref{p3e:thm:born-orbit-eh-identification} proves this statement.
Thus \((\mathrm{BD}_{\rm EH})\) is a theorem-status certification, not an
additional branch/interface hypothesis.  No equality of the two off-diagonal
metrics or finite matrix traces is asserted.

\item \emph{Identity-stage response/EH section and the momentum-cell test.}
Paper~III-B Theorem~\ref{p3b:thm:generated-action-response-owner} retains
$U=\Lambda_{\rm EH}^2$ at the identity stage of this actual action.
This is the initial coefficient identity, not contact transport.
The intrinsic momentum curvature and Hurwitz scalarizer are
\[
 \mathcal K_p=6/\Rmom^2,\qquad \rho_H=3/\pi.
\]
Their values do not alone identify the three-dimensional curvature owner
with the four-dimensional EH coefficient. The complete actual-action
comparison must pass the dimensionless residual
\begin{equation}
\label{p3e:eq:v2211-momentum-eh-port}
 \boxed{\varepsilon_{K\to U}:=
 \log\frac{U}{\rho_H\mathcal K_p}=0.}
\end{equation}
For the minimal scalar heat coefficient this condition is equivalent to
\begin{equation}
\label{p3e:eq:hurwitz-cell-normalized-eh-handoff}
 \boxed{U=\rho_H\mathcal K_p=\frac{18}{\pi\Rmom^2}
 \quad\Longleftrightarrow\quad
 \Lambda_g\Rmom=24\sqrt{3\pi}.}
\end{equation}
Indeed $U=\Lambda_g^2/(96\pi^2)$ and
$\rho_H\mathcal K_p=18/(\pi\Rmom^2)$ give exactly this equivalence.
It is an executable action/geometry matching test, not a new axiom setting
its residual to zero or a consequence of line duality. Denote a successful
origin-certified test by $(\mathrm{CELL}_{K\to U})$.

\begin{corollary}[Typed scalar Planck coordinate and conditional cell normalization]
\label{p3e:cor:typed-planck-normalization}
For the scalar-resolution coordinate
$G_{\rm scal}=c_{\rm op}^3/(2\pi\hbar U)$ put
$\ell_{\rm scal}^2=\hbar G_{\rm scal}/c_{\rm op}^3$. Then
\[
 \boxed{2\pi U=\ell_{\rm scal}^{-2}.}
\]
Only after $(\mathrm{CELL}_{K\to U})$ does one additionally have
\[
 2\pi U=6\mathcal K_p,\qquad
 \Rmom=6\ell_{\rm scal},\qquad
 4\pi^2U^2=36\mathcal K_p^2=\ell_{\rm scal}^{-4}.
\]
This scalar coordinate is not silently the Planck length of the raw joint
or the final physically matched action.
\end{corollary}
\begin{proof}
Substitute the definition of $G_{\rm scal}$, then the evaluated cell test.
No trace rescaling or measured Newton value enters the algebra.
\end{proof}

\item \emph{Certified cross-role positive parent $(\mathrm{CRP4})$.}
The action's evaluation pairing identifies its coefficient with the dual of
$\mathcal D_{\rm act}$. Its identification with Paper~III-C's contact
associated line requires the actual $(\mathrm{ACD}_H)$ certificate,
including origin, root normalization and transport compatibility.
Definition~\ref{p3e:def:crp4-cross-role-positive-parent} records that
successful port. The compatible positive fourth-root action is then unique.
The historical aliases $(\mathrm{OWN}_{56})$,
$(\mathrm{OWN}_{\rm resp/det})$ and $(\mathrm{ANC}_E)$ consume this same
port; no numerical coincidence or arbitrary flat gauge can replace it.

\item \emph{Fixed-shell realization family \((\mathrm{REG}_{\sigma})\).}
Import the unresolved fixed-shell family
\[
 \boxed{\mathfrak A_{\rm REG}
 =\{\mathrm{REG}_{\sigma}:\sigma_{\rm fin}\in\mathbb R\}},
 \qquad
 \delta_{\sigma}=\frac1{6\sqrt2}+\sigma_{\rm fin}.
\]
For every frozen member the determinant-depth character and neutral dual
ratios are theorem-level functions of \(\sigma_{\rm fin}\).  Ordered physical
scale ratios additionally require
\(\mathrm{REG}_\sigma\in\mathfrak A_{\rm REG}^{+}\) and the certified response
owner in \((\mathrm{CRP4})\).  On the generated \(2O\)
normalizer/interface branch, Paper~III-C
Theorem~\ref{p3:prop:conditional-2O-owner-REG0-closure} fixes
\[
 \boxed{(\mathrm{REG}_0):\sigma_{\rm fin}=0}
\]
as the unique anchored member.  Without that interface anchor the affine
family remains unresolved; the theorem does not manufacture an origin on
unanchored branches.

\item \emph{Derived owner-typed rank-seven ratio certification
\((\mathrm{MATCH}_{56})\), legacy \((\mathrm{HSC}_{56})\).}
The occurring MRESP composite contains the common-rate depth role, the
positive upper reference \(U\), and the certified CRP4 response ancestry.  At fixed
decreasing shell scheme
\(\mathrm{REG}_\sigma\in\mathfrak A_{\rm REG}^{+}\), Paper~III-C's neutral
dual-character theorem and owner-typed inheritance corollary therefore give
\[
 \boxed{\frac{U}{C}
 =H_{56}^{(\sigma)}
 =\alpha_{\rm hier}(\sigma_{\rm fin})^7.}
\]
Thus \(C=U/H_{56}^{(\sigma)}\) and
\(\scal_B=\sqrt{CU}=U/(H_{56}^{(\sigma)})^{1/2}\).
The source specialization uses the explicitly retained member \((\mathrm{REG}_0)\).
At fixed \(U\),
\[
 C(\sigma_{\rm fin})=C(0)e^{-7\sigma_{\rm fin}},\qquad
 \scal_B(\sigma_{\rm fin})=\scal_B(0)e^{-7\sigma_{\rm fin}/2}.
\]
The tags \((\mathrm{MATCH}_{56})\) and
\((\mathrm{HSC}_{56})\) certify this inherited identification and are not
placement axioms.  Without the ancestry port, only the neutral ratio is
available and these physical equations do not follow.  A nonzero empirical log-mismatch would deform the
generated character and is therefore a new response datum, not a hidden
scheme choice.

\item \emph{Typed terminal metric-response family \((\mathrm{TERM}_f)\).}
For the total metric response \(\mathfrak f\), let
\(\mathcal L_{\scal}\) denote the one-dimensional real value line with the
same physical dimension as \(\scal\).  The complete endpoint atlas is
\[
 \boxed{\mathfrak A_{\rm TERM}
 =\{(F_C,Z_C):F_C\in\mathcal L_{\scal},\ Z_C\in\mathbb R\}},
 \qquad
 F_C:=\mathfrak f(C),\quad Z_C:=\mathfrak f_{\scal}(C).
\]
After a curvature unit is fixed this product has numerical coordinates
in \(\mathbb R^2\), but its two factors are not dimensionally interchangeable:
\(F_C\) is curvature-valued and \(Z_C\) is dimensionless.
The combined source member is
\(\mathrm{TERM}_f^0:(F_C,Z_C)=(0,1)\), equivalently
\(C_{\rm sp}(C)=-C\) and \(C_{\rm sp}'(C)=0\) when
\(\mathfrak f(\scal)=\scal+C_{\rm sp}(\scal)\).
The two entries are independent endpoint data: \(F_C=0\) is a value
condition and \(Z_C=1\) is a first-derivative normalization.  Neither is
derived from the Hessian transport below, and neither supplies the other.
A vacuum constant-curvature endpoint additionally requires
\[
 \boxed{F_C=\frac{C Z_C}{2};}
\]
it is not implied by the source member.  The former
\((\mathrm{TERM}_H)\) Hamiltonian/Legendre narrative may motivate the source
choice but is not part of the mathematical boundary datum.

\item \emph{Minimal and enlarged metric-response realization families.}
The generated-action MRESP domain contains generated/co-realized constituents,
the certified CRP4 owner, irreducible endpoint data, and the derived
same-action quadratic owner.  Enlarging the parent by an independent
response or quadratic line restores the older calibration atlases.
Import the III-D exact response branch
\((\mathrm{RED})+({\rm UV})+({\rm IR})\) and the III-D \((\mathrm{NEUT}_7)\)
generated family of the III-C common-rate neutral protected semigroup.  The
complete visibility process, closed trace, normalized state, entropy
postprocessor, and determinant-depth coordinate are retained as a finite
generated-constituent presentation of the occurring response.  Add a frozen
decreasing member
\(\mathrm{REG}_{\sigma}\in\mathfrak A_{\rm REG}^{+}\),
\((\mathrm{CAL}_{\rm resp/EH})\), the response ancestry in
\((\mathrm{CRP4})\), and \((\mathrm{TERM}_f)\).
The generated neutral dual determinant-depth law and owner inheritance then
supply the scale interpolation
\[
 u(\tau):=e^{-\tau},\qquad
 q_{\rm sp}(\tau):=7(1-e^{-\tau}),\qquad
 \scal(\tau):=
 U\alpha_{\rm hier}(\sigma_{\rm fin})^{-q_{\rm sp}(\tau)},
\]
so \(\scal(0)=U\) and \(\scal(\tau)\downarrow C\).  On the protected
non-scalar block let
\[
 V_\tau|_W=e^{-\tau}I_W,
 \qquad A_\tau:=V_\tau^{1/2}|_W=e^{-\tau/2}I_W.
\]

The \emph{unmarked minimal response signature}
\((\mathrm{MRESP}_{\min})\) makes only the following type commitments.
The Hessian is an ordinary covariant tensor in
\(\operatorname{Sym}^2(W^*)\), its scale change is the tensor pullback along
the generated amplitude map \(A_\tau\), and no independent density character
or response line is inserted.  The scalar active count is a continuous
entropy-sufficient functional on the tensor-closed class of
finite-dimensional density matrices, is multiplicative under tensor products,
and equals eight on the maximally mixed protected octet.  Its value one on
pure states is then a theorem, not a premise.  Finally the scalar susceptibility is
positive and preserves the dimensionless homothety of
Theorem~\ref{p3e:thm:response-homothety}.  The classification theorems below
then force
\[
 \mathcal T_\tau(H)=A_\tau^*HA_\tau=e^{-\tau}H,
 \qquad
 w_{\rm nl}(\rho)=\frac{e^{S(\rho)}}8,
 \qquad
 \chi_{\rm UV}=\frac{\gamma_\sigma}{U},\quad \gamma_\sigma>0.
\]
Thus the unrestricted residual minimal atlas is only
\[
 \boxed{\mathfrak A_{\rm MRESP}^{\min}
 =\{\gamma_\sigma\in\mathbb R_{>0}\}},
\]
On the successfully evaluated same-action PQTI comparison the locus is
\[
 \boxed{\gamma_\sigma
 =16\sqrt{\alpha_{\rm hier}(\sigma_{\rm fin})},\qquad
 \chi_{\rm UV}
 =\frac{16\sqrt{\alpha_{\rm hier}(\sigma_{\rm fin})}}
        {\Lambda_{\rm EH}^2}.}
\]
The response Hessian is consequently
\[
 C_{\rm sp}''(\scal(\tau))
 =\frac{e^{S(\rho_\tau)}}8\,
   \frac{\gamma_\sigma}{U}\,e^{-\tau}.
\]
No first-variation transport is inferred from the Hessian pullback; at this
Hessian-only stage the terminal slope is represented by the independent datum
\(Z_C\) in \((\mathrm{TERM}_f)\).  Paper~III-F later goes one constructor
level upstream: the generated coefficient-family connection removes the
constant affine mode on the bare same-action branch and thereby derives
\(Z_C=1\).  If tensorial pullback, entropy sufficiency,
multiplicativity, or the unmarked-owner condition is dropped, the enlarged
positive/completely-positive transport and coupling atlas
\(\mathfrak T_+\times\mathfrak K_+\) reappears.  That enlarged atlas is not
used in the main response theorem.  The scale interpolation is inherited from
the certified response ancestry in \((\mathrm{CRP4})\) once a frozen
\(\mathrm{REG}_{\sigma}\in\mathfrak A_{\rm REG}^{+}\) is fixed.  The certified CRP4 port is
not a consequence of homothety or coefficient duality; it consumes the
evaluated complete actual-action residual. Likewise PQTI consumes the
actual normalized second jet rather than only its abstract tensor type.  The endpoint pair and an
actual solution/history remain outside these owner theorems; the REG member
is fixed only on the generated octahedral-interface branch.

\item \emph{Derived hierarchy-active state; legacy \((\mathrm{MON})\) retired.}
For the physical interval \(C\le\scal\le U\), inversion of the owner-typed
scale law in \((\mathrm{MRESP}_{\min})\) gives
\[
 u(\scal)=e^{-\tau_{\rm H}(\scal)}
 =\frac{\ln(\scal/C)}{\ln(U/C)}.
\]
The normalized active state
\[
 \sigma_{\rm act}(\scal)
 =\rho_{\tau_{\rm H}(\scal)}
 =\frac{P_0+u(\scal)P_W}{1+7u(\scal)}
\]
is a generated constituent of the occurring MRESP composite.  Paper-II-A
co-realization therefore supplies it without a separate occurrence
hypothesis.  Chronologically this is a \texttt{DERIVED-LATER} discharge of the
physical response-scale state left unresolved in Paper~III-D; it does not
identify the state with Paper~III-D's auxiliary secant-matched state.  Independent tensor-product composition, if later used, retains
the usual entropy additivity as a composition-domain statement; it is not
needed for this one-interval identification.
\end{enumerate}
\end{definition}

\begin{theorem}[Uniform ordered response bundle over the complete REG locus]
\label{p3e:thm:uniform-reg-response-bundle}
Let
\[
 I_+:=\left(-4\pi^2-\frac1{6\sqrt2},\,\infty\right),
 \qquad
 L_\sigma:=4\pi^2+\frac1{6\sqrt2}+\sigma,
 \qquad
 \alpha_\sigma:=e^{L_\sigma}.
\]
Fix the single CRP4 occurrence and one upper magnitude \(U>0\), but do not
choose a point of \(I_+\).  Then the complete
ordered response family is the smooth bundle
\[
 C_\sigma=U\alpha_\sigma^{-7},\qquad
 \scal_\sigma(\tau)
 =U\alpha_\sigma^{-7(1-e^{-\tau})},
 \qquad (\sigma,\tau)\in I_+\times[0,\infty).
\]
For every \(\sigma\in I_+\), simultaneously,
\[
 0<C_\sigma<\scal_\sigma(\tau)\le U,
 \quad
 \partial_\tau\scal_\sigma(\tau)
 =-7L_\sigma e^{-\tau}\scal_\sigma(\tau)<0,
 \quad
 \lim_{\tau\to\infty}\scal_\sigma(\tau)=C_\sigma,
\]
and the six finite hits are
\[
 \tau_j=-\log(1-j/7),\qquad
 \scal_\sigma(\tau_j)=U\alpha_\sigma^{-j},
 \qquad j=1,\ldots,6.
\]
Consequently every theorem whose conclusion is quantified over this bundle
(ordering, endpoint convergence, finite-hit counting, and the response
reconstruction for arbitrary positive \(\gamma_\sigma\)) requires no
selected REG member.  A numerical value of \(C/U\), or a point prediction at
one hierarchy scale, still requires the evaluation map
\(\operatorname{ev}_\sigma\); no theorem in the present corpus selects that
evaluation.  On this unanchored domain \(\mathrm{REG}_0\) is not selected. The occurring Hurwitz/Hopf branch instead imports Theorem~\ref{p3c:thm:occurring-hurwitz-reg-origin}.
\end{theorem}

\begin{proof}
The definition of \(I_+\) is exactly \(L_\sigma>0\), equivalently
\(\alpha_\sigma>1\), which Paper~III-D proves to be the complete ordered
owner-placement locus.  The displayed formulas are the generated neutral
character after the one owner port.  Positivity and strict ordering follow
from \(\alpha_\sigma>1\).  Differentiation gives the stated derivative, and
\(e^{-\tau}\to0\) gives the endpoint.  Solving
\(7(1-e^{-\tau})=j\) gives the six and only six finite integer solutions;
the value seven requires \(e^{-\tau}=0\) and hence occurs only in the limit.
Every calculation used \(\sigma\) as a free variable on \(I_+\), so the
result is a universal statement and contains no evaluation morphism.  Since
different \(\sigma\) give
\(C_\sigma/C_0=e^{-7\sigma}\), the family is nonconstant; therefore the
universal theorem cannot itself identify one member, including zero.
\end{proof}

\begin{corollary}[Owner-typed canonical active spectral profile]
\label{p3e:cor:owner-typed-active-spectral-profile}
On the complete ordered response bundle of
Theorem~\ref{p3e:thm:uniform-reg-response-bundle}, define
\[
 s_\sigma(\scal)
 :=\log_{\alpha_\sigma}\!\left(\frac U\scal\right)
 =\frac{\log(U/\scal)}{L_\sigma},
 \qquad C_\sigma\le\scal\le U.
\]
Then \(0\le s_\sigma(\scal)\le7\), and for every member of the unmarked
active orbit the response-scale pullback
\begin{equation}
\label{p3e:eq:owner-typed-active-spectral-profile}
 \boxed{
 P_{{\rm spec},\sigma}(\scal;u)
 :=\mathbf1_{[s_\sigma(\scal),7]}
       (\mathcal Q_{\rm act}(u))}
\end{equation}
is the canonical upper spectral filtration of Paper~III-C.  Its four exact
trace-defect anchors are
\begin{equation}
\label{p3e:eq:response-scale-active-spectral-anchors}
\begin{array}{c|c|c}
 d & \scal & P_{{\rm spec},\sigma}(\scal;u)\\ \hline
 0&U&\Pi_8\\
 4&U\alpha_\sigma^{-4}&\Pi_4\\
 5&U\alpha_\sigma^{-5}&\Pi_3(u)\\
 7&C_\sigma=U\alpha_\sigma^{-7}&P_0.
\end{array}
\end{equation}
The complete orbit-equivariant family is derived with no
\((\mathrm{ACT}_{\rm Cl})\) premise.  A particular rank-three projector still
requires a marked \(u\), and the physical scale formula still has exactly the
owner, REG, and upper-reference dependencies of the parent response bundle.
The family is stepwise: between the four anchors its projector codimension
need not equal \(s_\sigma(\scal)\).  It is neither the smooth visibility
operator \(V_\tau\), whose finite-time rank is always eight, nor an
operator-continuous rank-changing path or an actual transition token.
\end{corollary}

\begin{proof}
The response-bundle identity
\(\scal_\sigma(\tau)=U\alpha_\sigma^{-q_{\rm sp}(\tau)}\) gives
\[
 s_\sigma(\scal_\sigma(\tau))=q_{\rm sp}(\tau).
\]
Substitution in Paper~III-C,
Theorem~\ref{p3c:thm:common-trace-defect-spectral-filtration}, proves
\eqref{p3e:eq:owner-typed-active-spectral-profile}.  Evaluating at the four
spectral values \(0,4,5,7\) gives
\eqref{p3e:eq:response-scale-active-spectral-anchors}.  The owner and
occurrence firewalls are inherited unchanged from the two composed theorems.
\end{proof}

\begin{remark}[Ordered hierarchy refinements]
The smooth protected rank-trace coordinate crosses only the first six integer
values at finite flow parameter; its seventh value is the
\(\tau\to\infty\) endpoint.  On the
\(\mathfrak A_{\rm REG}^{+}\) branch and through the proved owner-ancestry
port, these are ordered scalar response surfaces without
\((\mathrm{FLAG})\).  A named
cosmological process at a surface still requires its downstream field-equation
predicate.  If \((\mathrm{FLAG})\) and its sharp intertwining data are
supplied, the same ledger may additionally index successive rank-one
quotients, but the seventh smooth-semigroup entry remains an asymptotic
endpoint rather than a finite crossing.
\end{remark}

\label{p3:part:dyn-grav}
\label{p3:part:gravity-master-action}

\subsection{Same-root graded ancestry and determinant-owner boundary}
\label{p3e:sec:graded-owner-update}

Paper~III-B now supplies an orientation-preserving spin lift of the generated
diagonal ratio-fibre projection.  Here write
$W_{\mathrm{spin}}:=W_L\oplus W_R$ for its complex half-spin target;
this is \emph{not} the seven-real-dimensional trace-free Clifford response
block $W_7:=\Cl(0,3)\ominus\mathbb R1$ used in the response Hessians.
On the protected spin constituents the transport gives
\[
 S_L\longrightarrow W_L,
 \qquad
 S_R\longrightarrow W_R.
\]
The present paper independently gives these same protected constituents their
branchwise local Lorentz/null-spinor role.  Therefore the relative grading is
a same-root ancestry relation rather than an independent bridge.

\begin{corollary}[Even external/internal grading ancestry]
\label{p3e:cor:even-external-internal-grading-ancestry}
Let $S_x=S_L\oplus S_R$ be the external spin carrier and let
$J_\Delta:S_x\to W_{\mathrm{spin}}=W_L\oplus W_R$ be the composition
of the local Lorentz-spinor ancestry with the Paper~III-B
orientation-preserving diagonal-fibre spin transport.  With
$\gamma_x|_{S_{L/R}}=\mp1$ and
$\Gamma_{W_{\mathrm{spin}}}|_{W_{L/R}}=\mp1$, respectively,
\[
 \boxed{J_\Delta\gamma_x=\Gamma_{W_{\mathrm{spin}}}J_\Delta.}
\]
No grading on the distinct real seven-dimensional response block $W_7$
is implied by this spin identity.
An odd relation would require an additional orientation reversal on exactly
one owner and is not generated by the diagonal graph.
\end{corollary}

Separately, if $L_1,L_2$ are independently constituted oriented positive
one-dimensional owners, the positive isomorphisms $L_1\to L_2$ form an
$\RR_{>0}$ torsor.  Equal dimension, positivity, or determinant weight cannot
select one.  The response-root/determinant bridge and the remaining
susceptibility normalization can therefore close only after both roles are
proved to be associated weight lines of one generated determinant-owner
functor (or after a canonical isometric normalization is generated).  This
retypes the residual as owner ancestry rather than a numerical coincidence.

\begin{proposition}[Exact C12 scheme--depth closure criterion]
\label{p3e:prop:c12-scheme-depth-closure-criterion}
Fix $\mathrm{REG}_\sigma\in\mathfrak A_{\rm REG}^{+}$ and write
$L_\sigma=\log\alpha_\sigma>0$.  The complete positive response-character
counterfamily is
\begin{equation}
 \label{p3e:eq:c12-character-counterfamily}
 \frac{\mathcal R_\lambda(q)}{U}
 =e^{-\lambda L_\sigma q},\qquad
 C_\lambda=Ue^{-7\lambda L_\sigma},\qquad \lambda>0.
\end{equation}
Every member has the same normalized hierarchy parameter
\[
 \frac{\log(\mathcal R_\lambda(q)/C_\lambda)}
      {\log(U/C_\lambda)}=1-\frac q7,
\]
and therefore preserves the MRESP visibility parameter, active state,
entropy law, covariant Hessian attenuation, endpoint ordering, and normalized
response reconstruction.  Those structures cannot select $\lambda$.

By Paper~III-C,
Proposition~\ref{p3c:prop:full-scheme-depth-determinant-connection}, every
weighted member also carries a flat connection and is connection-isomorphic
to the fundamental determinant-dual member on the contractible
scheme--depth base.  Hence neither a pointwise positive-line isomorphism nor
flat connection preservation closes C12.

Within this character family, the weight-one relation
\(
 \lambda_{RD}=1
\)
is forced exactly when the physical response comparison is induced before
trivialization by the same fundamental determinant-dual associated
representation, with no affine-origin choice or depth-dependent gauge
deformation.  On that branch the identity intertwiner gives
\[
 \frac{\mathcal R(q)}U=e^{-L_\sigma q},
 \qquad
 \frac CU=e^{-7L_\sigma}=\alpha_\sigma^{-7}.
\]
Neither post-response MRESP nor the scalar evaluation pairing alone
constructs the contact/action associated-functor morphism. Paper~III-B's
actual-action derivative theorem
\ref{p3b:thm:v2211-action-contact-certificate} supplies the exact test:
the complete coefficient must have $\eta_U=0$, with a generated normalized
port. Thus C12 closes on the successfully evaluated $(\mathrm{ACD}_H)$
domain; until then its physical status is an unresolved action/owner test.
The neutral contact character and the conditional representation theorem
remain unchanged. Paper~III-G, Appendix~\ref{p3e:app:c12-flat-gauge-countermodel}
retains the counterfamily.
\end{proposition}

\begin{proof}
Substitution in the displayed logarithmic quotient cancels
$\lambda L_\sigma$.  The visibility parameter is consequently $1-q/7$ for
every $\lambda>0$, and all listed normalized descendants depend on that
parameter rather than on $\lambda$ separately.  The connection statements
and affine-origin test are precisely Paper~III-C,
Proposition~\ref{p3c:prop:full-scheme-depth-determinant-connection}.  A
same-representation morphism has the fundamental dual weight by construction,
so it gives $\lambda=1$; conversely $\lambda=1$ admits the identity
intertwiner.  The independently typed response atlas supplies no such
morphism, which proves its scope boundary; the action-density duality fixes only its own coefficient slot. The
actual-action residual and normalized origin-certified port supply the
stronger comparison exactly when their evaluated certificate succeeds.
\end{proof}

\begin{definition}[CRP4: certified actual-action/contact parent]
\label{p3e:def:crp4-cross-role-positive-parent}
Let $\mathcal D_{\rm act}$ be the actual scalar-action density line of
Paper~III-B, so its coefficient line is
$\mathcal L_U=\mathcal D_{\rm act}^{-1}$. Let
$\mathcal L_H^{-1}$ be Paper~III-C's fundamental contact determinant-dual
line over $B=I_+\times[0,7]$, with transport
$H_{b_1b_0}=e^{-[\Phi(b_1)-\Phi(b_0)]}$.
The tag $(\mathrm{CRP4})$ certifies an actual generated normalized port
\[
 a_U:\mathcal L_U\longrightarrow\mathcal L_H^{-1},\qquad
 a_U(b_1)T^U_{b_1b_0}=H_{b_1b_0}a_U(b_0),
\]
including the operator/density/mark origin, root comparison and the successful
$(\mathrm{ACD}_H)$ test of
Theorem~\ref{p3b:thm:v2211-action-contact-certificate}. In its retained
compatible presentation it may be written $a_U=\operatorname{id}$;
that expression records the certified comparison and does not prove it.
The historical aliases $(\mathrm{OWN}_{56})$ and
$(\mathrm{OWN}_{\rm resp/det})$ refer to this same port. On the fixed
upper section $U=\Lambda_{\rm EH}^2$, the rank-four action and root below
are derived descendants of the certified character. Without the port the
contact hierarchy is neutral and physical weight-one placement remains open.
\end{definition}

\begin{theorem}[Uniqueness and connection of the CRP4 action]
\label{p3e:thm:crp4-uniqueness}
Let \(M\) be a local operational spacetime patch, write
\(\pi_B:B\times M\to B\) and \(\pi_M:B\times M\to M\), and let
\(\mathscr W^*\to M\) be Paper~III-A's rank-four causal generator bundle.
On
\[
 \mathscr P_4
 :=\pi_B^*\mathcal L_U\oplus\pi_M^*\mathscr W^*
 \longrightarrow B\times M,
\]
the certified CRP4 ancestry generates a unique continuous unital
\emph{$B$-direction} cocycle of typed, summand-preserving actions at fixed
$x\in M$, namely
\begin{equation}
 \label{p3e:eq:crp4-action}
 A_{b_1b_0,x}\bigl(r;(E,p)\bigr)
 =\bigl(H_{b_1b_0}r;
        H_{b_1b_0}^{1/4}(E,p)\bigr).
\end{equation}
Consequently the induced \emph{relative $B$-connection} on
$\mathscr P_4$ in the displayed positive trivializations is
\[
 \nabla_B^{\mathrm{CRP4}}
 =\left(d_B+d_B\Phi\right)
  \oplus\left(d_B+\frac14d_B\Phi\right),
\]
and the fourth exterior density of the generator block has fundamental
$B$-transport $H$.  This certificate does not uniquely determine a full
connection along $B\times M$: an $M$-direction connection on
$\mathscr W^*$ must be independently inherited or constructed.  If such a
connection $\nabla^M$ is supplied, its pullback can be added in the
$M$-directions without changing the displayed CRP4 $B$-transport.
\end{theorem}

\begin{proof}
By the successfully evaluated port in
Definition~\ref{p3e:def:crp4-cross-role-positive-parent}, $a_U$
transports the actual response section by the fundamental character $H$.
Paper~III-A proves that determinant is an isomorphism from the identity
component of the \(O(3)\)-equivariant causal-cone automorphism group to
\(\mathbb R_{>0}\), with inverse \(H\mapsto H^{1/4}I_4\).  Paper~III-C,
Theorem~\ref{p3c:thm:positive-associated-line-fourth-root}, independently
identifies the unique positive associated fourth root and its one-quarter
connection.  Applying those theorem-level operations to the character carried
by \(a_U\) gives \eqref{p3e:eq:crp4-action}.  The positive fourth-root multiplicativity and scalar determinant
\(a^4\) are the shared algebra of Paper~I-C and Paper~I-D,
Propositions~\ref{p1b:prop:trace-coordinate}
and~\ref{p1b:prop:local-coframe}~\cite{Kugelmann2026Paper1C,Kugelmann2026Paper1D}.
Applied to the two summands constructed here, they give continuity,
the unit, the cocycle law, cone preservation, determinant compatibility
and identity-component membership.  Finally
\(H_{b_1b_0}=e^{-(\Phi(b_1)-\Phi(b_0))}\) gives the displayed
partial $B$-connection.  Distinct $M$-connections share that same
$B$-transport, so full connection uniqueness does not follow.
\end{proof}

\begin{corollary}[Certified actual-action port closes C12 on its evaluated domain]
\label{p3e:cor:crp4-c12-discharge}
On $\mathrm{REG}_\sigma\in\mathfrak A_{\rm REG}^{+}$, assume the
actual normalized $(\mathrm{ACD}_H)$ comparison certified by
$(\mathrm{CRP4})$. Then
\[
 \boxed{\lambda_{RD}=1,\qquad
 \mathcal R(q)/U=e^{-qL_\sigma},\qquad
 C/U=e^{-7L_\sigma}=\alpha_\sigma^{-7}.}
\]
The physical conclusion is conditional on that evaluated action/owner port.
The initial coefficient duality, M8 geometry, flatness or numerical equality
alone do not close C12. The scalar heat counterfamily in
Corollary~\ref{p3b:cor:v2211-heat-contact-counterfamily} shows the residual
can be nonzero even with the same upper coefficient and process profile.
\end{corollary}
\begin{proof}
The actual-action residual theorem integrates $\eta_U=0$ with its normalized
root to the fundamental character. Restrict to fixed scheme and evaluate at
$q=7$. The prior counterfamilies prove the necessity of the additional
physical source/owner evaluation.
\end{proof}

\begin{corollary}[Common hierarchy as normalized response and energy ratios]
\label{p3e:cor:v179-response-energy-quarter-character}
On the same-owner CRP4 branch of
Corollary~\ref{p3e:cor:crp4-c12-discharge}, represent the positive
quarter-character with reference \(E_0>0\) by
\(E(8q)/E_0=\alpha_{\rm hier}^{-q/4}\).  Then
\[
 \frac{\mathcal R(q)}U
   =\left(\frac{E(8q)}{E_0}\right)^4,\qquad
 \frac CU=\left(\frac{E_{56}}{E_0}\right)^4 .
\]
\end{corollary}
\begin{proof}
The scalar-action response ratio is \(\alpha_{\rm hier}^{-q}\);
the positive associated fourth root has ratio
\(\alpha_{\rm hier}^{-q/4}\).  Raise that ratio to the fourth power and
specialize \(q=7\).
\end{proof}

These dimensionless identities explain the common determinant ancestry
of cosmological response depths and the first-order hierarchy used by
the finite spectrum.  Each dimensionful quantity retains its own unit
owner.  Neither an energy-to-the-fourth dimensional equality with a
curvature coordinate nor a particle-row assignment is inferred.
The explicit Clifford transition of
Theorem~\ref{p3c:thm:v179-balanced-hierarchy-response} uses this
quarter-character without adding an active midpoint projector.


\subsection{owner-normal-form reductions}
\label{p3e:sec:post-v125-owner-normal-forms}

This revision sharpens three residual scalar-response interfaces without
altering the generated response graph.

\subparagraph{Response/determinant character.}
Before any owner bridge is imposed, the MRESP coordinate already gives
\[
 \frac{\mathcal R(q)}{U}=e^{-\ell_R q},\qquad
 \ell_R:=\frac1{7}\log\frac{U}{C}.
\]
At fixed $\mathrm{REG}_\sigma$, the neutral determinant dual instead carries
$e^{-L_\sigma q}$, $L_\sigma=\log\alpha_\sigma$.  Thus the only invariant
comparison is
\[
 \lambda_{RD}:=\frac{\ell_R}{L_\sigma}
 =\frac{\log(U/C)}{7\log\alpha_\sigma}.
\]
The weight-one owner statement is exactly $\lambda_{RD}=1$.  For every
$\lambda>0$, however, the ordered roots
$C_\lambda=U\alpha_\sigma^{-7\lambda}$ satisfy the same MRESP construction and
give $\mathcal R(q)/U=\alpha_\sigma^{-\lambda q}$.  Hence a post-response
MRESP classification by itself does not select the weight-one determinant
ancestry.  Definition~\ref{p3e:def:crp4-cross-role-positive-parent} records
the earlier scalar-action coefficient owner needed for a valid closure, and
Corollary~\ref{p3e:cor:crp4-c12-discharge} then derives \(\lambda_{RD}=1\).
A pointwise line identification or connection-preserving gauge equivalence
would remain insufficient: the latter exists for every positive weight by
Proposition~\ref{p3e:prop:c12-scheme-depth-closure-criterion}.

\subparagraph{Ultraviolet susceptibility.}
The homothety theorem fixes only
\[
 \chi_{\rm UV}(U,\sigma)=\frac{\gamma_\sigma}{U},\qquad \gamma_\sigma>0.
\]
For every positive $g$, replacing $\gamma_\sigma$ by $g$ leaves the generated
visibility, entropy coupling, tensorial Hessian pullback, convexity, terminal
Hessian decay and response reconstruction intact.  Under the same-action
positive rescaling of Paper~III-F, the invariant coordinate is
\[
 \widehat\gamma_\sigma:=\gamma_\sigma/Z_C.
\]
Thus MRESP and homothety alone leave a one-parameter counterfamily.  PQTI is
the generated same-action relation that removes this freedom by tying the
unit primitive response to the tensor square of the CRP4 generator and by
using the already derived octet/scalar trace ratio.

\begin{proposition}[PQTI susceptibility law and invariant amplitude]
\label{p3e:prop:conditional-susceptibility-half-weight}
On \((\mathrm{CRP4})+(\mathrm{PQTI})\), the former susceptibility-owner
symbol \(a_\chi\) is the composite generated by \(b_Q\) and the unique
positive half-root functor, rather than an independent ancestry premise.
Moreover
\[
 \boxed{\widehat\gamma_\sigma
       =16\sqrt{\alpha_\sigma},\qquad
 A_\chi:=\widehat\gamma_\sigma\alpha_\sigma^{-1/2}
       =16.}
\]
The value \(A_\chi=16\) is therefore a generated invariant on this branch,
not an additional root.
\end{proposition}

\begin{proof}
Theorem~\ref{p3e:thm:pqti-quadratic-normalization} gives the invariant law
directly.  Dividing by \(\sqrt{\alpha_\sigma}\) gives
\(A_\chi=16\).  The counterfamily in
Paper~III-G, Appendix~\ref{p3e:app:remaining-owner-counterfamilies} remains a proof that
MRESP without PQTI does not determine this quantity; it is not a
counterexample to the PQTI branch.
\end{proof}

\subparagraph{Response/EH placement.}
The dimensionless ratio
\[
 \kappa_U:=\frac{U}{\Lambda_{\rm EH}^2}
\]
is a relation between two already constituted positive curvature-scale
sections, not a unit convention.  Paper~III-B
Theorem~\ref{p3b:thm:generated-action-response-owner} supplies only the
same-action identity-stage section:
\[
 \boxed{\kappa_U=1.}
\]
The response homothety still gives
$U_\kappa=\kappa\Lambda_{\rm EH}^2$ on the unrestricted
independent-response atlas, with
\(\gamma_\kappa=16\kappa\sqrt{\alpha_\sigma}\).  Those members independently
normalize the response section and therefore change the frozen action parent.
The family is recorded in
Paper~III-G, Appendix~\ref{p3e:app:remaining-owner-counterfamilies} as a mutation control,
not as a countermodel to the generated branch.


\subsection{Response-atlas decomposition}
\label{p3e:sec:v236-mresp-atlas}
\begin{proposition}[Minimal MRESP is not one monolithic occurrence root]
\label{p3e:prop:v236-mresp-decomposition}
On \((\mathrm{MRESP}_{\min})\), the tensorial Hessian transport, entropy
sufficiency and response homothety are structural consequences of the declared
minimal response grammar.  The remaining physical interfaces decompose into:
\begin{enumerate}[label=(\roman*),leftmargin=2.2em]
\item a same-action/owner port identifying the physical response Hessian with
the Hessian (or stationary Schur descendant) of an already owned action;
\item the actual background/source member on which that response is evaluated;
\item the positive scale calibration \(\gamma_\sigma\);
\item a REG point/member only for claims requiring a numerical finite-part
value; and
\item terminal/boundary members used by the particular readout.
\end{enumerate}
Thus an occurring physical response need not be counted again for every
structural transport law generated inside \((\mathrm{MRESP}_{\min})\).
If the same-action port is later proved, the Hessian component is
\(\mathrm{DL}\)/\(\mathrm{GEN}\), while the actual source/background and
calibration retain their own statuses.
\end{proposition}


\subsection{Current K16 regularity and scalar/tensor interface (rank four)}
\label{p3e:sec:v2219-k16-regular}
On the historical $n_4\mid n_7$ response comparator the point $C_2$ is strictly inside
an analytic branch of the coefficient. Consequently $[F]_{C_2}=0$,
$[\chi]_{C_2}=0$, and $[\partial_R\chi]_{C_2}=0$. A source-free
regular crossing satisfies $[\nabla_n F]_{C_2}=0$ and, because
$\chi(C_2)>0$, also $[\nabla_n R]_{C_2}=0$.  It is not a physical
restriction on the separate oscillating UV scalaron trajectory.
The old depth-two K16 finite-third-jet statement in the following section
is retained only as the previous comparison theorem.


\paragraph{Technical proof ownership and constituted scope.}
Every extended auxiliary derivation, explicit matrix verification,
and optional comparison moved during this reorganization is
preserved in Paper~III-G, Section~\ref{p3g:sec:gravity}.
The physical momentum carrier is inherited from Paper~II-B:
$\mathcal H_p\cong\mathcal H_1$ and $m_p=1$, without later
scalar, action, regulator, or source data reselecting it.

\section{Geometric gravitational action, signature, and normalization}
\label{paper3e:sec:organized-gravitational}
\subsection{Derivation of the Gravitational Lagrangian}
\label{p3:sec:lagrangian}

Table~\ref{tab:p3-gravity-ledger} records the gravitational-dynamics results discharged in this part.

\begingroup
\small
\setlength{\tabcolsep}{3pt}
\renewcommand{\arraystretch}{1.15}
\begin{longtable}{@{}
  L{3.5cm}
  L{8.5cm}
  L{\dimexpr\linewidth-12.0cm-4\tabcolsep\relax}
@{}}
\caption[Gravitational dynamics ledger]{Gravitational dynamics ledger.}
\label{tab:p3-gravity-ledger}\\
\toprule
\textbf{Ledger row} & \textbf{Claim} & \textbf{Discharged by} \\
\midrule
\endfirsthead
\multicolumn{3}{@{}l}{\small\emph{Gravitational dynamics ledger (continued)}}\\
\toprule
\textbf{Ledger row} & \textbf{Claim} & \textbf{Discharged by} \\
\midrule
\endhead
\midrule
\multicolumn{3}{r@{}}{\small\emph{Continued on next page}}\\
\endfoot
\bottomrule
\endlastfoot
Einstein--Hilbert channel
& The gravitational scalar channel gives the \(a_2\) Einstein--Hilbert term.
& Proposition~\ref{p3:prop:EH_from_a2} \\

Fixed scalar weights
& Ambient-octet provenance and the Born diagonal give
\(w_{\rm vac}=w_{\rm EH}=1/8\); the bare vacuum coefficient is not an
observed-dark-energy determination.
& Theorem~\ref{p3e:thm:fixed-scalar-and-active-weights} \\

Generated gravitational heat closure
& The least unmarked local second-order unit-preserving resolution role fixes
\(F_g(u)=e^{-u}\), \(f_{2,g}=(4\pi)^{-2}\), and
\(\Lambda_g^2=96\pi^2\Lambda_{\rm EH}^2\), without a separate process
occurrence; the conditional cell-port rewrite is \(\Rp:=U=(3/\pi)\mathcal K_p\), with \(\mathcal K_p=6/\Rmom^2\),
while the raw joint coefficient separately satisfies
\(\UEHjoint^{\rm raw}=96U\).
& Theorem~\ref{p3e:thm:generated-gravitational-heat-role},
Corollary~\ref{p3e:cor:grav-heat-branch} \\

Optional profile extension
& If the locality/leastness signature is enlarged to an arbitrary positive
profile, the scalar-action compiler fixes only its moment--scale class;
Newton-coordinate comparison adds no selector.
& Theorem~\ref{p3:thm:EH-normalization-gauge} \\

Lorentzian sign chain
& The III-A clocked-Hopf composite and its least generated doctrine fix the
sharp operational cone and Lorentzian conformal
sign; the declared passive local second-order scalar realization gives the
finite massive mode relation.
& Theorem~\ref{p3:thm:lorentzian-signature},
Theorem~\ref{p3:thm:p3-p4-lorentzian-chain-closure} \\

Celestial/null-direction bridge
& The Paper-II chiral projective sphere is identified with the projectivized
future null boundary of the independently derived local causal cone.
& Paper~II-E, Theorem~\ref{p2e:thm:celestial-hermitian-ray-bridge},
Corollary~\ref{p3:cor:chiral-projective-null-directions} \\

Conformal cutoff no-go
& Smooth conformal compactification cannot preserve the finite spectral cutoff.
& Theorem~\ref{p3:thm:conformal-compactification-cutoff} \\

Imported reduced response
& Paper~III-D supplies the exact two-root coupling, Riccati identity, secant index, and entropy-matched protected state.
& Paper~III-D, Theorem~\ref{p3:thm:geff_from_reciprocity},
Theorem~\ref{p3:thm:riccati}, Proposition~\ref{p3:prop:n_closed} \\

Physical shell transport
& The occurring MRESP composite co-realizes
\(\sigma_{\rm act}(\scal)=\rho_{\tau_{\rm H}(\scal)}\); hence
\(N_{\rm eff}=8^{n_{\rm act}}\) and
\(w_{\rm nl}=(1/8)8^{n_{\rm act}}\), with no MON occurrence premise.  Its
minimal entropy coupling is unique.
& Theorem~\ref{p3:thm:index-locking},
Theorem~\ref{p3e:thm:entropy-effective-count-uniqueness} \\

Canonical support filtration
& The common trace defect pulls the active-depth upper spectral family through
the ordered owner bundle, giving four exact scale/support anchors without
\(({\rm ACT}_{\rm Cl})\); the profile is stepwise and orbit-valued.
& Corollary~\ref{p3e:cor:owner-typed-active-spectral-profile} \\

Smooth master response
& Under \(({\rm MRESP}_{\min})\), tensorial Hessian pullback and the entropy
law are derived; arbitrary endpoint data and one positive dimensionless UV
calibration reconstruct the complete convex family, whose graph inherits the
joint input witness.
& Theorem~\ref{p3e:thm:complete-minimal-smooth-response-family},
Corollary~\ref{p3e:cor:mresp-composite-witness-transport} \\

Terminal generated response
& Under the generated response/EH identity, owner-ancestry port, and
\(\mathrm{REG}_{\sigma}\in\mathfrak A_{\rm REG}^{+}\),
\(C=U\alpha_{\rm hier}(\sigma_{\rm fin})^{-7}\) with
\(\mathcal K_p=6/\Rmom^2\), \(U=\Lambda_{\rm EH}^2\), with $U=(3/\pi)\mathcal K_p$ only after the actual cell test. The conditional smooth branch has a
decoupling scalaron there.
& Definition~\ref{p3:def:causal-wick-response-hypotheses}, Theorem~\ref{p3:thm:scalar-threshold} \\
\end{longtable}
\TableNotes{Each row states its own dependency boundary.  The heat-kernel
\(a_2\) identity is standard inside the declared scalar channel; causal
signature is the derived \(({\rm CAUS})\) domain.  Any additional
Euclidean-correlation reconstruction claim not proved inside Paper~III-E uses
the claim-scoped \(({\rm WICK}_\bullet)\) tag appropriate to that conclusion;
Paper~IV later supplies finite certificates on stated narrower domains.  The
nonlinear response requires \(({\rm MRESP}_{\min})\), and the single-copy
Einstein--Hilbert coefficient uses the derived Born-orbit pullback certified
by \((\mathrm{BD}_{\rm EH})\).}
\endgroup

\subsubsection{The Einstein--Hilbert Term}
\label{p3:sec:EH_term}

\begin{definition}[Generated ambient gravitational role and exact-readout domain]
\label{p3e:def:grav-closed-occurrence}
On each constituted reciprocal scalar metric parent
\((\mathcal H_s,g_s)\), \(s\in\{x,p\}\), let
\[
 L_{g,s}:=\Delta_{+,g_s}^{\rm scal},\qquad
 A_{g,s}(\Lambda_g):=
 e^{-L_{g,s}/\Lambda_g^2},\qquad \Lambda_g>0.
\]
The family is formed on the ambient regular Clifford octet before the derived
scalar projection.  It is transported by the constituted Born map and has no
external state, source, field, grading, projector, or channel insertion.

The historical tag
\[
 (\mathrm{OCC}_{g}^{\rm amb})
\]
and its legacy synonym \((\mathrm{GRAV}_{\circ})\) now certify the generated
role proved in Theorem~\ref{p3e:thm:generated-gravitational-heat-role}; they
do not declare an arbitrary profile occurrence.  On a finite window or compact smooth parent, its ordinary scalar closure
\[
 r_{g,s}:=\operatorname{Tr}_{\mathcal H_s}A_{g,s}:
 \mathbf1\to\mathbf1
\]
satisfies \(\operatorname{FSR}_{\Sigma}(r_{g,s})\), and its closing
provenance is unmarked \emph{on the ordinary trace-class branch}.
Trace class is automatic on a finite window and, for $\Lambda_g>0$,
on a compact smooth parent.  A noncompact parent needs a separately supplied
localized, relative or trace-per-volume functional $\tau_{g,\mathrm{rel}}$;
its finite value $\tau_{g,\mathrm{rel}}(A_{g,s})$ does not entail
$A_{g,s}\in\mathcal L^1(\mathcal H_s)$ or an ordinary full-carrier trace.  A different positive profile
\(F_g\ne e^{-u}\) is an optional enlarged process root and is not denoted by
\((\mathrm{GRAV}_{\circ})\) in the premise-minimal branch.
\end{definition}

\begin{definition}[Smooth gravitational expansion domain]
\label{p3e:def:grav-smooth-expansion-domain}
The tag \((\mathrm{GRAV\!-\!ASYM})\) supplements the generated heat role with
the following sufficient domain for the expansion used below: each smooth
parent is a compact boundaryless four-dimensional Riemannian manifold and
\(L_{g,s}=\Delta_{+,g_s}^{\rm scal}\) has its inherited smooth metric
principal symbol.  The generated profile \(e^{-u}\) is already nonnegative,
smooth, and rapidly decreasing with all required moments finite.  Manifolds with
boundary require separately stated elliptic boundary conditions and their
boundary coefficients.  This domain is unnecessary for the exact closed
trace and is never inferred from generated-role co-realization alone.
\end{definition}

\begin{theorem}[Born-orbit natural scalar-density identification]
\label{p3e:thm:born-orbit-eh-identification}
Let \(\beta:\mathcal M_x\to\mathcal M_p\) be the constituted smooth
common-model Born identification on the reciprocal scalar parents.  Let
\(\ell(g)\) be any diffeomorphism-natural local scalar built only from the
metric and its Levi--Civita curvature data, and define
\[
 \mathcal I_{s,\ell}(g_s):=\int_{\mathcal M_s}
 \ell(g_s)\,d\mu_{g_s},
 \qquad s\in\{x,p\},
\]
and their invariant average
\[
 \mathcal I_{{\rm B},\ell}(g_x,g_p)
 :=\frac12\mathcal I_{x,\ell}(g_x)
   +\frac12\mathcal I_{p,\ell}(g_p).
\]
On the generated Born-diagonal embedding
\[
 j_{\rm B}:\mathscr G_x\longrightarrow\mathscr G_x\times\mathscr G_p,
 \qquad
 j_{\rm B}(g):=(g,\beta_*g),
\]
one has
\[
\boxed{j_{\rm B}^*\mathcal I_{{\rm B},\ell}(g)
=\mathcal I_{x,\ell}(g).}
\]
Thus the two half-weighted representatives give one full natural scalar-density
copy on the Born diagonal.  The theorem applies in particular to
\(\ell=1\), \(\ell=\scal\), every local pure-metric heat coefficient, and a
declared local \(\ell=f(\scal)\).  This identity is differential-geometric and
requires no process occurrence.  On \((\mathrm{GRAV\!-\!ASYM})\), the
generated gravitational heat role contains the corresponding \(a_0,a_2,\ldots\)
constituents, so the same identity is their physical Born-diagonal pullback.
The legacy tag \((\mathrm{BD}_{\rm EH})\) records the
\(\ell=\scal\) specialization; it is not an additional identification
hypothesis.
\end{theorem}

\begin{proof}
The scalar \(\ell\) and Riemannian volume density are natural under
diffeomorphism:
\[
 \ell(\beta_*g)=\ell(g)\circ\beta^{-1},
 \qquad
 d\mu_{\beta_*g}=\beta_*d\mu_g.
\]
Change of variables therefore gives
\(\mathcal I_{p,\ell}(\beta_*g)=\mathcal I_{x,\ell}(g)\).  Substitution into the Born
average yields
\[
 j_{\rm B}^*\mathcal I_{{\rm B},\ell}(g)
 =\tfrac12\mathcal I_{x,\ell}(g)+\tfrac12\mathcal I_{x,\ell}(g)
 =\mathcal I_{x,\ell}(g).
\]
This proves the mathematical identity.  On the additional
\((\mathrm{GRAV\!-\!ASYM})\) domain, the two sector terms are finite
nonredundant Born-orbit constituents of the generated composite, so Paper~II-A
Theorem~\ref{p2:thm:generated-constituent-corealization} retains both.
Their sector labels are exchanged by Born parity and their common diagonal
signature identifies the orbit average.  No equality is claimed for
arbitrary off-diagonal pairs \((g_x,g_p)\), and no equality of the unequal
finite sector matrix traces is used.
\end{proof}

In particular \(\beta^*R_{\beta_*g}=R_g\), with the corresponding
natural density pullback. This geometric identity is distinct from the
numerical response inversion \(R\mapsto CU/R\). Thus naturality of
the gravitational density alone does not assign a weighted odd action
coefficient or an even scalaron-potential role under that numerical map.


\begin{mdframed}[nobreak=false]
\noindent\textbf{Fixed scalar-corner domain.}
The physical protected scalar-corner calculation in this subsection uses
\[
 \boxed{\mathsf C_{\rm IR}^{\rm der}
 +({\rm GRAV\!-\!ASYM}).}
\]
The finite protected block and scalar infrared remnant are upstream derived
outputs.  Paper~III-B's co-realization theorem derives the ambient rank-one
coefficient certified by \((\mathrm{CORN})\), while
Theorem~\ref{p3e:thm:born-orbit-eh-identification} derives the diagonal
full-copy statement certified by \((\mathrm{BD}_{\rm EH})\).  Neither tag is
an extra antecedent.  The labels \(({\rm GRAV}_{\circ})\) and
\(({\rm GRAV\!-\!HEAT})\) certify the generated process theorem and are not
extra antecedents either.  The displayed \(a_0\) and \(a_2\) coefficients are
smooth-parent semiclassical/heat-kernel extractions of the generated process,
not an exact finite-matrix heat-kernel coefficient.
\end{mdframed}

Use exactly the two-stage trace interface of Paper~III-B,
Definition~\ref{p3b:def:sector-trace-born-state}.  Ordinary carrier traces
first form \(\mathfrak T_{\rm sec}(A_x\oplus A_p)\in Z_{\rm role}\), and
only then does the normalized reciprocal role scalarizer \(\lambda_{\rm B}\) act on the two
role labels.  Its equal reciprocal weights are inherited from Paper~II-A through the III-B scalar-action owner, while its unit normalization fixes their common magnitude.  Thus the
reduced sectorwise physical action is the extensive composite
\(\lambda_{\rm B}\circ\mathfrak T_{\rm sec}\):
\begin{equation}
S_{\mathrm{grav}}[\Lambda_g]
=\frac12\,\Tr_{\mathcal H_x}\!\left(e^{-D_x^2/\Lambda_g^2}\right)
+\frac12\,\Tr_{\mathcal H_p}\!\left(e^{-D_p^2/\Lambda_g^2}\right),
\label{p3:eq:Sphys_optionA_recall}
\end{equation}
with one transported generated heat law.  The half weights belong to the normalized role scalarizer on \(Z_{\rm role}\), not to a prepared state and not to the unequal finite matrix blocks; on
the current \(5\oplus15\) seed the same extensive composite evaluates the two
central block identities as \(5/2\) and \(15/2\).

\begin{remark}[Exact finite action versus smooth Born-diagonal extraction]
\label{p3e:rem:v2211-finite-smooth-action-port}
The Born-diagonal density theorem compares two transported smooth metric
parents on $(\mathrm{GRAV\!-\!ASYM})$. It does not equate their traces
with the literal fixed momentum matrix trace on $m_p=1$. The controlled
joint-window theorem of Paper~III-B licenses the position/parent comparison;
no physical $m_p\to\infty$ limit is used. An export of the full-copy smooth
EH coefficient to the exact finite physical action therefore requires the
actual operator/effective-action coefficient port, including any scalar-circle
metric, localizer and normalization. Without that evaluated port the formulas
below are exact coefficients of the stated smooth-parent comparison, not
strict heat coefficients of the fixed finite matrix.
Section~\ref{p3e:sec:v2218-fixed-momentum-action-port} evaluates two direct
fixed-momentum promotions and the complete shared-cutoff heat comparison.
The resulting logarithmic and clock identities are executable tests of a
proposed port, rather than an identification of the finite momentum trace
with a four-dimensional metric density.
\end{remark}

\begin{definition}[Gravitational process moment \texorpdfstring{$f_{2,g}$}{f2g}]
\label{p3:def:f2_rescaled}
On the premise-minimal generated branch, define the dimensionless second
moment by
\begin{equation}
f_{2,g} \;:=\; \frac{1}{(4\pi)^2}\int_{0}^{\infty} e^{-u}\,du
=\frac1{(4\pi)^2}.
\label{p3:eq:f2_def}
\end{equation}
This convention absorbs the universal $(4\pi)^{-2}$ factor that appears in the $4$-dimensional
Seeley--DeWitt coefficients.  The earlier notation is recovered by
\(f_2=f_{2,g}\) and \(\Lcut=\Lambda_g\) only within this gravitational branch;
no equality with a matter or gauge process moment is implied.  On an optional
enlarged profile branch, the same symbol may be replaced explicitly by
\((4\pi)^{-2}\int_0^\infty F_g(u)\,du\); that replacement is not the
premise-minimal gravitational process.
\end{definition}


\begin{definition}[Unmarked minimal scalar Laplace-type resolution signature]
\label{p3:def:projected-grav-scalar-operator}
On \(\mathsf C_{\rm IR}^{\rm der}\), let
\(\Pi_0:\Cl(0,3)\to\mathbb R\mathbf1\) be the scalar infrared conditional
expectation on the physical protected block.  A gravitational resolution
generator is required to be a real, formally self-adjoint, metric-natural,
local scalar Laplace-type operator whose principal symbol is
\(a g^{\mu\nu}\) for one constant positive rate \(a>0\), with no auxiliary
line connection, potential, bundle, or channel insertion.  Before
normalization it has the form
\[
 L_{\rm grav}=a\Delta_+^{\rm scal}+E_{\rm grav},\qquad a>0.
\]
Its least positive resolution semigroup preserves the fixed unit:
\begin{equation}
\label{p3:eq:unit-preserving-grav-flow}
 e^{-tL_{\rm grav}}\mathbf1=\mathbf1
 \qquad(t\ge0).
\end{equation}
\end{definition}

\begin{lemma}[The minimal signature and unit preservation force the scalar Laplacian]
\label{p3:lem:scalar-remnant-laplacian}
The operator of Definition~\ref{p3:def:projected-grav-scalar-operator} satisfies
\begin{equation}
\label{p3:eq:scalar-remnant-E-zero}
 \boxed{L_{\rm grav}\simeq\Delta_+^{\rm scal},\qquad E_{\rm grav}=0,}
\end{equation}
where \(\simeq\) means equality after absorbing the positive constant \(a\)
into semigroup time, equivalently into \(\Lambda_g\).
\end{lemma}

\begin{proof}
On the real scalar line, formal self-adjoint Laplace type with no auxiliary
line connection has the unique Levi--Civita scalar connection, so there is no
independent first-order drift.  The requirement that the principal symbol be
the inherited cometric up to one constant rate excludes curvature-dependent
principal multipliers; allowing either enlargement would define a different
resolution signature.  Differentiate
\eqref{p3:eq:unit-preserving-grav-flow} at \(t=0\):
\[
 0=-L_{\rm grav}\mathbf1.
\]
Since \(\Delta_+^{\rm scal}\mathbf1=0\), one has
\(L_{\rm grav}\mathbf1=E_{\rm grav}\mathbf1\).  Therefore
\(E_{\rm grav}=0\) pointwise.  A term such as \(\xi\scal\), a mass, or any
other scalar potential would damp or weight the fixed unit and would no longer
represent the unital infrared conditional-expectation channel.  The spinorial
Lichnerowicz term belongs to the unprojected spinor representation and is not
the generator of this scalar unit-preserving channel.  This is also the
zero-mode specialization of Paper~II-D
Theorem~\ref{p2:thm:operator-rigidity}: invariant scalar second-order
principal symbols in this minimal scalar Laplace-type class are Laplacian up
to scale and a constant.  Positivity gives \(a>0\), unit preservation removes
the constant, and
\(t\mapsto at\) or \(\Lambda_g^2\mapsto\Lambda_g^2/a\) removes the remaining
rate convention.
\end{proof}

\begin{theorem}[Choice-free gravitational heat-role closure]
\label{p3e:thm:generated-gravitational-heat-role}
Let the scalar metric/protected root
\[
 d_g:=(g,A_{\rm prot},P_0,o_+)
\]
occur, where \(P_0\) is the derived terminal scalar projector and \(o_+\)
is the inherited positive resolution orientation.  In the unmarked minimal
real scalar Laplace-type domain of
Definition~\ref{p3:def:projected-grav-scalar-operator}, the gravitational
process role is a finitely presented generated-role descendant of \(d_g\).
Up to the positive calibration absorbed into \(\Lambda_g\), it is uniquely
\[
\boxed{
 F_g(u)=e^{-u}\quad(u\ge0),\qquad
 A_g(\Lambda_g)=e^{-\Delta_+^{\rm scal}/\Lambda_g^2}.}
\]
The historical labels \((\mathrm{GRAV}_{\circ})\) and
\((\mathrm{GRAV\!-\!HEAT})\) are therefore certifications of this theorem,
not independent occurrence assumptions.  The conclusion concerns the
generated process family and its law; it does not assert a separate
contingent token for every resolution time or fix the numerical calibration
\(\Lambda_g\).
\end{theorem}

\begin{proof}
The closure is the following eight-step construction.
\begin{enumerate}[label=\arabic*.,leftmargin=2.2em]
\item \emph{Protected resolution role.}
Paper~III-C gives
\(A_{\rm prot}=\mathbb R\mathbf1\oplus W\),
\(P_W=1-P_0\), and proves that the choice-free family
\[
 V_\tau=e^{-\tau P_W}=P_0+e^{-\tau}P_W
\]
is the unique unmarked unit-preserving common-rate semigroup after positive
time calibration.  It is co-realized from the occurring split rather than
introduced as a new process token.

\item \emph{Local scalar generator.}
The declared real, formally self-adjoint scalar Laplace-type signature has
\(L=a\Delta_+^{\rm scal}+E\), with one constant \(a>0\) and no auxiliary
line connection.  Lemma~\ref{p3:lem:scalar-remnant-laplacian} removes \(E\),
while \(a\) is absorbed into the scale.  Hence the spacetime generator is
\(\Delta_+^{\rm scal}\).

\item \emph{Global process law.}
The least continuous semigroup generated by this operator satisfies
\[
 H_{t+s}=H_tH_s,\qquad H_0=1,\qquad
 -\left.\partial_tH_t f\right|_{t=0+}=\Delta_+^{\rm scal}f
 \quad(f\in\operatorname{Dom}\Delta_+^{\rm scal}),
\]
and therefore \(H_t=e^{-t\Delta_+^{\rm scal}}\).  At
\(t=\Lambda_g^{-2}\), spectral calculus gives \(F_g(u)=e^{-u}\) on the
whole half-line, not merely at the eigenvalues of one finite carrier.
Nonlocal subordinations \(e^{-t\psi(\Delta)}\) or higher-order generators
leave the declared local second-order signature unless \(\psi(u)=au\).

\item \emph{Choice-free descent.}
The map
\[
 \mathcal F_g:d_g\longmapsto
 \bigl(e^{-t\Delta_g},e^{-\tau P_W}\bigr)
\]
is intrinsic and result-independent.  It is functorial under scalar-metric
isometries and protected orthogonal isomorphisms because functional calculus
commutes with conjugation.  It is faithful: differentiation at zero recovers
\(\Delta_g\) and \(P_W\), while
\(\lim_{\tau\to\infty}e^{-\tau P_W}=P_0\).  Its operation law, strong-operator topology, generator domain, and
target sorts have an exact presentation.  The derivative at zero is
strong on the generator domain, not an operator-norm derivative for an
unbounded scalar Laplacian (Paper~III-B,
Proposition~\ref{p3b:prop:v284-heat-generator-domain}).

\item \emph{G0 co-realization.}
Step 4 verifies Paper~II-A
Definition~\ref{p2:def:finitely-presented-generated-role-descendant}.
Theorem~\ref{p2:thm:generated-role-corealization} therefore gives
\[
 \PhysOcc(d_g)\wedge\mathcal F_g(d_g)\preccurlyeq_\Sigma d_g
 \Longrightarrow
 \operatorname{CoReal}(\mathcal F_g(d_g)\mid d_g).
\]
No second actualization principle or profile selector enters.

\item \emph{Closed readout.}
The gravitational role is unmarked and closed, so Paper~II-A's
loop--boundary exhaustion and Paper~III-B's trace theorem give
\[
 S_{\rm grav}(\Lambda_g)
 =\operatorname{Tr}e^{-\Delta_+^{\rm scal}/\Lambda_g^2}.
\]
On a compact smooth parent this operator is positive trace class for every
\(\Lambda_g>0\).  The reciprocal role average is formed only after the two
ordinary sector traces.

\item \emph{Stage-faithful weights.}
The occurring root is formed on the ambient protected octet before scalar
projection.  Host faithfulness therefore fixes
\[
 \tau_8(P_0)=\frac18.
\]
The Born orbit contributes
\(\tfrac12\mathcal I_{x,\ell}+\tfrac12\mathcal I_{p,\ell}\), and
Theorem~\ref{p3e:thm:born-orbit-eh-identification} pulls this to one full
copy for every natural pure-metric scalar density \(\ell\).  Thus the
reciprocal average contributes no second factor \(1/2\) on the Born diagonal;
the fixed vacuum and Einstein channels both retain the internal weight
\(1/8\).  The active ultraviolet state instead occupies the full octet and
has \(e^{S}/8=8/8=1\).

\item \emph{Canonical identification.}
The generated constituents have separating signatures:
fixed scalar support/zero derivatives, fixed scalar support/two derivatives,
and running \(1\oplus7\) support/nonlinear response.  Paper~II-A
Corollary~\ref{p2:cor:canonical-constituent-identification} uniquely routes
them to the vacuum, Einstein--Hilbert, and active nonlinear terms.  This
identifies the common exponential \emph{process-law role} on the internal and
metric carriers; it does not identify \(P_W\) with
\(\Delta_+^{\rm scal}\).  The active-state/Hessian law is supplied separately
by the minimal response classification, not by the heat-role theorem.
\end{enumerate}
\end{proof}

\begin{theorem}[Derived closed gravitational heat readout]
\label{p3e:thm:general-closed-gravitational-readout}
On the generated role of
Theorem~\ref{p3e:thm:generated-gravitational-heat-role}, \emph{and on a
finite-dimensional scalar carrier or a compact smooth scalar parent}, write
$\mathcal H_g$ for its declared scalar carrier.  Then the ordinary-trace
action form is derived rather than independently declared:
\[
 \boxed{
 A_g:=e^{-\Delta_+^{\rm scal}/\Lambda_g^2}
 \in\mathcal L^1(\mathcal H_g)_+,
 \qquad
 S_{\rm grav}=r_g
 :=\operatorname{Tr}_{\mathcal H_g}A_g.}
\]
Unitary naturality, positivity, extensive normalization, direct-sum
additivity on licensed blocks, and exclusion of ungenerated insertions follow
from the closed role.  The process law and moment are fixed; the numerical
scale calibration and the analytic domain of a chosen expansion are not.
For a noncompact parent, only a separately licensed finite
$\tau_{g,\mathrm{rel}}(A_g)$ can be claimed; it is a relative/localized/
volume-trace output, not the displayed ordinary $\mathcal L^1$ theorem.
\end{theorem}

\begin{proof}
Lemma~\ref{p3:lem:scalar-remnant-laplacian} and
Theorem~\ref{p3e:thm:generated-gravitational-heat-role} give
\(L_{\rm grav}=\Delta_+^{\rm scal}\) after scale calibration and
\(A_g=e^{-L_{\rm grav}/\Lambda_g^2}\).  It is positive trace class on the finite/compact branch by
Definition~\ref{p3e:def:grav-closed-occurrence} and standard finite/compact
heat-trace theory.  The defining generated role alone does not give ordinary
trace class on an unlocalized noncompact parent.  On the former branch
that definition freezes
the scalar diagram \(r_g=\operatorname{Tr}_{\mathcal H_g}A_g\), asserts
\(\operatorname{FSR}_{\Sigma}(r_g)\), and gives it unmarked closing
provenance.  Paper~II-A's loop--boundary exhaustion therefore puts \(r_g\)
in \(\Rloop\), and Paper~III-B
Theorem~\ref{p3:thm:gp2-dcl-derivation} identifies its closure with the
ordinary carrier trace.
The remaining properties are the trace and functional-calculus consequences
proved there.  No additional function or occurrence is selected.
\end{proof}

\begin{corollary}[Native heat action as a generated action-constructor descendant]
\label{p3e:cor:native-heat-action-constructor}
On the premise-minimal gravitational heat role \emph{with a finite or
compact ordinary-trace carrier}, the chain
\[
 d_g\longmapsto\Delta_+^{\rm scal}
 \longmapsto e^{-t\Delta_+^{\rm scal}}
 \longmapsto
 \operatorname{Tr}e^{-\Delta_+^{\rm scal}/\Lambda_g^2}
\]
is an explicit instance of Paper~III-A
Theorem~\ref{p3a:thm:generated-action-constructor-closure}.  Therefore an
independent replacement
$\operatorname{Tr}F(\Delta_+^{\rm scal}/\Lambda_g^2)$ is not another
realization of the same generated process unless the enlarged profile $F$ has
its own generated pre-action origin.  The conclusion is action selection on
this typed heat role; it is not a general rule that ranks arbitrary invariant
functionals by simplicity.
\end{corollary}

\begin{proof}
Theorem~\ref{p3e:thm:generated-gravitational-heat-role} generates the local
operator and the full continuous semigroup, while
Theorem~\ref{p3e:thm:general-closed-gravitational-readout} generates the ordinary closed trace on that branch.  A noncompact
relative/localized readout requires its own trace and closure certificate.
Hence the action-producing constructor is already
frozen before any alternate profile is evaluated, and Paper~III-A's
constructor-closure theorem applies.
\end{proof}

\begin{proposition}[Semiclassical scalar-channel Einstein--Hilbert coefficient]
\label{p3:prop:EH_from_a2}
Assume the fixed scalar-corner domain displayed above.  In the smooth-parent
semiclassical heat-kernel expansion of the generated heat process,
the \(a_2\) contribution of \(L_{\rm grav}\) to each reciprocal sector is
\begin{equation}
\Tr\!\left(e^{-L_{\rm grav}/\Lambda_g^2}\right)\supset f_{2,g}\,\Lambda_g^2\cdot
\frac{1}{6}\int \scal_s\,\sqrt{|g_s|}\,d^4x,
\qquad s\in\{x,p\},
\label{p3:eq:EH_sector_raw}
\end{equation}
where the universal factor \(1/6\) is the standard \(a_2\) coefficient for the minimally coupled scalar Laplacian in four dimensions \cite{berline1992heat,ChamseddineConnes1997}.  After the scalar infrared projection, Paper~III-B's derived ambient rank-one constituent, and the derived Born-diagonal pullback of Theorem~\ref{p3e:thm:born-orbit-eh-identification}, the linear curvature contribution to the effective gravitational action is
\begin{equation}
S_{\mathrm{EH}}
=\frac{f_{2,g}\Lambda_g^2}{48}\int \scal\,\sqrt{|g|}\,d^4x
=\frac{\Lambda_{\rm EH}^2}{8}\int \scal\,\sqrt{|g|}\,d^4x,
\label{p3:eq:EH_projected}
\end{equation}
where
\begin{equation}
\Lambda_{\rm EH}^2:=\frac{f_{2,g}}{6}\Lambda_g^2.
\label{p3:eq:Lambda_EH_def}
\end{equation}
The later Lorentzian section independently fixes the operational Lorentzian
sign from the GP1 cone; no Wick reconstruction enters the Euclidean
heat-kernel coefficient.
\end{proposition}

\begin{proof}
By Definition~\ref{p3:def:projected-grav-scalar-operator} and Lemma~\ref{p3:lem:scalar-remnant-laplacian}, the gravitational scalar channel is the projected minimal scalar Laplacian with \(E=0\).  The heat-kernel expansion for a scalar Laplace-type operator in four dimensions therefore gives
\[
a_2=\frac{1}{(4\pi)^2}\cdot\frac{1}{6}\int \scal\,\sqrt{|g|}\,d^4x .
\]
With the rescaled moment Definition~\ref{p3:def:f2_rescaled}, the
\((4\pi)^{-2}\) is absorbed into \(f_{2,g}\), yielding
\eqref{p3:eq:EH_sector_raw}.  The generated Born orbit contains both
half-weighted sector constituents.  Its diagonal pullback is, by
Theorem~\ref{p3e:thm:born-orbit-eh-identification},
\(\frac12+\frac12=1\) full scalar-curvature copy.  Off the Born diagonal the
two integrals remain distinct and the action remains their orbit average.
The protected scalar projection is derived on
\(\mathsf C_{\rm IR}^{\rm der}\).  Because the occurring
gravitational composite is formed on the ambient octet before that projection,
Paper~III-B Theorem~\ref{p3:thm:derived-scalar-corner-readout} derives its
physical constituent factor \(1/8\).  Therefore
\[
\frac{1}{8}\cdot\frac{1}{6}f_{2,g}\Lambda_g^2=\frac{f_{2,g}\Lambda_g^2}{48}
=\frac{\Lambda_{\rm EH}^2}{8},
\]
which gives \eqref{p3:eq:EH_projected}.
\end{proof}

\begin{theorem}[Fixed scalar and active-octet weights]
\label{p3e:thm:fixed-scalar-and-active-weights}
On the generated heat branch and
\((\mathrm{GRAV\!-\!ASYM})\), the Born-diagonal fixed scalar contribution has
the four-dimensional expansion
\begin{equation}
\label{p3e:eq:fixed-scalar-heat-expansion}
 \boxed{
 S_{\rm fixed}\sim
 \frac1{8(4\pi)^2}\int\sqrt{|g|}\,
 \left[\Lambda_g^4+\frac{\Lambda_g^2}{6}\scal
 +O(\scal^2,R_{\mu\nu}R^{\mu\nu},
 R_{\mu\nu\rho\sigma}R^{\mu\nu\rho\sigma},\nabla^2\scal)\right]d^4x.}
\end{equation}
Hence the protected internal support weights of the zero-derivative and
linear-curvature channels are
\[
 \boxed{w_{\rm vac}=w_{\rm EH}=\frac18.}
\]
For the co-realized active state of
Definition~\ref{p3:def:normalized-reduced-shell-length},
\[
 \boxed{w_{\rm nl}(U)=\frac18\,8^1=1,\qquad
        w_{\rm nl}(C)=\frac18\,8^0=\frac18.}
\]
The ultraviolet value is the full-octet weight of the nonlinear
inflation-candidate channel.  It does not by itself prove an inflationary
background, slow roll, exit, or reheating.  At the infrared endpoint the
MRESP Hessian visibility still vanishes,
\[
 C_{\rm sp}''(C)
 =w_{\rm nl}(C)\frac{\gamma_\sigma}{U}u(C)=0,
\]
so the terminal \(1/8\) is not a surviving nonlinear dark-energy term.
The zero-derivative coefficient in
\eqref{p3e:eq:fixed-scalar-heat-expansion} is a bare scalar vacuum channel.
Its renormalized sign and magnitude, and its identification with observed dark
energy, remain part of \((\mathrm{REN}_0)\) and the background/vacuum domain.
\end{theorem}

\begin{proof}
For \(F_g(u)=e^{-u}\), the scalar heat trace in four dimensions is
\[
 \Tr e^{-\Delta/\Lambda_g^2}
 \sim\frac1{(4\pi)^2}\int\sqrt{|g|}
 \left[\Lambda_g^4+\frac{\Lambda_g^2}{6}\scal+\cdots\right]d^4x.
\]
The process is formed on the ambient octet, so the generated scalar
constituent contributes \(\tau_8(P_0)=1/8\).  Theorem~\ref{p3e:thm:born-orbit-eh-identification}
applied to \(\ell=1\) and \(\ell=\scal\) turns the two half-weighted
reciprocal representatives into one full Born-diagonal copy; no further
factor \(1/2\) occurs.  This proves
\eqref{p3e:eq:fixed-scalar-heat-expansion} and the two fixed weights.

At \(U\), \(u=1\) and
\(\sigma_{\rm act}=I_8/8\), so \(S=\ln8\) and
\(e^S/8=1\).  At \(C\), \(u=0\) and
\(\sigma_{\rm act}=P_0\), so \(S=0\) and \(e^S/8=1/8\).
The minimal-response Hessian law contains the additional visibility \(u\);
hence its terminal value is zero.  Support counting fixes neither a
renormalized vacuum coefficient nor a cosmological solution, proving the
scope clauses.
\end{proof}

\begin{remark}[Which trace the factor \(1/8\) belongs to]
\label{p3:rem:corner-trace-convention}
Grade extraction alone does not fix this normalization, because two different
operators are available on \(\Cl(0,3)\) and they differ by exactly eight.  For
the regular representation of a scalar, \(\tau_8(L_{a_0\mathbf 1})=a_0\);
for the rank-one scalar projector, \(\tau_8(a_0P_0)=a_0/8\).  The coefficient
\(f_{2,g}/48\) is forced on the present branch because
Definition~\ref{p3e:def:grav-closed-occurrence} freezes the ambient protected
presentation before scalar extraction, and Paper~III-B's constituent theorem
then gives
\((\Tr_{\rm geom}\otimes\tau_8)(\,\cdot\otimes P_0)\).
The coefficient \(f_{2,g}/6\) belongs to a different branch in which the
one-dimensional scalar remnant is first isolated and then intrinsically
renormalized.  It is not a second realization of the same frozen
constituent presentation.  The same ambient presentation must be preserved
for the gauge channel of Paper~III-F.  Equality of the gravitational and
gauge process moments, however, requires that paper's separate
\((\mathrm{SAMEPROC}_{g/{\rm YM}})\) bridge; the common corner convention
alone does not identify their profiles.
\end{remark}

\begin{remark}[Optional enlarged-profile normalization]
If the premise-minimal local heat role is enlarged by an independently
occurring positive profile \(F_g\), its overall normalization is not fixed
on that optional asymptotic Einstein--Hilbert branch; rescaling
\(F_g\mapsto a_gF_g\) rescales
\(f_{2,g}\mapsto a_g f_{2,g}\).  The scalar-action compiler fixes only the product
\[
U=\Lambda_{\rm EH}^2=\frac{f_{2,g}}{6}\Lambda_g^2,
\]
not \(f_{2,g}\) and \(\Lambda_g\) separately.  If a physical Planck/Newton
legacy alias \(\Rp:=U\) is used, the convention \(\Lambda_g^2=U\) is
equivalent to \(f_{2,g}=6\), while \(f_{2,g}=1\) is equivalent to
\(\Lambda_g^2=6U\).  No observed Newton value selects a structural owner; it only calibrates the
presentation or the owner ratio.
\end{remark}


\begin{theorem}[Optional enlarged-profile normalization class and matched-EH coordinate]
\label{p3:thm:EH-normalization-gauge}
For the Einstein--Hilbert term on
\((\mathrm{OCC}_{g}^{\rm amb}[F_g])+(\mathrm{GRAV\!-\!ASYM})\) with an
independently supplied \(F_g\), two
presentations of the occurring process moment and resolution scale,
\[
(f_{2,g},\Lambda_g)\qquad\hbox{and}\qquad
(f_{2,g}',\Lambda_g'),
\]
are physically indistinguishable exactly when they have the same normalized
Einstein--Hilbert scale
\begin{equation}
\label{p3:eq:EH-normalization-equivalence}
\frac{f_{2,g}}{6}\Lambda_g^2
=\frac{f_{2,g}'}{6}(\Lambda_g')^2.
\end{equation}
Thus Newton's constant fixes the equivalence class
\begin{equation}
\label{p3:eq:EH-normalization-class}
\left[ f_{2,g},\Lambda_g\right]_{\rm EH}
:=\left\{(f_{2,g}',\Lambda_g'):
\frac{f_{2,g}'}6(\Lambda_g')^2=\Rp\right\},
\end{equation}
not the two entries separately.  In particular,
\begin{equation}
\label{p3:eq:f2-cutoff-gauge-examples}
\Lambda_g^2=\Rp\quad\Longleftrightarrow\quad f_{2,g}=6,
\qquad
f_{2,g}=1\quad\Longleftrightarrow\quad \Lambda_g^2=6\Rp,
\end{equation}
are two representatives of the same Einstein--Hilbert normalization class.
With the legacy alias \(\Rp:=U\), the same EH normalization class may be
written using \(\Rp\), but this is no longer a cross-owner condition.  Numerical
Newton comparison is metrological.  The unconditional class invariant remains
\((f_{2,g}/6)\Lambda_g^2=\Lambda_{\rm EH}^2=U\).
\end{theorem}

\begin{proof}
By Proposition~\ref{p3:prop:EH_from_a2}, the coefficient of the projected
Einstein--Hilbert term is
\[
\frac18\left(\frac{f_{2,g}}{6}\Lambda_g^2\right).
\]
Therefore two pairs \((f_{2,g},\Lambda_g)\) and
\((f_{2,g}',\Lambda_g')\) give the same Newton
coupling if and only if they obey \eqref{p3:eq:EH-normalization-equivalence}.
Using the legacy alias \(\Rp:=U\), the same common coefficient may be written
as \(\Rp\) by \eqref{p3e:eq:conditional-newton-matching}, giving
\eqref{p3:eq:EH-normalization-class} without a separate owner-selection step.
The examples in \eqref{p3:eq:f2-cutoff-gauge-examples} are immediate special
choices of representative.
\end{proof}


\textbf{Absolute-cutoff closure.}
The moment--scale equivalence class remains correct on the optional enlarged
profile branch.  On the premise-minimal branch,
Theorem~\ref{p3e:thm:generated-gravitational-heat-role} fixes
\(f_{2,g}=(4\pi)^{-2}\), and hence fixes
\(\Lambda_g^2=96\pi^2\Lambda_{\rm EH}^2\).  Conditional on the evaluated $(\mathrm{CELL}_{K\to U})$ port, with the momentum-curvature
comparison \(\mathcal K_p=6/\Rmom^2\), \(U=\Lambda_{\rm EH}^2=(3/\pi)\mathcal K_p\), this gives
\(\Lambda_g\Rmom=24\sqrt{3\pi}\).  The raw joint coefficient remains
\(\UEHjoint^{\rm raw}=96U\) and is not reinterpreted as a second Planck
curvature.


\begin{remark}[Separate origins of \(6\), \(2\pi\), and the Newton coordinate]
The scalar heat coefficient carries the local Seeley--DeWitt factor \(6\),
whereas the primitive Born phase cell carries \(\hBorn=2\pi\hbar\).  They are
distinct generated/represented structures.  For the joint-EH owner,
\(\UEHjoint=8\kappa_J\) and
\(G_{J,\rm raw}=c_{\rm op}^3/(\hBorn\UEHjoint)\) imply identically
\[
 \hbar\kappa_J=\frac{c_{\rm op}^3}{16\pi G_{J,\rm raw}}.
\]
Thus \(16\pi=8\times2\pi\) is bookkeeping of the joint-EH coefficient
coordinate and the Born phase-circle representative, not an extra fit factor.
The raw common-heat coefficient is instead transported downstream by
\(\UEHjoint^{\rm raw}=96\Lambda_{\rm EH}^2\).  The scalar coefficient $U$ and intrinsic curvature $\mathcal K_p$ have
a separate actual comparison port. The raw joint coordinate $G_{J,\rm raw}$
and $G_{\rm scal}$ differ by96; empirical comparison uses the finally
matched $G_{\rm phys}$.
\end{remark}

\begin{corollary}[Generated atomic gravitational heat branch]
\label{p3e:cor:grav-heat-branch}
On the generated role of
Theorem~\ref{p3e:thm:generated-gravitational-heat-role},
\((\mathrm{GRAV\!-\!HEAT})\) certifies
\(F_g(u)=e^{-u}\) for every \(u\ge0\), not merely an operator equality on one
spectrum at one resolution time.  Therefore
\[
\boxed{
\begin{gathered}
 S_{\rm grav}^{\rm heat}
 =\operatorname{Tr}e^{-\Delta_+^{\rm scal}/\Lambda_g^2},\\
 f_{2,g}^{\rm heat}=\frac1{(4\pi)^2},\qquad
 \Lambda_{\rm EH}^2=\frac{\Lambda_g^2}{6(4\pi)^2}.
\end{gathered}}
\]
Equivalently, before any classical action-role bridge,
\[
 \boxed{\Lambda_g^2=6(4\pi)^2\Lambda_{\rm EH}^2.}
\]
The later common-heat carrier theorem gives the raw represented transport
\(\UEHjoint^{\rm raw}=96\Lambda_{\rm EH}^2\).  This does not define a second
primitive curvature.  Conditional on the actual $(\mathrm{CELL}_{K\to U})$ comparison,
with $\mathcal K_p=6/\Rmom^2$ and $U=\Lambda_{\rm EH}^2=(3/\pi)\mathcal K_p$,
\[
 \boxed{\Lambda_g^2=96\pi^2U,\qquad \Lambda_g\Rmom=24\sqrt{3\pi}.}
\]
Thus the scalar moment--cutoff relation is fixed on the minimal heat
domain, while its momentum-curvature comparison remains conditional on the
evaluated cell port and the represented ordinary-trace
coefficient remains distinct.  No gauge or matter
process is identified with it without a separate same-process bridge.
\end{corollary}

\begin{proof}
Step 3 of
Theorem~\ref{p3e:thm:generated-gravitational-heat-role} uses spectral
calculus for the nonnegative self-adjoint smooth-parent generator and fixes
the process law on the whole half-line.  No finite-dimensional rigidity
theorem is needed.  Hence
\[
 f_{2,g}
 =\frac1{(4\pi)^2}\int_0^\infty e^{-u}\,du
 =\frac1{(4\pi)^2}.
\]
Substitution into
\(\Lambda_{\rm EH}^2=(f_{2,g}/6)\Lambda_g^2\) gives the second displayed
identity.  On the evaluated cell-port domain, combining it with
\(U=(3/\pi)\mathcal K_p\) and \(\mathcal K_p=6/\Rmom^2\) gives
\(\Lambda_g\Rmom=24\sqrt{3\pi}\).  The equality between the action process and heat
semigroup is supplied by the generated-role theorem, not by mathematical
heat-flow existence alone.
\end{proof}

\begin{corollary}[Singleton heat-profile fibre and descendant invariance]
\label{p3e:cor:singleton-heat-profile-fibre}
For every attained value \(\lambda=\Lambda_{\rm EH}^2>0\), the premise-minimal
gravitational heat role has the one-class dynamical profile fibre
\begin{equation}
 \label{p3e:eq:singleton-heat-profile-fibre}
 \boxed{
 \mathfrak P^{\rm heat}_\lambda
 =\left\{
  \left[
   F_g(u)=e^{-u},\quad
   f_{2,g}=\frac1{(4\pi)^2},\quad
   \Lambda_g^2=96\pi^2\lambda
  \right]
 \right\}.}
\end{equation}
Hence GP3 has no unresolved member on this role.  Every licensed descendant
that is generated from the same heat process satisfies Paper~III-B's profile
factorization theorem automatically.  On the Newton-matched branch one may
set \(\lambda=R_P\), but that empirical coefficient identification is not used
to select the profile.
\end{corollary}

\begin{proof}
Theorem~\ref{p3e:thm:generated-gravitational-heat-role} and
Corollary~\ref{p3e:cor:grav-heat-branch} fix the process law on the whole
half-line and give
\[
 F_g=e^{-u},\qquad f_{2,g}=(4\pi)^{-2}.
\]
The coefficient relation
\(\lambda=(f_{2,g}/6)\Lambda_g^2\) then has the unique positive solution
\(\Lambda_g^2=96\pi^2\lambda\).  The heat family has totally bounded realized
finite-window readout image, so it is GP1-admissible.  GP2 exactness retains
the full semigroup and excludes a different profile as an ungenerated process
root.  The exact moduli therefore contain one structured class, which GP3
retains.  Paper~III-B
Theorem~\ref{p3b:thm:gp-profile-invariance-factorization} now gives the
descendant statement.
\end{proof}

\begin{remark}[Optional cutoff profiles outside EH normalization]
Examples of sharp or later-model cutoff-profile representatives do not enter
the Einstein--Hilbert normalization theorem.  Their mathematical comparison is in
Paper~III-G, Appendix~\ref{p3e:app:profile-representatives}.
\end{remark}

\begin{remark}[Kinematic carrier density does not modify the projection derivation]
The Einstein--Hilbert coefficient \eqref{p3:eq:EH_projected} is fixed entirely by the heat-kernel coefficient and the Clifford scalar projection.  The Paper-II carrier density factor
\[
\nu_{\rm car}^{(s)}=(\lambda_{\rm car}^{(s)})^{3/2}
\]
may be used later to compare relative sector contributions after the reciprocal sector has been closed, but it does not alter the algebraic derivation of the factor $1/8$ in the scalar channel.
\end{remark}


\begin{remark}[Born--Liouville phase-volume identity and calibration condition]
The calibrated local cotangent phase-volume identity is supplementary to the
positive Einstein--Hilbert/causal/response chain and is treated in
Paper~III-G, Appendix~\ref{p3e:app:local-born-liouville}.  It requires
\((\mathrm{CAL}_{\hbar})\).  The canonical position and momentum
polarizations are Lagrangian and are not separately assigned symplectic
state counts.  No one-dimensional uncertainty-width statement is claimed
without a separate resolution-width/cell bridge.
\end{remark}

\subsection{Derivation of Lorentzian Signature}
\label{p3:sec:signature}

Within the derived \((\mathrm{CAUS})\) domain, the
round vector parent-role space of the derived rank-one branch and its current
GP2-exact/host-faithful realization is
\(S^3\cong SU(2)\) with a positive round metric.
Paper~I-A supplies the positive primitive Reeb lattice, Paper~II-A transports it
through the actual scalar-succession provenance occurrence, and Paper~III-A
supplies the oriented scalar elapsed cover, its reciprocal dual, and the
clocked-Hopf composite acting on the reciprocal round parent.  The least
constructor closure is the propagation doctrine itself, not an independently
occurring family.  Theorem~\ref{p3:thm:gp1-speed-calibration} therefore
supplies the sharp finite bound
\(d_h\le c_{\rm op}|\Delta\widetilde t|\).
The local cone and Lorentzian conformal sign follow from that bound before
numerical calibration or a field Hamiltonian is selected.  A local
second-order scalar field, when such
a sector is realized, must subsequently respect the already constituted
cone.  The inherited succession orientation is distinct from the conventional
simultaneous overall sign choice for a representative of the Lorentzian
conformal class.

\begin{theorem}[Lorentzian conformal class from the GP1 operation ceiling]
\label{p3:thm:lorentzian-signature}
Assume the complete \((\mathrm{CAUS})\) domain of
Definition~\ref{p3:def:causal-wick-response-hypotheses}.  On the local
universal-cover spacetime/role chart
\[
        \RR_{\widetilde t}\times\RR^3_x,
\]
let \(h\) be the positive isotropic spatial metric and orient
\(\widetilde t\) by the scalar succession order.  The closure of the
GP1-admissible future tangent operations is
\begin{equation}
\label{p3:eq:gp1-operational-causal-cone}
\mathcal C^+
:=
\left\{(\tau,\xi)\in\mathbb R\oplus\mathbb R^3:
\tau\ge0,\ h(\xi,\xi)\le c_{\rm op}^2\tau^2\right\}.
\end{equation}
Its sharp boundary is the null cone of the conformal metric
\[
 g_{\rm op}=-c_{\rm op}^2d\widetilde t^{\,2}+h.
\]
Hence the local operational conformal class has signature
\[
        \boxed{(-,+,+,+)}
\]
up to the conventional simultaneous overall reversal.  Among
\(O(3)\)-invariant nondegenerate quadratic forms, this conformal class is
uniquely determined by the sharp GP1 cone.  This statement neither creates a
field sector nor identifies the cone with an empirical light cone.  Replacing
\((h,\widetilde t,c_{\rm op})\) by a jointly homothetic representative leaves
\(\mathcal C^+\) and \([g_{\rm op}]\) unchanged; \(({\rm CAL}_{xp})\) only
assigns numerical units.
\end{theorem}

\begin{proof}
For every admissible propagation history,
Theorem~\ref{p3:thm:gp1-speed-calibration} gives
\(d_h(x(t),x(s))\le c_{\rm op}|t-s|\).  Paper~III-A,
Corollary~\ref{p3:cor:gp1-round-causal-envelope}, then uses homogeneity,
actual \(O(3)\) tangent isotropy, and the stationary operation to identify the
closed convex hull of the physical tangent operations with the entire
radius-\(c_{\rm op}\) ball.  The seed and its isometry conjugates attain the
boundary sphere; convexification fills the mathematical interior without
asserting memberwise physical occurrence.
Homogeneous rescaling by the nonnegative
elapsed tangent \(\tau\) is exactly
\eqref{p3:eq:gp1-operational-causal-cone}.

Apply Paper~I-D, Proposition~\ref{p1b:prop:local-coframe}~\cite{Kugelmann2026Paper1D},
to this already derived sharp cone, with
\((t,c,h)=(\widetilde t,c_{\rm op},h)\).  It uniquely determines the
isotropic conformal quadratic form up to overall sign.  Choosing positive
spatial sign gives \(-c_{\rm op}^2d\widetilde t^{\,2}+h\), hence the
displayed Lorentzian signature.
\end{proof}

\begin{theorem}[Derived clock--Hopf null-ray geometry]
\label{p3e:thm:clock-hopf-null-ray-geometry}
Assume the complete \((\mathrm{CAUS})\) domain, the clock calibration
\((\mathrm{CAL}_t)\), and the actual provenance-oriented clock--Hopf word of
Paper~III-A, Definition~\ref{p3a:def:clocked-hopf-propagation-doctrine}.
Let \(\xi_{0,x}\) be the \(2\pi\)-periodic infinitesimal generator of its
position-side Hopf action.  Transport the Paper-II-D principal connection as
the oriented pair \((\eta_{0,x},\xi_{0,x})\), normalized by
\[
 \eta_{0,x}(\xi_{0,x})=1,
 \qquad d\eta_{0,x}=\pi_x^*\omega_x .
\]
Here the orientation is defined by the actual contragredient action itself;
this convention absorbs its inverse-action sign and is not an untyped naive
pushforward from the momentum parent.  For a round representative of radius
\(R_x>0\), put
\[
 \eta_x:=R_x\eta_{0,x},
 \qquad e_H:=R_x^{-1}\xi_{0,x},
 \qquad \eta_x(e_H)=1,
 \qquad h(e_H,e_H)=1 .
\]
On the elapsed cover
\(\widetilde{\mathcal M}_{\rm op}=\RR_{\widetilde t}\times S_x^3\), let
\(u_t=\partial_{\widetilde t}\), so one quotient clock turn has duration
\(T_t\), and recall
\(c_{\rm op}=2\pi R_x/T_t\).  Then the retained clock, Hopf operation, and
vector-bundle operations canonically generate
\begin{equation}
\label{p3e:eq:clock-hopf-null-fields}
 \boxed{
 N_\pm:=\frac1{\sqrt2}
 \left(c_{\rm op}^{-1}u_t\pm e_H\right).}
\end{equation}
For the representative
\(g_{\rm op}=-c_{\rm op}^2d\widetilde t^{\,2}+h\),
\[
 g_{\rm op}(N_\pm,N_\pm)=0,
 \qquad g_{\rm op}(N_+,N_-)=-1.
\]
Both fields are complete.  Their path families
\begin{equation}
\label{p3e:eq:clock-hopf-null-path-family}
 \gamma_{\pm,z}(s):=\operatorname{Fl}^{N_\pm}_s(z),
 \qquad z\in\widetilde{\mathcal M}_{\rm op},
\end{equation}
are therefore generated null-ray families.  On the compact clock quotient
the diagonal and antidiagonal operations close after the common parameter
period \(\sqrt2\,c_{\rm op}T_t=\sqrt2\,2\pi R_x\); on the elapsed cover they
are nonclosed null helices.

The pair and the full basepoint-indexed path family are structural outputs of
the same rooted clock--Hopf component.  A particular basepoint, a hadron/beam
alignment, or an actually traversed ray is token-level data and is not
selected by this theorem.
\end{theorem}

\begin{proof}
Paper~II-D, Corollary~\ref{p2d:cor:generated-hopf-normalization-graph},
retains the principal/metric Hopf normalization graph, while Paper~III-A
retains the provenance-fixed position-side action and its primitive period.
Defining the transported connection by the actual oriented action keeps the
generator and connection in one native operation diagram, including the
contragredient sign.  The canonical tangent splitting of the product and its
vector-space operations then generate \eqref{p3e:eq:clock-hopf-null-fields}.
Product orthogonality gives
\[
 g_{\rm op}(c_{\rm op}^{-1}u_t,c_{\rm op}^{-1}u_t)=-1,
 \qquad g_{\rm op}(e_H,e_H)=1,
\]
which proves both metric identities.  Completeness follows from completeness
of the clock translation and compact Hopf action.  During the displayed
common parameter interval the clock component makes one quotient turn and
the spatial component makes one Hopf turn, proving the closure statement.
Paper~II-A,
Theorem~\ref{p2a:thm:universal-corealization-bridge-normal-form}, retains the
complete constructor incidence; it does not select a member of the path
family.
\end{proof}

\begin{corollary}[Null coframe, local ray chart, and no global light front]
\label{p3e:cor:null-coframe-frobenius-boundary}
The vertical clock--Hopf coframe dual to \(N_\pm\) is
\begin{equation}
\label{p3e:eq:clock-hopf-null-coframe}
 \boxed{
 \ell^\pm:=\frac1{\sqrt2}
 \left(c_{\rm op}d\widetilde t\pm\eta_x\right),
 \qquad d\ell^\pm=\pm\frac1{\sqrt2}d\eta_x .}
\end{equation}
It satisfies \(\ell^\alpha(N_\beta)=\delta^\alpha_\beta\) on the vertical
two-plane, but
\begin{equation}
\label{p3e:eq:clock-hopf-frobenius-obstruction}
 \ell^\pm\wedge d\ell^\pm
 =\frac12\left(
 \pm c_{\rm op}d\widetilde t\wedge d\eta_x
 +\eta_x\wedge d\eta_x\right)\ne0 .
\end{equation}
Thus the global null operations do not generate a null-hyperplane foliation
or global flat coordinates \(x^\pm\).  On a two-dimensional clock--Hopf orbit
cylinder the curvature pulls back to zero; with physical clock coordinate
\(t\) and Hopf angle \(\theta_H\),
\[
 x^\pm=\frac1{\sqrt2}
 \left(c_{\rm op}t\pm R_x\theta_H\right)
\]
obeys \(dx^\pm=\ell^\pm\) after restriction.  The theorem therefore supplies
null-ray geometry, not an unproved global light-front quantization surface.
\end{corollary}

\begin{proof}
The clock form is closed and the Hopf curvature is horizontal, which gives
\eqref{p3e:eq:clock-hopf-null-coframe}.  The two summands in
\eqref{p3e:eq:clock-hopf-frobenius-obstruction} have respectively a clock leg
and no clock leg, while \(d\eta_x\) is nonzero and
\(\eta_x\wedge d\eta_x\) is the Hopf contact volume.  Their sum cannot
vanish.  Frobenius therefore forbids the claimed hypersurface foliation.
Restriction to the clock--fiber cylinder removes the horizontal curvature,
and direct differentiation gives the local two-dimensional coordinates.
\end{proof}

\begin{theorem}[Positive local null-generator components]
\label{p3e:thm:positive-null-generator-components}
Let \((E,p)\) lie in the local positive polar of Paper~III-A,
\[
 E\ge c_{\rm op}\|p\|_{h^{-1}},
\]
and put \(p_H:=p(e_H)\).  The derived local components
\begin{equation}
\label{p3e:eq:positive-null-generator-components}
 \boxed{
 p^\pm:=\frac1{\sqrt2}
 \left(E/c_{\rm op}\pm p_H\right)}
\end{equation}
satisfy
\[
 p^\pm\ge0,
 \qquad p^++p^-=\sqrt2\,E/c_{\rm op},
 \qquad p^+-p^-=\sqrt2\,p_H .
\]
They are cotangent generator-germ components under
Theorem~\ref{p3a:thm:canonical-generator-germ-localization}; they are neither
global affine coordinates on the compact reciprocal carrier nor finite-QCD
operators.  If a separately proved local mass-shell relation holds, its
usual null-component identity is a conditional algebraic corollary, not part
of this theorem.
\end{theorem}

\begin{proof}
Dual Cauchy--Schwarz and the positive-polar inequality give
\[
 |p_H|\le\|p\|_{h^{-1}}\le E/c_{\rm op},
\]
which proves positivity; addition and subtraction give the remaining
identities.  The scope statement is exactly the local/global boundary in
Paper~III-A, Remark~\ref{p3:rem:polar-dual-not-global-pair}.
\end{proof}

\begin{definition}[Derived null-ray package and alignment boundary]
\label{p3e:def:lf-ray-geometry-package}
On the preceding theorem domain define
\[
 \boxed{
 \LFRayGeom
 :=\left(N_\pm,\ell^\pm,
          \{\gamma_{\pm,z}\}_{z\in\widetilde{\mathcal M}_{\rm op}},
          p^\pm\right).}
\]
This package is derived.  It contains no selected basepoint or hadron state,
no global null foliation, no finite-QCD representation, no color connection
or holonomy, and no continuum or renormalization limit.  A hadronic readout
may call one member ``plus'' only after either a declared state/beam alignment
maps its momentum direction to the retained Hopf ray or exact covariance
proves independence of the conjugate ray.  That alignment is not supplied by
the simultaneous occurrence of the geometry and the state.
\end{definition}

\begin{corollary}[Chiral projective realization of local null directions]
\label{p3:cor:chiral-projective-null-directions}
Fix a spatial orientation and an \(h\)-orthonormal frame at \(x\).  The
protected representation already co-realizes
\[
 \mathcal S_{\rm prot}=\mathcal S_{{\rm prot},L}
 \oplus\mathcal S_{{\rm prot},R}.
\]
For either already co-realized constituent \(\chi=L,R\), write
\(\mathcal S_{{\rm prot},\chi}\cong\mathbb C^2\) and, for
\(0\ne z\in\mathcal S_{{\rm prot},\chi}\), write
\[
        zz^\dagger
        = k^0\mathbf1+\sum_{i=1}^3 k^i\sigma_i .
\]
Then
\[
        [z]\longmapsto
        \left[\left(\frac{k^0}{c_{\rm op}},\,\sum_{i=1}^3k^i e_i\right)\right]
\]
is a bijection
\[
 \mathbb P(\mathcal S_{{\rm prot},\chi})\cong\mathbb{CP}^1
 \xrightarrow{\ \cong\ }
 \mathbb P\!\left(\partial\mathcal C_x^+\setminus\{0\}\right).
\]
Under this identification use the branchwise complexified action
\[
 \rho_L(A)z:=Az,
 \qquad
 \rho_R(A)z:=(A^\dagger)^{-1}z,
 \qquad A\in SL(2,\mathbb C).
\]
For either \(\chi=L,R\), the induced projective action
\([z]\mapsto[\rho_\chi(A)z]\) descends through \(\{\pm I\}\) to the
corresponding standard proper orthochronous Lorentz action on future null rays,
\[
        PSL(2,\mathbb C)\cong SO^+(1,3),
\]
while its compact subgroup
\(PSU(2)\cong SO(3)\) is the spatial rotation action.  This is a local
null-direction and Lorentz-action bridge.  It does not constitute a
finite-dimensional unitary boost dynamics, select a preferred physical
chirality, or identify the projective rays with quantum states.  The left and
right branches are branchwise realizations on constituents that already
co-realize upstream.
\end{corollary}

\begin{proof}
Paper~II-E,
Theorem~\ref{p2e:thm:celestial-hermitian-ray-bridge}, proves the purely
algebraic bijection
\[
        [z]\longmapsto[zz^\dagger]
\]
between \(\mathbb{CP}^1\) and the positive rank-one rays in
\(\operatorname{Herm}_2(\mathbb C)\) on the left constituent, and the
analogous conjugate-dual construction on the right constituent, together with
\[
        \det(zz^\dagger)
        =(k^0)^2-\sum_i(k^i)^2=0,
        \qquad k^0>0,
\]
and branchwise equivariance: \(K\mapsto AKA^\dagger\) on the left and
\(K\mapsto(A^\dagger)^{-1}KA^{-1}\) on the right.
Theorem~\ref{p3:thm:lorentzian-signature} independently identifies the
physical future null boundary, in the chosen orthonormal frame, by
\[
        h(\xi,\xi)=c_{\rm op}^2\tau^2,\qquad \tau>0.
\]
Setting \(k^0=c_{\rm op}\tau\) and
\(\xi=\sum_i k^i e_i\) composes the two
bijections.  The standard kernel \(\{\pm I\}\) gives the projective
isomorphisms stated above.  Since only \(PSU(2)\) preserves the chosen
positive Hermitian norm, the noncompact directions are not thereby promoted
to finite-dimensional unitary operations.
\end{proof}

\begin{corollary}[Protected projective ray as primal and dual null direction]
\label{p3e:cor:projective-primal-dual-null-ray}
On the domain of Corollary~\ref{p3:cor:chiral-projective-null-directions} and
Theorem~\ref{p3:thm:gp1-operation-generator-polar-dual}, one has a canonical
projective identification
\[
 \boxed{\mathbb P(\mathcal S_{{\rm prot},\chi})\cong\mathbb{CP}^1
 \cong \mathbb P(\partial\mathcal C_x^+\setminus\{0\})
 \cong \mathbb P(\partial(\mathcal C_x^+)^\vee\setminus\{0\}).}
\]
For a unit spatial direction \(n\), the normalized primal boundary ray may be
represented by \(v_n=(1,c_{\rm op}n)\), while the supporting dual boundary
ray is represented by \(q_n=(c_{\rm op},n^\flat)\); then
\[
 \langle q_n,v_n\rangle=0.
\]
Thus the protected spinor ray fixes the common null direction on both the
operation and generator sides, but not its energy/frequency scale.
\end{corollary}

\begin{corollary}[Local passive scalar principal part on the derived cone]
\label{p3:cor:causal-scalar-principal-part}
Suppose a theorem-domain sector realizes one real, passive, unlabelled,
local second-order scalar field whose associated tangent field-front boundary
exhausts the sharp boundary of \(\mathcal C^+\), and whose Hamiltonian
principal part is positive.  Let \(c_{\rm field}\) denote its front speed.
The exhaustion hypothesis gives
\[
        c_{\rm field}=c_{\rm op}.
\]
After canonical normalization of the scalar pair
\((\varphi,\pi)\), its principal Hamiltonian and Lagrangian densities are
\begin{align}
 H_2[\varphi,\pi]
 &=
 \frac12\int\!\left(\pi^2+
 c_{\rm field}^2h^{ij}\partial_i\varphi\,\partial_j\varphi\right)d\mu_h,
 \label{p3:eq:minimal-positive-scalar-hamiltonian}\\
 \mathcal L_2
 &=
 \frac12\left((\partial_{\widetilde t}\varphi)^2-
 c_{\rm field}^2h^{ij}\partial_i\varphi\,\partial_j\varphi\right).
 \label{p3:eq:minimal-lorentzian-scalar-lagrangian}
\end{align}
The equality \(c_{\rm field}=c_{\rm op}\) is a consequence of the displayed
exhaustion hypothesis, not of GP1 alone.  A medium or interacting effective
sector whose characteristics lie strictly inside \(\mathcal C^+\) may have
\(c_{\rm field}<c_{\rm op}\).
\end{corollary}

\begin{proof}
Before canonical normalization write the positive principal Hamiltonian as
\[
H_{\rm pr}
=\frac12\int\!\left(
a\pi^2+b\,h^{ij}\partial_i\varphi\,\partial_j\varphi
\right)d\mu_h,
\qquad a,b>0.
\]
Spatial isotropy leaves only the displayed gradient term, and the absence of
a preferred spatial vector excludes a mixed principal term.  With
\(A\ge0\) the nonnegative scalar Laplacian, Hamilton's equations give
\(\partial_{\widetilde t}^2\varphi=-abA\varphi\), so
\(c_{\rm field}^2=ab\), not a ratio of the two coefficients.
Equality of the tangent field-front boundary with
Theorem~\ref{p3:thm:lorentzian-signature} therefore fixes
\(c_{\rm field}=c_{\rm op}\).
The canonical change
\(\varphi'=\varphi/\sqrt a\), \(\pi'=\sqrt a\,\pi\) preserves the symplectic
pairing and gives
\[
H_{\rm pr}
=\frac12\int\!\left(
(\pi')^2+c_{\rm field}^2h^{ij}\partial_i\varphi'\,\partial_j\varphi'
\right)d\mu_h.
\]
Its Legendre transform is the displayed Lagrangian.
\end{proof}

\begin{theorem}[Minimal local scalar cone-faithful principal order]
\label{p3:thm:minimal-local-scalar-cone-order}
Assume an already constituted scalar propagation role on the smooth stationary
parent is local differential of finite order, has real scalar principal
symbol, is spatially isotropic, is hyperbolic with a nonempty spatially
propagating characteristic set, and is cone-faithful in the typed sense
defined below.  Write a tangent operation as
\(v=(\tau,\xi)\in T_xM\cong\mathbb R\oplus T_xX\) and a principal covector
as \(\zeta=(\omega,k)\in T_x^*M\cong\mathbb R\oplus T_x^*X\).

No scalar realization of order zero or one can realize the Lorentzian
characteristic boundary.  The least possible local scalar principal order is
therefore two.  At that order spatial isotropy forces, after an overall
positive multiplication, the cotangent principal symbol
\[
  p_x(\zeta)=\sigma_2(P)(\omega,k)
  =a\omega^2-b\,h^{-1}(k,k),
  \qquad a,b>0.
\]
For the positive characteristic sheet
\(\operatorname{Char}_x^+(P)=\{p_x=0,\ \omega>0\}\), define the tangent
field-front cone by the closed convex hull of the base projections of its
positive Hamilton rays,
\begin{equation}
\label{p3e:eq:typed-field-front-cone}
 \mathcal F_x^+(P)
 :=\overline{\operatorname{cone}\operatorname{conv}}
 \left\{\pi_*H_{p_x}(\zeta):
 \zeta\in\operatorname{Char}_x^+(P)\right\}.
\end{equation}
Then, writing \(c_{\rm field}^2=b/a\),
\[
 \boxed{\mathcal F_x^+(P)
 =\{(\tau,\xi):\tau\ge0,\
       \|\xi\|_h\le c_{\rm field}\tau\}.}
\]
Cone faithfulness means the now type-correct inclusion
\(\mathcal F_x^+(P)\subseteq\mathcal C_x^+\), and it is equivalent to
\(0<c_{\rm field}\le c_{\rm op}\).  Equivalently, the polar inclusions on
the cotangent generator side reverse:
\[
 (\mathcal C_x^+)^\vee\subseteq(\mathcal F_x^+(P))^\vee,
 \qquad
 (\mathcal F_x^+(P))^\vee
 =\{(E,p):E\ge c_{\rm field}\|p\|_{h^{-1}}\}.
\]
If, as an additional condition, the field-front boundary equals the sharp
generated tangent boundary, then
\[
  c_{\rm field}=c_{\rm op},
  \qquad
  \sigma_2(P)(\omega,k)
  \propto \omega^2-c_{\rm op}^2h^{-1}(k,k).
\]
The theorem fixes only the least principal differential structure.  It neither
constructs the propagation role from the causal cone nor removes admissible
lower-order scalar data.
\end{theorem}

\begin{proof}
Order zero has no nontrivial characteristic propagation.  For order one,
spatial isotropy forces an invariant scalar principal symbol
\(a\omega+b^ik_i\) to have \(b^i=0\).  Its Hamilton base projection has no
spatial component, so it cannot realize a Lorentzian front boundary.  Hence
the least spatially propagating order is
two.  At degree two the only spatially isotropic scalar quadratic invariants are
\(\omega^2\) and \(h^{-1}(k,k)\); the mixed scalar-vector term is
excluded by the absence of a preferred spatial covector.  Comparing the
signs with the assumed nonempty real hyperbolic characteristic set gives,
after an overall sign choice, \(a,b>0\).

In canonical coordinates the base projection of the Hamilton vector is
\[
 \pi_*H_{p_x}(\omega,k)
 =(2a\omega,-2b\,k^\sharp).
\]
On \(p_x=0\) with \(\omega>0\),
\(\omega=\sqrt{b/a}\,\|k\|_{h^{-1}}\).  Therefore every nonzero Hamilton
ray satisfies
\[
 \frac{\|\xi\|_h}{\tau}
 =\frac{b\|k\|_{h^{-1}}}{a\omega}
 =\sqrt{\frac ba}=c_{\rm field}.
\]
Spatial isotropy supplies every boundary direction, and the closed convex
conic hull fills exactly the stated field-front cone.  Comparing two tangent
cones now gives
\(\mathcal F_x^+(P)\subseteq\mathcal C_x^+\) exactly when
\(c_{\rm field}\le c_{\rm op}\).  The polar formula follows directly from
the pairing \(E\tau-p(\xi)\), and polarity reverses inclusions.  Equality of
the two boundaries forces equality of their radii and hence of the speeds.
\end{proof}


\begin{theorem}[Primitive unlabelled scalar support ancestry]
\label{p3e:thm:primitive-scalar-support-ancestry}
Let an already constituted local scalar propagation role be real, nonzero,
primitive, neutral, and unlabelled: its fiber support is generated by the
dagger-fixed scalar tensor-unit constructor and carries no independent
species, multiplicity, charge, grading, variance, or carrier owner.  Then its
support is the same structured neutral support as the host tensor-unit line,
\[
 \boxed{\operatorname{Supp}(\varphi)=I_{\mathbb R}=\mathbb R,}
\]
and a second neutral scalar carrier is not an admissible realization of this
branch.  This theorem fixes support ancestry only; it does not fix a kinetic
operator, principal symbol, or speed.
\end{theorem}

\begin{proof}
Paper~II-A,
Corollary~\ref{p2a:cor:unique-neutral-tensor-unit-support}, identifies
\(I_{\mathbb R}\) as the unique nonzero irreducible object in the unlabelled
real neutral support sort.  The scalar fiber constructor supplies its
canonical nonzero host-equivariant tensor-unit map.  Paper~II-A,
Theorem~\ref{p2a:thm:unique-neutral-support-ancestry}, therefore makes that
map surjective onto the host neutral support.  Any second copy requires a
species or multiplicity owner and is outside the stated branch.  The
same-support speed countermodel of Paper~II-A shows that no principal-speed
conclusion follows at this step.
\end{proof}


\begin{theorem}[Metric-principal descendant saturation]
\label{p3e:thm:metric-principal-descendant-saturation}
Assume the hypotheses of
Theorem~\ref{p3:thm:minimal-local-scalar-cone-order} and the primitive
unlabelled branch of
Theorem~\ref{p3e:thm:primitive-scalar-support-ancestry}.  Suppose in addition
that the scalar role's principal quadratic relation has already been generated
as a same-owner descendant of the operational Lorentzian quadratic relation,
with no generated constitutive tensor, independent clock, anisotropic
principal tensor, or other deformation of that principal owner.  Work on the
smooth globally hyperbolic product representative of the local operational
domain, with smooth lower-order coefficients.  Then, up to
 a positive overall scalar multiplier,
\[
 \boxed{
 \sigma_2(P)(\omega,k)
 \propto
 \omega^2-c_{\rm op}^2h^{-1}(k,k),}
\]
so its characteristic boundary is exactly the sharp operational null boundary
and
\[
 \boxed{c_{\rm field}=c_{\rm op}.}
\]
Thus boundary exhaustion is a theorem on the metric-principal descendant
branch rather than an additional field-speed premise.  This theorem does not
construct the field role itself and does not apply to a medium/effective
sector carrying a separately generated constitutive deformation.
\end{theorem}

\begin{proof}
The support theorem discharges support ancestry but, by itself, says nothing
about the cone.  The inverse of the operational conformal metric fixes the
cotangent parent principal relation
\(\omega^2-c_{\rm op}^2h^{-1}(k,k)\) up to an overall positive conformal
factor.  The additional hypothesis places the child relation in that same
cotangent principal owner with no deformation.  Paper~II-A,
Theorem~\ref{p2a:thm:undeformed-host-principal-inheritance}, therefore gives
an invertible overall owner factor multiplying the parent quadratic form.
For the real scalar line this factor is a positive scalar.  It does not change
the characteristic zero set, so the two null boundaries coincide.  The
resulting scalar operator is normally hyperbolic on the stated globally
hyperbolic product domain; the standard Cauchy/finite-propagation theorem
therefore identifies its wave-front boundary with that null cone
\cite[Theorems~3.2.11--3.2.12]{BarGinouxPfaffle2007}.  This discharges the analytic hypothesis in
Paper~II-A's inheritance theorem and gives equality of the positive front
speeds.
\end{proof}

\begin{definition}[Local scalar free-law domain with derived minimal principal order]
\label{p3:def:paper-iv-finite-field-input}
Let \(A=\Delta_{S^3,h}\ge0\) denote the normalized nonnegative round scalar
Laplacian of
Paper~II, Theorem~\ref{p2:thm:operator-rigidity}, with
\[
\sigma_2(A)(x,\xi)=h^{ij}\xi_i\xi_j,
\qquad A\mathbf1=0.
\]
All locality and differential-order statements below concern this smooth
parent operator before restriction to a finite spectral subspace.  The
\emph{local scalar free-law domain}
 consists of real conservative passive sectors with time-translation invariance,
spatial isotropy, an already constituted local finite-order cone-faithful
scalar propagation role, and a scalar invariant zeroth-order term. The
smooth-parent free operator is required to be formally self-adjoint on the
invariant product measure
$d\widetilde t\,d\operatorname{vol}_h$, on compactly supported time tests
and the closed spatial manifold. In particular this is a no-time-drift
domain. Passivity, stationarity and isotropy alone would also allow a
dissipative term $\gamma\partial_{\widetilde t}$ and do not imply the
free dispersion relation below. By
Theorem~\ref{p3:thm:minimal-local-scalar-cone-order}, its least nondegenerate
principal realization is second order.  Its
characteristic speed is denoted \(c_{\rm field}\), so containment means
\[
        0<c_{\rm field}\le c_{\rm op}.
\]
Its \emph{primitive unimpeded subdomain} is the primitive unlabelled support
branch of Theorem~\ref{p3e:thm:primitive-scalar-support-ancestry} whose
principal quadratic relation is a same-owner descendant of the operational
metric relation and carries no generated principal deformation.  Theorem~\ref{p3e:thm:metric-principal-descendant-saturation}
then derives
\(c_{\rm field}=c_{\rm op}\) and discharges the boundary-exhaustion
hypothesis of Corollary~\ref{p3:cor:causal-scalar-principal-part}.  On the least primitive realization, locality together with
Theorem~\ref{p3:thm:minimal-local-scalar-cone-order} excludes \(A^j\) for
\(j\ge2\) and nonpolynomial spectral profiles from the principal law; such
operators may belong to enlarged effective roles, and finite functional
calculus alone would not exclude them.  The rest-stable subdomain requires the
zeroth-order coefficient itself to be nonnegative; this does not follow
merely from positivity after compression to a block that omits the zero mode.
On the primitive subdomain it is parametrized by a sector rest value
\(\mu\ge0\).
\end{definition}

\begin{theorem}[Papers II--III finite massive dispersion chain]
\label{p3:thm:p3-p4-lorentzian-chain-closure}
On every finite spectral block in the primitive unimpeded subdomain of
Definition~\ref{p3:def:paper-iv-finite-field-input}, a separated mode
\[
\psi(\widetilde t,x)=e^{-i\omega\widetilde t}\phi_\lambda(x),
\qquad
A\phi_\lambda=\lambda\phi_\lambda,
\]
satisfies
\[
        \omega^2=c_{\rm field}^2\lambda+
        \frac{\mu^2c_{\rm field}^4}{\hbar^2}.
\]
Equivalently, with \(\widetilde E=\hbar\omega\) and
\(\|p\|_h^2=\hbar^2\lambda\),
\begin{equation}
\label{p3:eq:derived-lorentzian-mass-shell}
\widetilde E^2=c_{\rm field}^2\|p\|_h^2+
\mu^2c_{\rm field}^4.
\end{equation}
Here \(c_{\rm field}=c_{\rm op}\) in the primitive unimpeded subdomain.
Paper~IV may realize this relation as a positive-frequency one-particle
Hamiltonian on a stable finite block and interpret its spectral bottom as a
threshold.  It is not an input to the Lorentzian-sign theorem.
\end{theorem}

\begin{proof}
Corollary~\ref{p3:cor:causal-scalar-principal-part} fixes the principal wave
operator to
\(\partial_{\widetilde t}^2+c_{\rm field}^2A\), with
\(c_{\rm field}=c_{\rm op}\) in the stated subdomain. Spatial isotropy
excludes an invariant spatial first-order vector term. Time-translation
invariance makes a possible real time-drift coefficient constant, but
$\gamma\partial_{\widetilde t}$ has formal adjoint
$-\gamma\partial_{\widetilde t}$ on the declared invariant product
measure. Formal self-adjointness therefore forces $\gamma=0$. This uses
the conservative smooth-parent domain before any finite compression;
passivity alone would not remove damping. The remaining lower-order free
term is the invariant scalar zeroth-order coefficient, written
\(\mu^2c_{\rm field}^4/\hbar^2\).  Hence
\[
\left(\partial_{\widetilde t}^2+c_{\rm field}^2A+
\frac{\mu^2c_{\rm field}^4}{\hbar^2}\right)\psi=0.
\]
Substitution of the separated mode gives the frequency relation; multiplying
by \(\hbar^2\) gives \eqref{p3:eq:derived-lorentzian-mass-shell}.
\end{proof}

\begin{corollary}[Dilation constraint on the \(({\rm CAUS})\) branch]
\label{p3:cor:closed-lorentzian-dilation-constraint}
On the generated Born-diagonal pure-metric subbranch of the derived
\(({\rm CAUS})\) Lorentzian branch, the effective-density and dilation
identities of Theorem~\ref{p3:thm:dilation_identity} apply.  Causality supplies
the Lorentzian signature only; the measure identity is supplied independently
by Born-diagonal metric naturality.  No Paper-IV input is needed for the cone,
sign, or determinant identity.
\end{corollary}

\begin{proof}
Theorem~\ref{p3:thm:lorentzian-signature} supplies the signature on
\(({\rm CAUS})\).  Corollary~\ref{p3e:cor:born-diagonal-effective-density-balance}
supplies the determinant-density identity after restriction along
\(j_{\rm B}\).  Theorem~\ref{p3:thm:dilation_identity} then gives the local
dilation relation.
The local free-sector dispersion theorem supplies additional field
interpretation but does not enter the determinant identity.
\end{proof}

\begin{corollary}[Gravitational Euclidean heat readout and separate claim-scoped reconstruction]\label{p3:cor:wick-from-gp-constitution}
On the generated gravitational heat role, the Euclidean spectral functional
is the scalar-Laplacian heat readout
\[
 S_{\rm grav}(\Lambda_g)
 =\Tr_{\mathcal H_g}\!\left(e^{-\Delta_+^{\rm scal}/\Lambda_g^2}\right).
\]
This heat trace is a closed spectral readout; Paper~III-E proves no
identification of it with a Euclidean quantum two-point covariance or higher
correlation hierarchy.  If such Euclidean quantum correlations are
independently supplied or constructed and a Lorentzian reconstruction claim
is made, cyclic two-point continuation is marked by
\(({\rm WICK}_{2{\rm cyc}})\), identification with the same Lorentzian field
variables by \(({\rm WICK}_{\rm same})\), finite OS reconstruction by
\(({\rm WICK}_{\rm OS,fin})\), and interacting infinite-volume reconstruction
by \(({\rm WICK}_{\rm int,\infty})\).  Within Paper~III-E only conclusions
whose required reconstruction conditions are available are obtained.  The
free Lorentzian equation of
Theorem~\ref{p3:thm:p3-p4-lorentzian-chain-closure} was derived directly from
the operational cone and does not use this corollary.
\end{corollary}
\begin{proof}
The displayed heat trace is well defined by the generated scalar elliptic
operator, spectral calculus, and the compact smooth-parent heat theory already
used above.  A scalar spectral trace is not, merely by that construction, a
specified quantum two-point kernel, and no same-process map from the
Paper~III gravitational heat role to a quantum correlation hierarchy is
proved here.  Consequently analytic continuation of a separately supplied
Euclidean correlation object to a unitary Lorentzian theory is a distinct
claim.  The claim-scoped tags record the analyticity, reflection-positivity,
compatibility, contour, and limiting data not proved in this paper.  Without
such an additional correlation/reconstruction domain, the present result
asserts only the generated Euclidean gravitational heat readout.
\end{proof}

\begin{remark}[Paper~IV closure of finite Euclidean-to-Hamiltonian certificates]
\label{p3:rem:p4-euclidean-hamiltonian-closure}
Corollary~\ref{p3:cor:wick-from-gp-constitution} separates the generated
gravitational heat trace from any separately specified Euclidean quantum
correlation hierarchy and records the claim-scoped \(({\rm WICK}_\bullet)\)
status inside Paper~III-E.  Paper~IV later proves finite cyclic reconstruction
for positive-spectral kernels, a same-carrier criterion, and on an explicitly
Hamiltonian-generated ordinary/even branch a forward Euclidean semigroup and
finite reflection-positivity certificate.  Those later theorems may discharge
the corresponding finite tags on their stated domains; they neither identify
\(S_{\rm grav}\) with a quantum covariance nor prove interacting
infinite-volume reconstruction, and they cannot be imported as unstated
Paper~III premises.
\end{remark}


\begin{remark}[Operational Lorentzian cone and claim-specific Wick scope]
Theorem~\ref{p3:thm:lorentzian-signature} and
Corollary~\ref{p3:cor:wick-from-gp-constitution} state the precise additional
status of the dynamics-side Lorentzian construction imported from Paper~II:
\(({\rm CAUS})\) is now the derived GP1 operational-cone tag, while each
relevant \(({\rm WICK}_\bullet)\) tag marks a reconstruction conclusion not
proved inside Paper~III-E.  Paper~IV later converts the finite tags into
checkable or automatic certificates on explicitly stated positive-kernel or
Hamiltonian-generated domains; the interacting infinite-volume tag remains a
later-layer condition.  The finite massive mode relation is proved in
Theorem~\ref{p3:thm:p3-p4-lorentzian-chain-closure}; Paper~IV supplies its
finite quantum/event realization rather than the Lorentzian sign.
\end{remark}

\begin{proposition}[From kinematic spectral gaps to finite-volume energy thresholds]
\label{p3:prop:spectral-gap-mass-upgrade}
This proposition records the Paper~III upgrade of the Paper~II compact
spectral gaps.  The gap data there were kinematic; in the minimal local
second-order free scalar domain they become finite-volume energy thresholds
after composing the scalar completion, the GP1 Lorentzian cone, and
Theorem~\ref{p3:thm:p3-p4-lorentzian-chain-closure}.  Paper~IV supplies the
positive-frequency one-particle and Born-event interpretation of those
thresholds.

Let \(\lambda\ge0\) denote an eigenvalue of the same normalized scalar
Laplacian \(A\) fixed in
Definition~\ref{p3:def:paper-iv-finite-field-input}.  On the scalar-cover
Lorentzian branch, without invoking any \(({\rm WICK}_\bullet)\) tag, a separated mode
\(\psi(\ttilde,x)=e^{-i\omega\ttilde}\phi_\lambda(x)\) satisfies
\[
\omega^2=c_{\rm field}^2\lambda
+\mu^2 c_{\rm field}^4/\hbar^2,
\]
where \(\mu\) is the sector's rest parameter and
\(c_{\rm field}=c_{\rm op}\) by the primitive-unimpeded hypothesis.
Consequently, a nonzero
kinematic spectral gap $\lambda_{\min}>0$ imported from Paper~II gives the
finite free-sector energy threshold
\[
E^2=\hbar^2 c_{\rm field}^2\lambda_{\min}
+\mu^2 c_{\rm field}^4.
\]
\end{proposition}

\begin{proof}
Before the scalar completion of Theorem~\ref{p3:thm:scalars}, the eigenvalues
of the compact spatial operators are purely kinematic numbers.  After the
GP1 cone/sign theorem is composed with the minimal local second-order free
domain of Theorem~\ref{p3:thm:p3-p4-lorentzian-chain-closure}, the
universal-cover time variable \(\ttilde\) enters with Lorentzian sign, so the
squared wave operator on a separated mode is
\[
-\partial_{\ttilde}^2\psi
=\bigl(c_{\rm field}^2\lambda
+\mu^2 c_{\rm field}^4/\hbar^2\bigr)\psi.
\]
Writing \(\psi(\ttilde,x)=e^{-i\omega\ttilde}\phi_\lambda(x)\) gives the
stated dispersion relation.  Thus the gaps acquire an energy-threshold
interpretation in the displayed local free domain.  Calling such a threshold
a particle rest mass additionally uses the stable one-particle/event
interpretation supplied by Paper~IV.
\end{proof}

\begin{remark}[Compact phases vs. causal evolution]
The scalar \emph{phase} is compact (\(\Sone\)), while the \(({\rm CAUS})\)
real-time variable is placed on its universal cover.  This separation avoids
interpreting the compact phase circle itself as a timelike spacetime loop; a
global causality theorem would require the full Lorentzian geometry and is not
claimed here.
\end{remark}


\subsection{The Dilation Constraint}
\label{p3:sec:dilation-constraint}

Paper~II records the momentum-primary sector density law
\(\nu_{\rm car}^{(s)}=(\lambda_{\rm car}^{(s)})^{3/2}\) on the compact
carrier.  Paper~III then introduces effective position- and momentum-side
scale factors in the reciprocal chart.  This subsection records the local
effective-density consequence on the generated Born-diagonal action measure.
The reciprocal product itself is derived by the premise-minimal Paper-II-C
intrinsic-evaluation/GP2 closure and completed in Paper~III-A.  On the
generated metric image, sectorwise measure balance is a conclusion of metric
volume-density naturality.  Exchange invariance of the normalized compact
Haar product is derived separately below and is not used to assert any
relation between independently parameterized effective densities.

\begin{theorem}[Normalized compact Haar Born balance]
\label{p3e:thm:compact-haar-born-balance}
Let \(A\cong SU(2)\cong S^3(\Rmom)\) be the constituted compact round role,
let \(B_A\) be its generated reciprocal compact role, and let
\(\beta_A:A\to B_A\) be the generated group isomorphism.  For normalized
Haar measures \(\mu_A\) and \(\mu_B\),
\[
 (\beta_A)_*\mu_A=\mu_B,
 \qquad
 \beta_A^*\mu_B=\mu_A.
\]
The Born exchange
\[
 \mathcal B_{\beta_A}(a,b)
 :=\bigl(\beta_A^{-1}(b),\beta_A(a)\bigr)
\]
therefore preserves the normalized product measure:
\[
 \boxed{
 (\mathcal B_{\beta_A})_*(\mu_A\otimes\mu_B)
 =\mu_A\otimes\mu_B.}
\]
Moreover the normalized constant
\(e_0=\operatorname{Vol}(S^3(\Rmom))^{-1/2}\), whose geometric scalar-ground
and density properties are exported by Paper~I-A, induces the same Haar
probability:
\[
 |e_0|^2d\operatorname{vol}
 =\frac{d\operatorname{vol}}
        {\operatorname{Vol}(S^3(\Rmom))}
 =\mu_A.
\]
No physical vacuum preparation is used in this measure identity.
\end{theorem}

\begin{proof}
Paper~II-A, Theorem~\ref{p2a:thm:generated-born-inheritance}, says that a
licensed product-measure descendant of the complete unmarked reciprocal pair
must transform equivariantly; the present Haar argument identifies that target
action concretely and shows that it is trivial.  The pushforward
\((\beta_A)_*\mu_A\) is a normalized left- and right-invariant measure on the
compact group \(B_A\).  Uniqueness of normalized Haar measure
therefore gives \((\beta_A)_*\mu_A=\mu_B\), and bijectivity gives the pullback
identity.  Under \(\mathcal B_{\beta_A}\), the first factor is transported
from \(B_A\) to \(A\) and the second from \(A\) to \(B_A\); the two preceding
identities prove product invariance.  The final equality is Paper~I-B,
Corollary~\ref{p1:cor:normalized-constant-haar-density}; importing it introduces
no preparation or effective-metric premise.
\end{proof}

\begin{corollary}[Haar/effective-density firewall]
\label{p3e:cor:haar-effective-balance-firewall}
Theorem~\ref{p3e:thm:compact-haar-born-balance} derives compact Haar transport
and product exchange invariance.  It does not imply equality of independently
parameterized local effective metric-volume densities.  That stronger
sectorwise statement follows only after a theorem identifies those densities
with the transported Haar densities.
\end{corollary}

\begin{proof}
Product invariance constrains \(\mu_A\otimes\mu_B\) as one measure on the
reciprocal product.  It does not equate arbitrary factor densities
\(d\nu_x^{\rm eff}\) and \(d\nu_p^{\rm eff}\).  Independent positive local
determinant factors can change the latter while leaving the normalized Haar
theorem untouched, proving the strict logical separation.
\end{proof}

\begin{corollary}[Effective-density balance on the generated Born diagonal]
\label{p3e:cor:born-diagonal-effective-density-balance}
Let \(\beta:M_x\to M_p\) be the constituted reciprocal diffeomorphism.  On
the pure-metric image
\[
             j_{\rm B}(g)=(g,\beta_*g)
\]
the effective volume densities obey
\[
             \boxed{\beta^*d\mu_{\beta_*g}=d\mu_g.}
\]
Thus \((\mathrm{BAL}_B^{\rm eff})\) is \texttt{DERIVED} on the generated
Born-diagonal metric branch.  The present theory makes no extension claim for
independently parameterized effective sector metrics.
\end{corollary}

\begin{proof}
Metric volume density is natural under diffeomorphism:
\(d\mu_{\beta_*g}=\beta_*d\mu_g\).  Pullback by \(\beta\) gives the
identity.  The scope is strict: the conformal family in
Paper~III-G, Appendix~\ref{p3e:app:remaining-owner-counterfamilies} changes the effective
density ratio while preserving the normalized Haar theorem, so an
off-diagonal extension would be a different conditional theorem.
\end{proof}

\begin{remark}[Scope of the Born-diagonal effective-density identity]
\label{p3e:def:born-measure-balance-interface}
On the generated metric image, the historical notation
\((\mathrm{BAL}_B^{\rm eff})\), or \((\mathrm{BAL}_B)\) when unambiguous,
denotes the theorem-level identity
\[
 \boxed{(\mathrm{BAL}_B^{\rm eff}):\qquad
 \mathcal B^*d\nu_p^{\rm eff}=d\nu_x^{\rm eff}.}
\]
It is the conclusion of
Corollary~\ref{p3e:cor:born-diagonal-effective-density-balance}, not an
antecedent.  For independently parameterized metrics the same equation would
select a proper balanced subfamily, but no active result in this paper or its
downstream consumers requires that extension.
\end{remark}

\begin{theorem}[Born-diagonal effective-density and dilation identity]
\label{p3:thm:dilation_identity}
\hfill \newline
Use the \(({\rm CAUS})\) Lorentzian-signature branch of
Theorem~\ref{p3:thm:lorentzian-signature}, equivalently work with absolute
determinant densities of scalar-plus-round-vector sector forms.  On a local
reciprocal chart adapted to the constituted diffeomorphism
\(\mathcal B:M_x\to M_p\), restrict to the generated Born diagonal
\(g_p=\mathcal B_*g_x\).  \(O(3)\)-isotropy of each vector sector and the
absence of an inserted preferred vector force the two effective
two-derivative sector forms to have the diagonal scalar-plus-round-vector form
\[
ds_x^2=-\lambda_t\,dt^2+\lambda_x\,d\vec{x}^{\,2},
\qquad
ds_p^2=-\lambda_E\,dE^2+\lambda_p\,d\vec{p}^{\,2},
\]
with \(\lambda_t,\lambda_x,\lambda_E,\lambda_p>0\) and \(\vec{x},\vec{p}\in\RR^3\).  Let
\[
 d\nu_x^{\rm eff}:=\sqrt{|g_x|}\,d^3x\,dt,
 \qquad
 d\nu_p^{\rm eff}:=\sqrt{|g_p|}\,d^3p\,dE
\]
be the resulting sector measures.  Define the absolute local Born-chart
Jacobian \(J_{\mathcal B}>0\) relative to the displayed coordinate volume
forms by
\begin{equation}
 \mathcal B^*(d^3p\,dE)
 =J_{\mathcal B}\,d^3x\,dt.
 \label{p3:eq:born-chart-jacobian}
\end{equation}
Born-diagonal metric naturality gives
\begin{equation}
\label{p3:eq:sectorwise-born-balanced-measure}
 \mathcal B^*d\nu_p^{\rm eff}=d\nu_x^{\rm eff}.
\end{equation}
Then the scale parameters satisfy the Jacobian-covariant local relation
\begin{equation}
\boxed{\ \lambda_x^{3}\lambda_t
       =J_{\mathcal B}^{\,2}\lambda_p^{3}\lambda_E.\ }
\label{p3:eq:dilation_cubic}
\end{equation}
The previously displayed equality without a Jacobian is precisely the
unit-Jacobian special case \(J_{\mathcal B}=1\); no such normalization is
assumed.  No relation is asserted for an independently parameterized pair
\((g_x,g_p)\).
\end{theorem}

\begin{proof}
First derive the displayed metric form.  In a reciprocal chart the vector part of each sector carries the isotropy group \(O(3)\), so Schur/rotation invariance permits only a scalar multiple of the round Euclidean quadratic form on the three vector coordinates.  The scalar Hamiltonian/Born completion supplies one scalar coordinate in each reciprocal sector.  A mixed term such as \(dt\,dx^i\) or \(dE\,dp^i\) would require a preferred vector in the corresponding \(O(3)\)-module, and a mixed \(dt\,dE\) or \(dx^i dp^j\) term would be an additional reciprocal-sector intertwiner not supplied by the completed datum.  Thus the effective sector metrics have exactly the diagonal forms stated above.

For
\(g_x=\mathrm{diag}(-\lambda_t,\lambda_x,\lambda_x,\lambda_x)\),
\[
 \sqrt{|g_x|}=\lambda_x^{3/2}\lambda_t^{1/2},
\]
and for
\(g_p=\mathrm{diag}(-\lambda_E,\lambda_p,\lambda_p,\lambda_p)\),
\[
 \sqrt{|g_p|}=\lambda_p^{3/2}\lambda_E^{1/2}.
\]
Using \eqref{p3:eq:born-chart-jacobian}, the pulled-back momentum-sector
measure is
\[
 \mathcal B^*d\nu_p^{\rm eff}
 =J_{\mathcal B}\lambda_p^{3/2}\lambda_E^{1/2}
   d^3x\,dt.
\]
Corollary~\ref{p3e:cor:born-diagonal-effective-density-balance} gives
\eqref{p3:eq:sectorwise-born-balanced-measure}.  It therefore follows,
pointwise at
paired chart points,
\[
 \lambda_x^{3/2}\lambda_t^{1/2}
 =J_{\mathcal B}\lambda_p^{3/2}\lambda_E^{1/2}.
\]
Squaring gives \eqref{p3:eq:dilation_cubic}.  The normalized Haar-product
identity of Theorem~\ref{p3e:thm:compact-haar-born-balance} does not identify
independently parameterized effective marginals and is not used to extend the
theorem beyond the generated diagonal.
\end{proof}

\begin{remark}[Product invariance versus sectorwise balance]
\label{p3:rem:product-vs-sectorwise-balance}
The effective local phase-space product
\[
 d\mu_{\rm ph}^{\rm eff}=d\nu_x^{\rm eff}\,d\nu_p^{\rm eff}
\]
is not identified with the compact Haar product merely because both live on
the reciprocal carrier.  Even a separately imposed product statement would
need its marginal normalization and transported-factor typing before it could
be converted into sector equality.  The dilation constraint instead uses
the generated identity \(g_p=\mathcal B_*g_x\); volume-density naturality then
identifies the two effective sector measures and retains the chart Jacobian
explicitly.  Logical use of the constraint comes
after the \(\Dcl\) classification and sectorwise-additivity result, even
though the elementary measure computation is recorded here for continuity of
the phase-space discussion.
\end{remark}

\begin{remark}[Local versus global momentum-measure implementation]
The Born-diagonal dilation theorem is local and does not depend on a global momentum
localization convention.  The detailed compact-
\(S^3\)/POVM implementation note is deferred to
Paper~III-G, Appendix~\ref{p3e:app:global-momentum-measure}.
\end{remark}

\begin{corollary}[Linear form in volume-normalized variables]
\hfill \newline
\label{p3:cor:dilation_linear}
Define
\[
\mu_x:=\lambda_x^{3},\qquad \mu_p:=\lambda_p^{3}.
\]
Then \eqref{p3:eq:dilation_cubic} is equivalent to the linear constraint
\[
\boxed{\ \mu_x\,\lambda_t
       =J_{\mathcal B}^{\,2}\mu_p\,\lambda_E.\ }
\]
\end{corollary}

\begin{proof}
Immediate from \eqref{p3:eq:dilation_cubic}.
\end{proof}


\subsection{Restoration of Conformal Symmetry}
\label{p3:subsec:conformal-restoration}

The microscopic scalar sector is an \(\Sone\) phase, while the standard smooth
conformal compactification uses a mathematical conformal-time circle
\(S^1_{\rm conf}\).  No identification of these two circles is made without an
additional bridge.  The comparison below concerns the standard conformal-time
compactification and must also be kept separate from the finite fuzzy cutoff
representation.

\begin{theorem}[Smooth conformal compactification and finite-cutoff no-go]
\label{p3:thm:conformal-compactification-cutoff}
The manifold \(S^3\times S^1_{\rm conf}\), up to the usual \(\ZZ_2\) identification, is the standard smooth conformal compactification of Minkowski space, and its universal cover is the Einstein static universe \(S^3\times\RR\).  The connected conformal group \(SO_0(2,4)\) acts conformally on this smooth/untruncated compactification.  However, a finite fuzzy carrier cannot support a nontrivial exact finite-dimensional unitary representation of the noncompact conformal group.  At finite cutoff this rules out an exact nontrivial finite-dimensional unitary realization of the connected noncompact conformal group on the fixed fuzzy Hilbert carrier. It does not rule out nonunitary finite-dimensional algebraic conformal representations, nor the projected cutoff covariance. The conformal action is therefore only the projected/cutoff action
\begin{equation}
\label{p3:eq:projected-conformal-action}
        U_m(g):=P_mU(g)P_m,
\end{equation}
which is exact only on transformations preserving the retained bandlimited subspace and otherwise gives cutoff covariance rather than a genuine group representation.
\end{theorem}

\begin{proof}
The first statement is the standard conformal compactification: replacing noncompact Minkowski time by the mathematical conformal-time circle gives \(S^3\times S^1_{\rm conf}\) modulo the antipodal \(\ZZ_2\) identification, and passing to the universal cover gives \(S^3\times\RR\).  The smooth conformal group is \(SO_0(2,4)\).

For the finite-cutoff statement, suppose a finite fuzzy carrier carried a nontrivial exact unitary representation of \(SO_0(2,4)\).  Differentiating gives a Lie-algebra homomorphism
\[
        \mathfrak{so}(2,4)\longrightarrow \mathfrak u(N).
\]
The source is noncompact simple, while every Lie subalgebra of \(\mathfrak u(N)\) is compact as a real Lie algebra with negative-semidefinite Killing form on the semisimple part.  The kernel of a homomorphism from a simple Lie algebra is either zero or the whole algebra; an injective homomorphism would embed a noncompact simple algebra into a compact one, impossible.  Hence the finite-dimensional unitary representation is trivial on the noncompact conformal directions.

Therefore the exact \(SO_0(2,4)\) action belongs to the smooth/untruncated compactification.  On a fuzzy sector one may project a smooth representation by \(P_m\), giving \eqref{p3:eq:projected-conformal-action}; in general \(P_mU(g_1g_2)P_m\ne P_mU(g_1)P_mU(g_2)P_m\) unless the retained bandlimited subspace is preserved.  This is the precise finite-cutoff content used in Paper~III.
\end{proof}


\section{Co-realized hierarchy and conditional scalar response}
\label{paper3e:sec:organized-scalar}
\subsection{Co-realized hierarchy-active state and protected action weights}
\label{p3:sec:index-locking}

Paper~III-D's secant exponent \(n_{\rm sec}\) is a geometric response
coordinate.  Its entropy-matched state is an auxiliary comparison and is not
used to define the active shell.  The occurring
\((\mathrm{MRESP}_{\min})\) composite
instead contains Paper~III-C's complete visibility family at the same
hierarchy parameter that defines \(\scal(\tau)\).  Constituent co-realization
therefore fixes the active state directly.

\begin{definition}[Hierarchy-active state and entropy index]
\label{p3:def:normalized-reduced-shell-length}
On \(C\le\scal\le U\), put
\begin{equation}
\label{p3:eq:active-log-visibility}
 u(\scal)
 :=\frac{\ln(\scal/C)}{\ln(U/C)}\in[0,1],
 \qquad
 \tau_{\rm H}(\scal):=-\ln u(\scal),
\end{equation}
with \(\tau_{\rm H}(C)=+\infty\).  Define
\begin{equation}
\label{p3:eq:hierarchy-active-state}
 \sigma_{\rm act}(\scal)
 :=\rho_{\tau_{\rm H}(\scal)}
 =\frac{P_0+u(\scal)P_W}{1+7u(\scal)}.
\end{equation}
Its entropy, normalized active index, effective count, and active fraction are
\[
\begin{aligned}
 S_{\rm act}(\scal)&:=S(\sigma_{\rm act}(\scal)),
 &n_{\rm act}(\scal)&:=\frac{S_{\rm act}(\scal)}{\ln8},\\
 N_{\rm eff}(\scal)&:=e^{S_{\rm act}(\scal)},
 &w_{\rm nl}(\scal)&:=\frac{N_{\rm eff}(\scal)}8.
\end{aligned}
\]
This is the active protected constituent of the hierarchy response, not a
redefinition of \(n_{\rm sec}\).
\end{definition}

\begin{theorem}[Uniqueness of the entropy-effective count]
\label{p3e:thm:entropy-effective-count-uniqueness}
Let \(\mathcal D_{\rm fin}\) be the class of finite-dimensional density
matrices, closed under tensor products.  Suppose
\(N:\mathcal D_{\rm fin}\to\mathbb R_{>0}\) is entropy-sufficient and
continuous, so \(N(\rho)=F(S(\rho))\) for a continuous
\(F:[0,\infty)\to\mathbb R_{>0}\), and suppose
\[
 N(\rho\otimes\eta)=N(\rho)N(\eta),
 \qquad N(I_8/8)=8.
\]
Then necessarily
\[
 \boxed{N(\rho)=e^{S(\rho)}}
\]
for every \(\rho\in\mathcal D_{\rm fin}\).  In particular \(N=1\) on pure
states and the uniquely octet-normalized active fraction is
\(w_{\rm nl}(\rho)=e^{S(\rho)}/8\).  Thus no functional
\(\mathcal K\)-selector survives in
\((\mathrm{MRESP}_{\min})\); abandoning entropy sufficiency or tensor
multiplicativity defines an enlarged response role.
\end{theorem}

\begin{proof}
Von Neumann entropy is additive under tensor products, so the hypotheses give
\(F(s+t)=F(s)F(t)\) whenever \(s,t\) occur as finite-state entropies.  Every
number \(s\ge0\) occurs: choose \(n\) with \(s/n\le\ln2\); the binary entropy
of \(\operatorname{diag}(p,1-p)\) continuously fills \([0,\ln2]\), and the
\(n\)-fold tensor power has entropy \(s\).  Hence the functional equation
holds on all of \(\mathbb R_{\ge0}\).

Positivity permits \(G:=\log F\).  Then \(G\) is continuous and additive on
\(\mathbb R_{\ge0}\).  For rational \(q\ge0\), additivity gives
\(G(q)=qG(1)\); continuity extends this to
\(G(s)=\kappa s\) for every \(s\ge0\).  The octet normalization yields
\(8=F(\ln8)=8^\kappa\), hence \(\kappa=1\).  Therefore \(F(s)=e^s\).
At \(s=0\) this also proves the pure-state normalization rather than requiring
it separately.
\end{proof}

\begin{theorem}[Derived active-state law and exact endpoint weights]
\label{p3:thm:index-locking}
Assume the occurring \((\mathrm{MRESP}_{\min})\) composite, including
\(\mathrm{REG}_{\sigma}\in\mathfrak A_{\rm REG}^{+}\) and the CRP4 response
summand.  Then
\(\sigma_{\rm act}\) of
Definition~\ref{p3:def:normalized-reduced-shell-length} is co-realized without
a separate MON occurrence tag, and
\begin{equation}
\label{p3:eq:entropy-index-locking-closed}
 \boxed{
 \begin{aligned}
 S_{\rm act}(\scal)
 &=\ln(1+7u)
   -\frac{7u}{1+7u}\ln u,\\
 n_{\rm act}(\scal)
 &=\frac{\displaystyle
      \ln(1+7u)-\frac{7u}{1+7u}\ln u}{\ln8},\\
 N_{\rm eff}(\scal)&=8^{n_{\rm act}(\scal)},\\
 w_{\rm nl}(\scal)&=\frac18\,8^{n_{\rm act}(\scal)}.
 \end{aligned}}
\end{equation}
Here \(u=u(\scal)\) and \(u\ln u:=0\) at \(u=0\).  The functions are
continuous and strictly increasing on the open interval, with
\[
 \boxed{
 n_{\rm act}(C)=0,\quad w_{\rm nl}(C)=\frac18,
 \qquad
 n_{\rm act}(U)=1,\quad w_{\rm nl}(U)=1.}
\]
In general \(n_{\rm act}\ne n_{\rm sec}\); in particular
\[
 n_{\rm act}(U)=1,
 \qquad
 n_{\rm sec}(U)=\frac{U-C}{U+C}<1.
\]
\end{theorem}

\begin{proof}
Paper~III-C's owner-typed inheritance corollary first gives
\(C=U\alpha_{\rm hier}(\sigma_{\rm fin})^{-7}\).  Together with
\(\scal(\tau)=U\alpha_{\rm hier}(\sigma_{\rm fin})^{-7(1-e^{-\tau})}\),
this yields
\[
 \ln(\scal/C)=7e^{-\tau}\ln\alpha_{\rm hier}(\sigma_{\rm fin}),
 \qquad
 \ln(U/C)=7\ln\alpha_{\rm hier}(\sigma_{\rm fin}).
\]
Hence \(e^{-\tau}=u(\scal)\), proving
\eqref{p3:eq:active-log-visibility}.  The complete visibility family is a
frozen generated-constituent presentation of the occurring MRESP response.
Paper~III-C
Corollary~\ref{p3c:cor:corealized-protected-visibility-family}, itself an
application of the Paper-II-A co-realization theorem, therefore realizes
\(\rho_{\tau_{\rm H}(\scal)}\) and its entropy without another occurrence
certificate.

The spectrum of \(\sigma_{\rm act}\) is
\[
 \frac1{1+7u}\quad\text{once},\qquad
 \frac{u}{1+7u}\quad\text{seven times}.
\]
Direct substitution gives the displayed entropy.  Moreover,
\[
 \frac{dS_{\rm act}}{du}
 =-\frac{7\ln u}{(1+7u)^2}>0
 \qquad(0<u<1).
\]
Since \(u(\scal)\) is strictly increasing, so are \(S_{\rm act}\),
\(n_{\rm act}\), \(N_{\rm eff}\), and \(w_{\rm nl}\).  At \(u=0\) the state
is \(P_0\) and the entropy is zero; at \(u=1\) it is \(I_8/8\) and the entropy
is \(\ln8\).  Exponentiating
\(S_{\rm act}=n_{\rm act}\ln8\) proves the remaining identities and endpoint
values.  The last comparison uses Paper~III-D's continuous secant value.
\end{proof}

\begin{theorem}[Canonical identification of the two protected action weights]
\label{p3:rem:index-locking-status}
On the occurring protected gravitational composite, the exhausted generated
constituents with signatures
\[
\begin{array}{c|c|c|c}
&\text{support}&\text{scale behavior}&\text{action role}\\ \hline
P_0&1&\text{fixed}&\text{linear EH scalar}\\
\sigma_{\rm act}&1\oplus7&\text{running}&\text{nonlinear protected response}
\end{array}
\]
are simultaneously realized and have a unique type-preserving
identification with the two corresponding action terms.  The ambient octet
trace therefore gives
\[
 \boxed{w_{\rm EH}=\frac18\quad\text{at every scale},\qquad
 w_{\rm nl}(\scal)=\frac18\,8^{n_{\rm act}(\scal)}.}
\]
The second factor never replaces or interpolates the first.
In \((\mathrm{MRESP}_{\min})\), its assignment to the nonlinear Hessian is
the unique entropy-sufficient, tensor-multiplicative coupling of
Theorem~\ref{p3e:thm:entropy-effective-count-uniqueness}, followed by the
unique type-preserving constituent identification.  Co-realization alone,
without that minimal coupling signature, would not imply the assignment.
\end{theorem}

\begin{proof}
Paper~III-B Theorem~\ref{p3:thm:derived-scalar-corner-readout} co-realizes
the fixed rank-one scalar constituent and gives its ambient weight \(1/8\).
Theorem~\ref{p3:thm:index-locking} co-realizes the scale-dependent active
state and gives its entropy.  Theorem~\ref{p3e:thm:entropy-effective-count-uniqueness}
forces the dimensionless active fraction
\(w_{\rm nl}:=e^{S}/8\) in the minimal composition signature.  The displayed signatures differ in
support rank, scale dependence, and action role.  Paper~II-A
Corollary~\ref{p2:cor:canonical-constituent-identification} therefore gives
the unique matching of this scalar readout to the nonlinear Hessian role.
Because the linear signature is fixed, no generated
constructor permits replacing it by the scale-dependent active fraction.
\end{proof}

\begin{definition}[Primitive-quadratic trace inheritance certificate]
\label{p3e:def:primitive-quadratic-trace-inheritance}
On a successfully certified CRP4 parent, the positive quarter-root has a
well-defined tensor square. An actual same-action second-variation port must
compare its complete normalized primitive response with that tensor-square
owner, while retaining the common trace, kinetic/source Gram form and EH
coefficient. We call a successful origin-certified comparison
$(\mathrm{PQTI})$. The abstract tensor-square type does not itself give
this physical port or evaluate its relative amplitude. A real-channel
$\tau_8$ constituent and a complex Hilbert trace retain the precise
normalizations proved by Paper~III-B; a projector cannot be replaced by a
complex-linear one by dimension counting.
On $Z_C>0$, the coefficient test for the claimed terminal-slope-normalized amplitude is
\[
 \boxed{\mathfrak R_{\rm PQ}
 :=\frac{U\chi_{\rm UV}^{\rm act}}{Z_C}
       -16\sqrt{\alpha_\sigma}=0.}
\]
Here $\chi_{\rm UV}^{\rm act}$ is the susceptibility extracted from the
independently supplied complete action/source second jet, not from a function
already reconstructed with the proposed amplitude. The test is necessary and
sufficient for the scalar amplitude equality once the actual normalized
owner/source comparison exists. Its vanishing is to be evaluated, not imposed
as a replacement definition of the physical Hessian. The theorem below states
the algebraic normalization on that successful domain.
\end{definition}

\begin{theorem}[PQTI normalization of the ultraviolet quadratic seed]
\label{p3e:thm:pqti-quadratic-normalization}
Assume \((\mathrm{CRP4})+(\mathrm{PQTI})\) and a frozen
\(\mathrm{REG}_\sigma\in\mathfrak A_{\rm REG}^{+}\).  In the direct sum of
the independently typed Einstein--Hilbert line and the nonlinear primitive
quotient by \(\ker D^2\), the protected action representative is
\[
 \mathcal L_{\rm prim,raw}
 =\frac18\Lambda_{\rm EH}^2\scal
  +\sqrt{\alpha_\sigma}\,\mathscr Q_{C,U}(\scal).
\]
After the single whole-action normalization that makes the
Einstein--Hilbert coefficient \(\Lambda_{\rm EH}^2\), this becomes
\begin{equation}
 \label{p3e:eq:pqti-eh-normalized-primitive-action}
 \boxed{\mathcal L_{\rm prim}
 =\Lambda_{\rm EH}^2\scal
  +8\sqrt{\alpha_\sigma}\,\mathscr Q_{C,U}(\scal)}
 \qquad\text{in the nonlinear primitive quotient.}
\end{equation}
Consequently the MRESP ultraviolet susceptibility and its smooth protected
continuation are
\begin{equation}
 \label{p3e:eq:pqti-smooth-hessian-law}
 \boxed{\chi_{\rm UV}
 =\frac{16\sqrt{\alpha_\sigma}}{\Lambda_{\rm EH}^2},\qquad
 \mathfrak f''(\scal)
 =w_{\rm nl}(\scal)u(\scal)
  \frac{16\sqrt{\alpha_\sigma}}{\Lambda_{\rm EH}^2}.}
\end{equation}
On the generated action-response branch,
\[
 \boxed{\gamma_\sigma=16\sqrt{\alpha_\sigma}}
\]
in the terminal-slope-one representative, or invariantly
\(\widehat\gamma_\sigma=\gamma_\sigma/Z_C
=16\sqrt{\alpha_\sigma}\) after a common positive action rescaling.
On the enlarged independent-response atlas
\(U_\kappa=\kappa\Lambda_{\rm EH}^2\), the diagnostic covariance remains
\(\gamma_{\sigma,\kappa}=16\kappa\sqrt{\alpha_\sigma}\).
No coefficient \(7\sqrt\alpha,6\sqrt\alpha,\ldots\) is generated at later
hierarchy supports: those supports are descendants of the same nonlinear
insertion and alter \(w_{\rm nl}u\), not the primitive owner or its relative
trace normalization.
\end{theorem}

\begin{proof}
Paper~III-A Corollary~\ref{p3a:cor:causal-root-quadratic-half-character}
and Paper~III-C Corollary~\ref{p3c:cor:positive-associated-line-half-root}
give the unique positive tensor-square character
\(\alpha_\sigma^{1/2}\) of the CRP4 quarter-root.  PQTI transfers that
character, with no additional scalar, to the primitive response owner.
Paper~III-D Corollary~\ref{p3d:cor:primitive-unit-quadratic-response} gives
\(\mathscr Q_{C,U}''=2\).  Finally
Paper~III-B Theorems~\ref{p3:thm:normalized-protected-octet-trace} and
\ref{p3:thm:derived-scalar-corner-readout} give
\(\tau_8(I_8)=1\) and \(\tau_8(P_0)=1/8\).  Thus the raw relative
coefficients are exactly \(\sqrt{\alpha_\sigma}\) and \(1/8\); multiplying
the complete action once by eight proves
\eqref{p3e:eq:pqti-eh-normalized-primitive-action}.  Differentiating the
nonlinear primitive twice and dividing by the Einstein--Hilbert scale gives
\(16\sqrt{\alpha_\sigma}/\Lambda_{\rm EH}^2\).  The already derived MRESP
visibility and active fraction supply \(u\) and \(w_{\rm nl}\), proving
\eqref{p3e:eq:pqti-smooth-hessian-law}.  Substitution of
The generated initial-stage identity
\(U=\Lambda_{\rm EH}^2\) gives the displayed \(\gamma_\sigma\).
The affine kernel is never used in this calculation.  Since later rank
supports are evaluations of the same generated constituent rather than new
primitive second-variation owners, they cannot repeat the trace comparison.
\end{proof}

\subparagraph{Full-response versus second-variation scope.}
Equation~\eqref{p3e:eq:pqti-eh-normalized-primitive-action} fixes a class,
not a distinguished affine representative. In particular, replacing
\(\mathscr Q_{C,U}\) by \(\mathscr Q_{C,U}-2CU\) leaves this class unchanged
but changes its zero-order value. The weighted Born action on normalized
coefficients does not descend through \(\ker D^2\), whereas the Euler
kernel removes only the linear direction (Paper~I-D,
Proposition~\ref{p1b:prop:born-hessian-quotient-obstruction}). Thus PQTI
remains a susceptibility theorem; it does not by itself identify the full
positive response with either an additive Lorentzian action term or the
same action's Legendre potential.  At the present Paper-III-E Hessian stage,
Theorem~\ref{p3e:thm:regular-two-scale-terminal-reconstruction} shows that
the finite response Hessian plus a complete terminal jet reconstructs the
same-action potential, while
Corollary~\ref{p3e:cor:two-scale-affine-junction-family} classifies the
one-parameter family left when only the terminal value is retained.  Paper~III-F,
Theorem~\ref{p3f:thm:born-coefficient-family-first-response} and
Corollary~\ref{p3f:cor:constructor-born-discharge-bare}, subsequently go one
constructor level upstream: the generated scheme/depth coefficient connection
removes the constant affine mode and derives the regular terminal slope on the
bare same-action branch.  Paper~III-G, Appendix~\ref{p3e:app:born-affine-audit} remains the
negative audit showing why the Hessian quotient or fixed-point centering alone
cannot choose the affine member.
The constructive refinement is the paired Hessian
\((F'',(\mathcal T_{CU}F)'')\), whose kernel is only the linear
direction and which reconstructs \(RF'-F\) exactly (Paper~I-D,
Theorem~\ref{p1b:thm:paired-hessian-euler}). For
\(F=UR+a_{\rm nl}R^2+\beta\), it retains both
\(2a_{\rm nl}\) and \(2\beta/(CU)\); PQTI's first readout does
not assign the second. Reciprocal parity relates the quadratic and
constant coefficients but does not remove the trace/character factor
\(8\sqrt{\alpha_\sigma}\).
For a separately identified Euler response \(Q_d\), one has
\(F''=Q_d'/R\), which agrees with \(Q_d''\) only for a quadratic
plus constant (Proposition~\ref{p1b:prop:euler-response-dressing-compatibility}).
This Paper~I-D distinction retains the current susceptibility normalization and
prevents its dressed continuation from being silently replaced.


\subsection{Exact response character and calibrated gravitational placement}
\label{p3:sec:master-lagrangian-closed}

\begin{proposition}[Exact two-root curvature-response character]
\label{p3:def:reduced-curvature-transport-factor}
On the imported Paper~III-D reduced response branch, a positive response
normalization representative \(\kappa>0\) may be absorbed into
\[
 \scal_{\rm resp}^{(\kappa)}=\frac{\kappa}{2}(C+U).
\]
Every such representative yields the same calibrated invariant curve
\[
 \boxed{g_{\rm eff}(\scal)
 =\frac{\scal^2+CU}{(C+U)\scal}.}
\]
Paper~III-E adopts \(\kappa=2\), for which
\[
 \boxed{\scal_{\rm resp}=C+U,\qquad
        \scal_{\rm resp}\,\scal\,g_{\rm eff}(\scal)=\scal^2+CU.}
\]
Thus the factor \(2\) and the displayed \(\scal_{\rm resp}=C+U\) are a
normalization representative, not a surviving physical selector.  The
invariant response curve requires no Einstein--Hilbert placement.
\end{proposition}

\begin{proof}
This is the exact III-D v1.26 two-root response theorem together with its
normalization-representative invariance.  The final displayed identity is the
\(\kappa=2\) representative multiplied by \((C+U)\scal\).
\end{proof}


\begin{theorem}[Regular terminal two-scale reconstruction]
\label{p3e:thm:regular-two-scale-terminal-reconstruction}
Let \(0<C<U\) and let
\[
 \mathscr H_{C,U}(\scal)
 :=\scal g_{\rm eff}(\scal)
 =\frac{\scal^2+CU}{C+U}
\]
be the finite two-root response of
Paper~III-D, Corollary~\ref{p3d:cor:two-scale-response-normalizations}.
On the bare undressed response lane suppose that the response Hessian is
\[
 h_{C,U}(\scal)=\mathscr H_{C,U}''(\scal)=\frac{2}{C+U},
\]
and choose the regular terminal source member
\(
( F_C,Z_C)=(0,1)
\)
of the metric-response endpoint space.  Then
Theorem~\ref{p3e:thm:general-hessian-reconstruction} gives the unique
response coefficient
\begin{equation}
\label{p3e:eq:regular-two-scale-f}
 \boxed{
 \mathfrak f_{\rm reg}(\scal)
 =(\scal-C)+\frac{(\scal-C)^2}{U+C}.}
\end{equation}
Its complete same-action Euler/Legendre remainder is
\begin{equation}
\label{p3e:eq:regular-two-scale-legendre}
 \boxed{
 \mathcal V_{\rm reg}(\scal)
 :=\scal\mathfrak f_{\rm reg}'(\scal)-\mathfrak f_{\rm reg}(\scal)
 =\frac{\scal^2+CU}{U+C}
 =\mathscr H_{C,U}(\scal)
 =\scal g_{\rm eff}(\scal).}
\end{equation}
Equivalently, after multiplying the metric coefficient by the independently
owned Einstein--Hilbert action normalization \(U\),
\[
 \boxed{U\mathcal V_{\rm reg}=U\mathscr H_{C,U}
 =\frac{U}{U+C}(\scal^2+CU).}
\]
No reciprocal action-jet parity premise is used.  The ultraviolet datum in
this theorem is the Hessian; the two affine integration data are supplied
separately by the infrared terminal jet.  The theorem does not assert that
the active dressed Laurent action has constant Hessian, nor does it select
the regular terminal member over a sharp junction member.
\end{theorem}
\begin{proof}
The reconstruction theorem gives
\[
 \mathfrak f(\scal)
 =F_C+Z_C(\scal-C)
  +\int_C^{\scal}(\scal-s)h_{C,U}(s)\,ds,
\]
which with \((F_C,Z_C)=(0,1)\) and
\(h_{C,U}=2/(U+C)\) is
\eqref{p3e:eq:regular-two-scale-f}.  Differentiation gives
\[
 \mathfrak f_{\rm reg}'
 =1+\frac{2(\scal-C)}{U+C}.
\]
Hence
\[
 \scal\mathfrak f_{\rm reg}'-\mathfrak f_{\rm reg}
 =\frac{\scal^2+CU}{U+C},
\]
proving \eqref{p3e:eq:regular-two-scale-legendre}.  Equivalently,
\(\mathcal V_{\rm reg}'=\scal h_{C,U}
=\mathscr H_{C,U}'\), while at the terminal point
\(\mathcal V_{\rm reg}(C)=C=\mathscr H_{C,U}(C)\); equality on the
connected bare interval follows from one integration plus this IR value.
\end{proof}

\begin{corollary}[Affine junction family at the Hessian stage]
\label{p3e:cor:two-scale-affine-junction-family}
Every coefficient with the same bare Hessian and terminal value is
\[
 \mathfrak f_{\kappa}(\scal)
 =\mathfrak f_{\rm reg}(\scal)+\kappa(\scal-C),
\]
with terminal slope \(\mathfrak f_\kappa'(C)=1+\kappa\).  Its same-action
potential is
\[
 \boxed{\mathcal V_\kappa(\scal)
 =\mathscr H_{C,U}(\scal)+\kappa C.}
\]
Therefore the Hessian fixes neither the terminal slope nor the additive
potential shift.  Within this Hessian classification the regular \(C^1\)
member is exactly \(\kappa=0\).  This is not a live downstream selector on
the generated bare same-action branch: Paper~III-F,
Corollary~\ref{p3f:cor:constructor-born-discharge-bare}, derives
\(\kappa=0\) from the coefficient-family constructor while retaining only the
already owned terminal value.  The present corollary remains the exact
classification of what the second derivative alone cannot determine.
For the sharp undressed TERM--REN representative
\[
 \mathfrak f_{\rm sharp}(\scal)
 =\scal+\mathscr H_{C,U}(\scal)-2C,
\]
one has
\[
 \boxed{\kappa_{\rm sharp}=\frac{2C}{U+C}},
\]
so the sharp and regular members have identical second derivative but
different terminal first derivative.
\end{corollary}
\begin{proof}
The affine-kernel statement is
Corollary~\ref{p3e:cor:terminal-data-minimality}.  Applying
\(\scal\partial_{\scal}-1\) to
\(\kappa(\scal-C)\) gives \(\kappa C\).  Finally direct subtraction yields
\[
 \mathfrak f_{\rm sharp}-\mathfrak f_{\rm reg}
 =\frac{2C}{U+C}(\scal-C).
\]
\end{proof}

\begin{theorem}[Canonical covariant Hessian pullback]
\label{p3e:thm:canonical-covariant-hessian-pullback}
Let \(W\) be a finite-dimensional real vector space, let
\(A:W\to W\) be a generated amplitude map, and let a Hessian carry its
ordinary unmarked covariant type
\(H\in\operatorname{Sym}^2(W^*)\).  The tensor functor supplies the unique
map induced by \(A\) without an additional scalar/density owner:
\[
 \boxed{\mathcal T_A(H)=\operatorname{Sym}^2(A^*)H,
 \qquad
 (\mathcal T_AH)(x,y)=H(Ax,Ay).}
\]
Equivalently, for every \(C^2\) scalar response \(\Phi\) near the origin,
\[
 \operatorname{Hess}_0(\Phi\circ A)
 =\mathcal T_A(\operatorname{Hess}_0\Phi).
\]
For \(A_\tau=V_\tau^{1/2}|_W=e^{-\tau/2}I_W\), this forces
\[
 \boxed{\mathcal T_\tau(H)=e^{-\tau}H.}
\]
The family preserves positive semidefiniteness and obeys
\(\mathcal T_{\tau+s}=\mathcal T_\tau\mathcal T_s\).  Hence
\(\mathcal T\) is not a selector in \((\mathrm{MRESP}_{\min})\).  A different
exponent or a general completely positive map requires an additional
character, tensor, or response owner and belongs to the enlarged role.
\end{theorem}

\begin{proof}
The dual map is \(A^*:W^*\to W^*\).  On a decomposable symmetric tensor it
acts as
\[
 \ell\odot m\longmapsto(A^*\ell)\odot(A^*m).
\]
Such tensors span \(\operatorname{Sym}^2(W^*)\), proving existence and
uniqueness of the induced linear tensor map.  The second-derivative chain rule
for the linear map \(A\) gives the Hessian identity; there is no term involving
\(D^2A\).  Substituting \(A_\tau=e^{-\tau/2}I_W\) contributes one factor on
each covariant leg, hence \(e^{-\tau}\).  Pullback preserves
positive-semidefinite forms, and functoriality of pullback proves the semigroup
law.  Multiplying this tensor pullback by any nontrivial scalar character is
not induced by \(A\) on \(\operatorname{Sym}^2(W^*)\); it is precisely an
extra one-dimensional owner excluded by the unmarked hypothesis.
\end{proof}

\begin{theorem}[Dimensionless homothety of a Hessian response]
\label{p3e:thm:response-homothety}
Put \(c=C/U\), \(x=\scal/U\), and suppose
\[
 h(\scal):=\mathfrak f''(\scal)
 =\frac1U\widehat h(x),\qquad
 F_C=U\widehat F_C,
\]
with \(Z_C=\mathfrak f'(C)\) dimensionless.  Then
\[
 \boxed{\mathfrak f(\scal)
 =U\left[
 \widehat F_C+Z_C(x-c)
 +\int_c^x(x-y)\widehat h(y)\,dy\right].}
\]
Consequently
\[
 (C,U,\scal)\longmapsto
 (\lambda C,\lambda U,\lambda\scal),\qquad \lambda>0,
\]
preserves the complete dimensionless response shape.  Moreover
\[
 \mathfrak f'(\scal)=\widehat f'(x),\qquad
 \mathfrak f''(\scal)=U^{-1}\widehat f''(x),
\]
so scalaron masses scale as \(U\) at the level of mass squared whenever the
standard scalaron formula applies.  This theorem is a homothety statement and
does not itself use the later generated identity-stage normalization
\(\kappa_U=1\).
\end{theorem}

\begin{proof}
Apply Theorem~\ref{p3e:thm:general-hessian-reconstruction} and substitute
\(s=Uy\), \(\scal=Ux\), and
\(h(Uy)=U^{-1}\widehat h(y)\).  The derivative formulas follow by the chain
rule, and every term in the scalaron mass-squared ratio acquires one factor
of \(U\).
\end{proof}

\begin{corollary}[Homothety classification of the ultraviolet susceptibility]
\label{p3e:cor:uv-susceptibility-homothety}
Suppose the response Hessian factors as
\[
 h_{U,C}(\scal)
 =w(u(\scal))\,u(\scal)\,\chi_{\rm UV}(U,\sigma),
\]
where \(w\) and \(u\) are dimensionless and invariant under
\((C,U,\scal)\mapsto(\lambda C,\lambda U,\lambda\scal)\).  Then the complete
homothety law of Theorem~\ref{p3e:thm:response-homothety} holds for all
\(\lambda>0\) if and only if
\[
 \boxed{\chi_{\rm UV}(U,\sigma)=\frac{\gamma_\sigma}{U}}
\]
for one dimensionless \(\gamma_\sigma\), positive exactly when the
susceptibility is positive.  Homothety fixes no numerical value of
\(\gamma_\sigma\).  This classifies the unrestricted MRESP atlas;
Theorem~\ref{p3e:thm:pqti-quadratic-normalization} later cuts out the PQTI
subfamily by an additional action-content relation.
\end{corollary}

\begin{proof}
Homothety requires
\(h_{\lambda U,\lambda C}(\lambda\scal)=\lambda^{-1}h_{U,C}(\scal)\).
The two dimensionless factors cancel, leaving
\(\chi_{\rm UV}(\lambda U,\sigma)=\lambda^{-1}
\chi_{\rm UV}(U,\sigma)\).  Thus
\((\lambda U)\chi_{\rm UV}(\lambda U,\sigma)
=U\chi_{\rm UV}(U,\sigma)\); for fixed \(\sigma\) the product is a constant
\(\gamma_\sigma\).  The converse follows by substitution.  Since every
positive \(\gamma_\sigma\) obeys the same covariance, no value is selected.
\end{proof}

\begin{corollary}[Generated-owner and scheme-explicit response]
\label{p3:lem:curvature-transport-character}
Under \((\mathrm{CAL}_{\rm resp/EH})\), \((\mathrm{CRP4})\), and
\(\mathrm{REG}_{\sigma}\in\mathfrak A_{\rm REG}^{+}\),
\[
 U=\Lambda_{\rm EH}^2,\qquad
 C=\frac{\Lambda_{\rm EH}^2}
 {\alpha_{\rm hier}(\sigma_{\rm fin})^{7}},\qquad
 \scal_{\rm resp}=\Lambda_{\rm EH}^2
 \bigl(1+\alpha_{\rm hier}(\sigma_{\rm fin})^{-7}\bigr),
\]
and
\[
 \scal_{\rm resp}\,\scal\,g_{\rm eff}(\scal)
 =\scal^2+(\Lambda_{\rm EH}^2)^2
 \alpha_{\rm hier}(\sigma_{\rm fin})^{-7}.
\]
The midpoint
\(\scal_B=\Lambda_{\rm EH}^2
\alpha_{\rm hier}(\sigma_{\rm fin})^{-7/2}\) is a Born-even response scale,
not a geometric radius unless another theorem supplies that interpretation.
With the legacy alias \(\Rp:=U=\Lambda_{\rm EH}^2\), these formulas may be
written equivalently with \(\Rp\).
On the generated octahedral-interface branch,
\((\mathrm{REG}_0):\sigma_{\rm fin}=0\) gives the anchored point values.
\end{corollary}

\begin{proof}
Paper~III-C's owner-typed inheritance corollary gives
\(C/U=\alpha_{\rm hier}(\sigma_{\rm fin})^{-7}\).  Substitute this ratio and
the generated identity-stage value \(U=\Lambda_{\rm EH}^2\) into the exact
response character above.  No use of \((\mathrm{CAL}_{xp})\), second ratio
match, or additional normalization is made.
\end{proof}

\begin{remark}[No monomiality claim from \(\Dcl\)]
\label{p3:lem:homogeneous-running-monomial}
The exact response character is the displayed two-term quantity
\(\scal^2+CU\).  Paper~III-E makes no claim that \(\Dcl\) alone excludes
every other nonlinear scalar function or forces a monomial interpolation.
Such a restriction would require an additional character/multiplicativity
hypothesis and is not used below.
\end{remark}

\subsection{Conditional smooth metric response}

\begin{theorem}[General Hessian reconstruction from typed endpoint data]
\label{p3e:thm:general-hessian-reconstruction}
Let \(0<C<U\), let \(h\in C([C,U])\), and specify a curvature-valued
\(F_C\in\mathcal L_{\scal}\) and a dimensionless \(Z_C\in\mathbb R\).
There is a unique
\(\mathfrak f\in C^2([C,U])\) satisfying
\[
 \mathfrak f''(\scal)=h(\scal),\qquad
 \mathfrak f(C)=F_C,\qquad
 \mathfrak f'(C)=Z_C,
\]
namely
\[
 \boxed{\mathfrak f(\scal)
 =F_C+Z_C(\scal-C)
 +\int_C^{\scal}(\scal-s)h(s)\,ds.}
\]
If \(h\ge0\), the reconstructed response is convex.  This theorem requires no
visibility semigroup, entropy coupling, hierarchy match, or Wick
reconstruction; those additional data specify the particular \(h,C,U,F_C,Z_C\) member.
\end{theorem}

\begin{proof}
Two applications of the fundamental theorem of calculus give the displayed
formula.  Conversely, differentiating it twice recovers \(h\) and the two
endpoint values.  The difference of two solutions is affine and has zero
value and derivative at \(C\), hence vanishes.  Nonnegative second derivative
is equivalent to convexity.
\end{proof}

\begin{corollary}[Exact affine kernel and minimality of terminal data]
\label{p3e:cor:terminal-data-minimality}
For fixed \(h\in C([C,U])\), the complete solution set of
\(\mathfrak f''=h\) is the two-dimensional affine space
\[
 \mathfrak f_h+\ker D^2
 =\mathfrak f_h+\operatorname{span}_{\mathbb R}\{1,\scal-C\}.
\]
The endpoint map
\[
 \operatorname{ev}_C:\{\mathfrak f:\mathfrak f''=h\}\longrightarrow
 \mathbb R^2,\qquad
 \mathfrak f\longmapsto(\mathfrak f(C),\mathfrak f'(C))
\]
is an affine bijection.  Consequently a unique reconstruction by continuous
linear scalar boundary conditions requires two independent conditions: one
consistent scalar condition leaves an affine solution set whose translation
kernel has dimension at least one.  Thus the raw
Paper-III-E pair \((F_C,Z_C)\) cannot be deduced from the Hessian.  It can be
reduced only by an additional physical relation or by an exact quotient in a
larger action.
\end{corollary}

\begin{proof}
If \(\mathfrak f_1''=\mathfrak f_2''=h\), then
\((\mathfrak f_1-\mathfrak f_2)''=0\), so their difference is
\(a+b(\scal-C)\).  Conversely every such affine addition preserves the
Hessian.  Evaluation sends this kernel basis to \((1,0)\) and \((0,1)\), so
it is bijective.  A system of \(m\) continuous linear scalar boundary
conditions restricts the two-dimensional kernel by a linear map to
\(\mathbb R^m\).  Whenever the inhomogeneous conditions are consistent, the
translation kernel has dimension \(2-r\), where \(r\) is the rank of that
map.  Uniqueness therefore requires rank two and at least two independent
conditions.
\end{proof}

\begin{theorem}[UV--IR moments, gluing and smooth interior nonuniqueness]
\label{p3e:thm:v2217-uv-ir-moment-nonuniqueness}
Let $0<C<U$, $f\in C^2([C,U])$, $F=f'$ and $\chi=f''\ge0$.
Prescribed endpoint values of $f,F$ necessarily obey
\[
 M_0=F(U)-F(C)=\int_C^U\chi(s)\,ds,\qquad
 M_1=[sF(s)-f(s)]_C^U=\int_C^U s\chi(s)\,ds.
\]
Conversely, a specified continuous $\chi$ and the lower endpoint data
match the upper endpoint exactly if and only if these two equalities hold.
For $M_0>0$ one has $CM_0\le M_1\le UM_0$; without prescribed collar
germs every strictly interior ratio $M_1/M_0\in(C,U)$ is attained by a
smooth nonnegative susceptibility. With prescribed collars, their moments
must first be subtracted and the residual test is on the free interval.
If $\chi>0$ on an interior open interval, then every nonzero
$u\in C_c^\infty$ of that interval gives, for sufficiently small real
$\epsilon$, a different $f_\epsilon=f+\epsilon u$ with the same complete
endpoint germs, both moments, all unchanged interface jets and nonnegative
susceptibility. On a connected overlap two $C^2$ action charts with equal
susceptibility agree exactly when their value and first derivative agree
at one common point.
\end{theorem}
\begin{proof}
Integrate $F'=\chi$ and $(RF-f)'=R\chi$. Conversely the first equality
matches $F(U)$, and the second then matches $f(U)$ by the reconstruction
formula. Positivity bounds the weighted mean of $s$ between $C,U$.
A nonnegative even smooth bump of mass $M_0$ centered at $M_1/M_0$,
with support inside $(C,U)$, proves the interior existence claim. Fixed
collars and any fixed smooth transitions contribute known moments and can
be subtracted; no claim dispenses with their matching jets.
On the compact support of $u$, let $m=\min\chi>0$. Choosing
$|\epsilon|\|u''\|_\infty<m$ preserves positivity there and leaves it
unchanged elsewhere. The bump is zero on endpoint collars and
$\int u''=0$, $\int s u''=[su'-u]_C^U=0$, proving all invariances.
It is not affine, so the action differs. Finally the difference of charts
with equal Hessian is affine; its zero value and slope at one point make
it zero everywhere on the connected overlap.
\end{proof}
The endpoint nonuniqueness is in the smooth category; a full analytic germ
on a connected analytic domain is a different uniqueness problem.
Constitutive completion does not certify the independent microscopic word.


\begin{lemma}[Protected curvature-Hessian law in the minimal response signature]
\label{p3:lem:protected-hessian-transport}
Assume \((\mathrm{MRESP}_{\min})\).  Let \(W\) be the seven-dimensional protected
non-scalar sector.  Then
\[
 \boxed{C_{\rm sp}''(\scal)
 =w_{\rm nl}(\scal)\,\frac{\gamma_\sigma}{U}\,u(\scal)},
 \qquad
 w_{\rm nl}(\scal)=\frac18\,8^{n_{\rm act}(\scal)}.
\]
Here \(u(\scal)=\ln(\scal/C)/\ln(U/C)\) is the smooth visibility coordinate
generated parametrically in \((\mathrm{MRESP}_{\min})\) and inverted in
Theorem~\ref{p3:thm:index-locking}.  Thus the Einstein--Hilbert constituent
keeps its fixed \(1/8\), while the nonlinear Hessian carries the separately
identified active fraction.
\end{lemma}

\begin{proof}
Theorem~\ref{p3e:thm:canonical-covariant-hessian-pullback} derives exactly
one tensorial visibility factor
\(e^{-\tau}=u(\scal)\).  Theorem~\ref{p3e:thm:entropy-effective-count-uniqueness}
derives the only entropy-sufficient multiplicative octet fraction,
\(w_{\rm nl}=e^S/8\), and
Theorem~\ref{p3:rem:index-locking-status} gives its unique type-preserving
nonlinear constituent.  Finally
Corollary~\ref{p3e:cor:uv-susceptibility-homothety} gives
\(\chi_{\rm UV}=\gamma_\sigma/U\).  Multiplication of these three typed scalar
readouts yields the displayed Hessian.  At \(U\), \(u=w_{\rm nl}=1\), so its
value is \(\gamma_\sigma/U\).  No arbitrary transport or coupling functional
 is selected.
\end{proof}

\begin{theorem}[Complete minimal smooth-response family]
\label{p3e:thm:complete-minimal-smooth-response-family}
Fix any \(\sigma\in I_+\), \(U>0\), the certified CRP4 response ancestry, ultraviolet calibration
\(\gamma_\sigma>0\), and endpoint pair
\((F_C,Z_C)\in\mathcal L_{\scal}\times\mathbb R\).  Put
\[
 C=U\alpha_\sigma^{-7},\qquad
 h_{\sigma,\gamma}(\scal)
 :=\frac18\,8^{n_{\rm act}(\scal)}
   \frac{\gamma_\sigma}{U}\,u(\scal).
\]
There is exactly one minimal response with these data:
\begin{equation}
\label{p3e:eq:complete-minimal-response-family}
 \boxed{\mathfrak f_{\sigma,\gamma,F,Z}(\scal)
 =F_C+Z_C(\scal-C)
 +\frac{\gamma_\sigma}{U}
  \int_C^\scal(\scal-s)\,
  \frac18\,8^{n_{\rm act}(s)}u(s)\,ds.}
\end{equation}
For every member of this parameter bundle it satisfies
\[
 \mathfrak f''_{\sigma,\gamma,F,Z}>0\quad(C<\scal\le U),
 \qquad
 \mathfrak f''_{\sigma,\gamma,F,Z}\longrightarrow0^+
 \quad(\scal\downarrow C),
\]
No uniform-convergence assertion over the unbounded parameter bundle is
intended.  Convexity and terminal Hessian decay require neither the source REG evaluation, the source endpoint
pair, nor the source ultraviolet calibration.
\end{theorem}

\begin{proof}
The state formula and \(u(\scal)\) are continuous on \([C,U]\) with
\(u(C)=0\), and \(w_{\rm nl}=e^{S}/8\) lies in \([1/8,1]\).
Lemma~\ref{p3:lem:protected-hessian-transport} therefore gives a continuous
Hessian, strictly positive on \((C,U]\), tending to zero at \(C\).
Theorem~\ref{p3e:thm:general-hessian-reconstruction} yields
\eqref{p3e:eq:complete-minimal-response-family} and uniqueness.  The proof
uses \(\sigma,\gamma_\sigma,F_C,Z_C\) as arbitrary typed parameters, so none
of their source values is required for the asserted properties.
\end{proof}

\begin{corollary}[Witness transport for the complete minimal response]
\label{p3e:cor:mresp-composite-witness-transport}
Let
\[
 D_{\rm MRESP}
 :=(\mathrm{REG}_\sigma,U,\mathrm{CRP4},
    \gamma_\sigma,F_C,Z_C,\mathcal V(P_0))
\]
be a jointly occurring/generated input tuple in the minimal response
signature, with
\(\sigma\in I_+\) and \(U,\gamma_\sigma>0\).  At the beginning of this
chronological stage, freeze the response-reconstruction schema licensing
multiplication of the retained response factors, affine endpoint-jet
addition, and twice-iterated integration.  Declare before evaluation the
target sort \(\mathsf{RespGraph}_{[C,U]}^{2,+}\) of twice continuously
differentiable convex response graphs on \([C,U]\) with the displayed endpoint
jet and generated Hessian.  Let \(w_{\rm MRESP}\) be the finite typed word
\[
 D_{\rm MRESP}\longmapsto
 \mathfrak f''_{\sigma,\gamma,F,Z}
 \longmapsto
 F_C+Z_C(\scal-C)
 +\int_C^\scal\!\int_C^u
   \mathfrak f''_{\sigma,\gamma,F,Z}(v)\,dv\,du .
\]
All intermediate factors and integrals are retained in their native sorts,
and the displayed target and required output are fixed before evaluation.
Thus
\[
 \operatorname{Orig}_{\Sigma_{\rm MRESP}}
 (w_{\rm MRESP};D_{\rm MRESP},
  \mathsf{RespGraph}_{[C,U]}^{2,+})
\]
holds.  Then the constructor
\[
 D_{\rm MRESP}\longmapsto
 \mathfrak f_{\sigma,\gamma,F,Z}
\]
of Theorem~\ref{p3e:thm:complete-minimal-smooth-response-family} transports
the complete input witness and co-realizes its unique response graph.  No
separate occurrence assumption for the reconstructed function is required.

This conclusion is conditional on occurrence of the complete tuple, not on
its proper subtuples.  It selects no value of
\(\sigma,\gamma_\sigma,F_C,Z_C\), supplies no missing owner ancestry, and
does not adopt the response as a Lorentzian gravitational action or assert a
solution history.  Those are distinct declared output sorts.
\end{corollary}

\begin{proof}
The frozen word \(w_{\rm MRESP}\) uses only multiplication, the already generated
visibility/entropy/Hessian factors, and the twice-iterated integral with the
given endpoint jet; its intermediate Hessian and two integral outputs remain
retained and reusable.  The affine-kernel theorem proves that its output is the
unique graph in the declared target sort, while the Hessian classification
proves that no transport or coupling function is inserted.  Evaluation and
integration are natural under the unmarked owner maps, and differentiating
twice together with evaluating at \(C\) reflects the complete input data.
Thus the constructor is conservative and its essential image is exactly the
displayed graph.  Continuity on the compact interval and the uniqueness
theorem exhibit a GP1-finite-and-bounded, GP2-exact-and-faithful realization,
so \(\operatorname{Obs}_{\rm FRC}(w_{\rm MRESP},D_{\rm MRESP})\) is false.
Paper~II-A,
Theorem~\ref{p2a:thm:rooted-generated-image-corealization}, gives the stated
co-realization.  Naturality and exactness verify the realization rather than
licensing the word.  Its declared-output boundary gives the final firewall.
\end{proof}

\begin{lemma}[Source terminal protected slope from \((\mathrm{TERM}_f)\)]
\label{p3:lem:terminal-protected-slope}
Under the explicit source endpoint member
\(\mathrm{TERM}_f^0:(F_C,Z_C)=(0,1)\), the nonlinear protected
response satisfies
\begin{equation}
\label{p3:eq:terminal-protected-zero-slope}
 \boxed{C_{\rm sp}'(C)=0.}
\end{equation}
\end{lemma}

\begin{proof}
By definition
\(\mathfrak f(\scal)=\scal+C_{\rm sp}(\scal)\).  Therefore
\[
 Z_C=\mathfrak f'(C)=1+C_{\rm sp}'(C).
\]
The independently declared source value \(Z_C=1\) gives
\(C_{\rm sp}'(C)=0\).  No semigroup or first-variation transport law is
used.
\end{proof}

\begin{lemma}[Source endpoint correction under \((\mathrm{TERM}_f)\)]
\label{p3:lem:terminal-hamiltonian-lagrangian-sign}
Under the explicit source member \(\mathrm{TERM}_f^0:(F_C,Z_C)=(0,1)\),
\begin{equation}
\label{p3:eq:terminal-legendre-sign}
 \boxed{C_{\rm sp}(C)=-C.}
\end{equation}
\end{lemma}

\begin{proof}
Writing \(\mathfrak f(\scal)=\scal+C_{\rm sp}(\scal)\), the typed datum
\(\mathfrak f(C)=F_C=0\) gives \(C_{\rm sp}(C)=-C\) algebraically.  No
Hamiltonian or Legendre narrative is required.  Such a narrative may be
supplied as a physical interpretation, but it does not strengthen the
mathematical endpoint statement.
\end{proof}

\begin{theorem}[Source-member specialization of the minimal response]
\label{p3:thm:master-lagrangian}
Assume \((\mathrm{MRESP}_{\min})\) and the source endpoint member.  The unique twice continuously differentiable
nonlinear protected correction with the Hessian of
Lemma~\ref{p3:lem:protected-hessian-transport} and terminal data
\(C_{\rm sp}(C)=-C\), \(C_{\rm sp}'(C)=0\) is
\begin{equation}
\label{p3:eq:spectral-response-reconstruction}
 \boxed{
 C_{\rm sp}(\scal)
 =-C+\int_C^{\scal}(\scal-s)
       \frac18\,8^{n_{\rm act}(s)}
       \frac{\gamma_\sigma}{U}\,u(s)\,ds .}
\end{equation}
Define
\[
 \boxed{\mathfrak f_{\rm sp}(\scal):=\scal+C_{\rm sp}(\scal).}
\]
Then
\[
 \mathfrak f_{\rm sp,\scal\scal}(\scal)>0
 \quad(C<\scal\le U),
 \qquad
 \mathfrak f_{\rm sp,\scal\scal}(\scal)\to0^+
 \quad(\scal\downarrow C),
\]
and
\[
 \mathfrak f_{\rm sp,\scal}(C)=1.
\]
This is the specialization of
Theorem~\ref{p3e:thm:complete-minimal-smooth-response-family} to
\(\mathrm{TERM}_f^0:(F_C,Z_C)=(0,1)\).  It is not a
classification of all generally covariant metric actions.
\end{theorem}

\begin{proof}
The factors \(w_{\rm nl}\), \(\gamma_\sigma/U\), and \(u\) are continuous and
strictly positive on \((C,U]\).  Since
\(w_{\rm nl}(C)=1/8\) and \(u(C)=0\), their product tends to zero at terminal
resolution.  Apply
Theorem~\ref{p3e:thm:complete-minimal-smooth-response-family} with the source
endpoint data; subtracting the linear Einstein--Hilbert term gives the
displayed correction.  The
linear Einstein--Hilbert term contributes derivative one, while the nonlinear
correction has zero terminal derivative, giving
\(\mathfrak f_{\rm sp,\scal}(C)=1\).
\end{proof}


\subsection{Carrier-selected rank-four ultraviolet coefficient}
\label{p3e:subsec:v2219-carrier4-current}
In the inherited grade-four response comparator, the first constituted
nonzero carrier is $P_4$, \emph{not} the optional least-support Laurent
choice. This positive-chart response is retained under its own hypotheses,
but it is not the current capped gravitational working density. In the
positive scalar-response chart $R>C_7$ put
\begin{equation}\label{p3e:eq:v2219-carrier-four-response}
\boxed{D_4(R)=Q^{n_4^+(R)},\quad
\chi_{4\mid7}(R)=D_4(R)
\frac{(1+n_7^+(R))(R^2-C_7^2)}{(U+C_7)R^2},\quad Q=8\sqrt\alpha.}
\end{equation}
The complete terminal-normalized one-variable \emph{response representative},
not a global two-sided Lorentzian UV action, is
\begin{equation}\label{p3e:eq:v2219-carrier-four-legendre}
\boxed{f_{4\mid7}^{+}(R)=R-C_7+
\int_{C_7}^{R}(R-s)\chi_{4\mid7}(s)\,ds.}
\end{equation}
\begin{theorem}[Carrier-four interior midpoint and displaced coefficient freeze]
\label{p3e:thm:v2219-carrier-four-jets}
On the positive active interval, $\chi_{4\mid7}>0$, the terminal jets
are $f(C_7)=0$, $f_R(C_7)=1$, and $f_{RR}(C_7)=0$.
The rank-four reciprocal Born midpoint is
\[
 \boxed{R_{B,4}=\sqrt{UC_4}=C_2=E_{32}^2,\quad
 D_4(C_2)/Q=Q^{n_4^+(C_2)-1},\quad
 [f_{RRR}]_{C_2}=0.}
\]
At the anchored $\mathrm{REG}_0$ scheme the preceding exact ratio
evaluates to $0.825729956674\ldots$; that decimal is \emph{not}
uniform over the positive $\mathrm{REG}_\sigma$ locus.
The first coefficient-induced third-jet jump occurs \emph{instead} at
$R=C_4$ (structural label $K_{\rm sh}=32$):
\begin{equation}\label{p3e:eq:v2219-c4-third-jet}
\boxed{[f_{RRR}]_{C_4}=
\frac{(U-C_4)\ln Q}{4C_4(U+C_4)\ln\alpha}\,\chi_7(C_4)>0,\quad
 [f_{RR}]_{C_4}=0.}
\end{equation}
Hence $K_{\rm sh}=16$ is a regular internal Born balance, not a coefficient
freeze; the physical positive-curvature Laurent coordinate of $R=C_2$ is
$K_R=32$, not $K_{\rm sh}=16$.
\end{theorem}
\begin{proof}
Twice integrating the continuous positive Hessian gives the claimed
terminal conditions and convexity.  The Laurent positive part vanishes
at $q=4$, where the right derivative of $n_4^+$ with respect to $R$ is
$\tanh(2\ln\alpha)/(4C_4\ln\alpha)$ and the derivative below $C_4$
vanishes.  At $q=2$, both $n_4$ and $n_7$ are interior analytic functions,
so no derivative jump exists.  The reciprocal fixed point follows from
$R^2=UC_4$.  No smoothing function or new carrier is inserted.
\end{proof}


\begin{proposition}[Capped gravity is not the grade-four Laurent density on a shared curvature chart]
\label{p3e:prop:v288-cap-laurent-nonidentification}
Retain the positive-chart grade-four response of
Theorem~\ref{p3e:thm:v2219-carrier-four-jets} and the complete capped
scalar density of Theorem~\ref{p3e:thm:v223-exact-cap-reduction}.
If the capped parameters satisfy $2B_U>d:=1-B_U/B_I$, then
\[
 \lim_{R\downarrow C_7}(f_{4\mid7}^{+})_{RR}(R)=0,
 \qquad (f_c)_{RR}(C_7)\ge 2B_U-d>0.
\]
Consequently the two densities cannot agree near $C_7$ in one and the
same curvature variable after an affine counterterm and a nonzero
constant overall normalization. Moreover $f_c$ is real analytic for
all real $R$, whereas the grade-four comparator has the nonzero
third-jet jump at $C_4$ of
Theorem~\ref{p3e:thm:v2219-carrier-four-jets}; this independently
forbids their equality on any common interval spanning $C_4$.
Neither conclusion forbids different reductions of a more complete
common action with a separately derived change of variables or owner.
\end{proposition}
\begin{proof}
The terminal Hessian of the inherited comparator is explicitly zero by
Theorem~\ref{p3e:thm:v2219-carrier-four-jets}.
Proposition~\ref{p3f:prop:v223-cap-legendre} gives the strictly positive
uniform lower bound for the cap Hessian. An affine counterterm has
zero second derivative, so no finite constant normalization can
identify these different Hessian limits. The square-root profile
$\sqrt{1+(R/R_s)^2}$ is analytic and positive for real $R$, as are
$B(R)$ and $B(R)^{-1}$; thus $f_c$ is real analytic. The displayed
finite nonzero third-jet jump of the Laurent comparator cannot occur
in an analytic function. Both arguments use one shared curvature chart.
\end{proof}

\begin{theorem}[Exact closed-three-sphere acceleration Hessian of an owned metric response]
\label{p3e:thm:v287-s3-acceleration-hessian}
Let $Z>0$ be the fixed normalization of one local four-dimensional
metric-$f(R)$ action, $f\in C^2(I)$, and restrict its metric density
$Z\sqrt{-g}\,f(R)$ to the $k=+1$ FLRW jet chart
\[
 ds^2=-N(t)^2dt^2+a(t)^2d\Omega_3^2,
 \qquad a>0,\quad N>0,\quad R\in I.
\]
With the unit round $S^3$ volume $2\pi^2$ and the paper's Lorentzian
curvature sign convention, its unreduced gravitational Lagrangian is
\begin{align}
 L_f&=2\pi^2ZNa^3 f(R),\nonumber\\
 R&=6\left(
 \frac{\ddot a}{aN^2}+\frac{\dot a^2}{a^2N^2}
 +\frac1{a^2}-\frac{\dot a\dot N}{aN^3}\right).
 \label{p3e:eq:v287-s3-flrw-curvature}
\end{align}
The \emph{local acceleration-jet coefficient} at fixed
$(a,\dot a,N,\dot N)$ is exactly
\begin{equation}\label{p3e:eq:v287-s3-acceleration-hessian}
 \boxed{\frac{\partial^2 L_f}{\partial(\ddot a)^2}
       =\frac{72\pi^2Za}{N^3}\,f_{RR}(R).}
\end{equation}
In particular it is nonzero exactly when $f_{RR}(R)\ne0$.
The identity survives the addition of a total time derivative
$dB(a,\dot a,N,\dot N,t)/dt$ with $B$ independent of accelerations,
because this changes $L_f$ at most affinely in $\ddot a$.
It does \emph{not} assert that this unreduced coefficient is a
propagating, constrained, or experimentally observed scalar degree of
freedom. Additional higher-derivative action terms must have their own
complete acceleration jets; an abstract response representative is not
thereby promoted to an occurring metric action.
\end{theorem}
\begin{proof}
The inherited round $S^3$ has volume $2\pi^2$; integrating
$\sqrt{-g}=Na^3\sqrt{\gamma}$ gives the stated $L_f$.
In the jet chart $R$ is affine in $\ddot a$ and
$\partial R/\partial\ddot a=6/(aN^2)$,
$\partial^2R/\partial(\ddot a)^2=0$.
Twice applying the ordinary chain rule at fixed remaining jets gives
\[
 L_{f,\ddot a\ddot a}
 =2\pi^2ZNa^3\,f_{RR}(R)
 \left(\frac6{aN^2}\right)^2
 =72\pi^2ZaN^{-3}f_{RR}(R).
\]
A total first-jet boundary derivative is affine in $\ddot a$.
The claims concern the stated local curvature-action chart, not its
physical occurrence, full constraint reduction, or additional sources.
\end{proof}

\begin{remark}[Rank-two comparison does not select the rank-seven response]
The following depth-two subsection, its named $K16$ third-jet
kink, and the $\Theta_2$ readout are retained solely as the
\emph{counterfactual $n_2\mid n_7$ comparison} and for label
traceability.  They are not the active rank-seven response branch.  Wherever
a rank-two comparison sentence calls $n_2$ ``active'', ``unique'' or the physical
inflation dressing, its status is restricted by the stated rank-two comparison scope.
\end{remark}


\subsection{Local source certification versus global metric response}
\label{p3e:sec:v204-local-global-response}
\begin{corollary}[Local source certification and global response have distinct scope]
\label{p3e:cor:v2216-local-global-response-scope}
A successful $(\mathrm{EXJET}_7)$ calculation certifies the actual frozen
source family's local second jet on its declared domain. Its use as the
Hessian of a global metric coefficient additionally requires an actual
second-jet transport/effective constructor, its physical normalization,
and the complete endpoint, source and boundary matching. It does not
identify every complete local source germ with that global coefficient.

In particular the terminally normalized metric candidate is
\[
 f_{2\mid7}(R)=R-C+\int_C^R(R-t)D_2(t)\chi_7^{\rm jet}(t)\,dt,
\]
with physical Euler readout
\[
 P_{2\mid7}(R)=C+\int_C^R tD_2(t)\chi_7^{\rm jet}(t)\,dt.
\]
These are unique on the supplied second-jet and endpoint data. The raw
Laurent readout is not their metric Euler source. Neither the positive
resolvent provider nor the finite $1\oplus7$ nonlinear trace of Paper~III-D
removes the actual-action matching premise.
\end{corollary}
\begin{proof}
Twice integrating the supplied continuous susceptibility fixes $f$ up to
an affine function; value zero and slope one at $C$ fix that kernel.
Differentiation of $Rf'-f$ gives $Rf''$, proving the Euler formula.
On the lower interval $D_2=1$,
Theorem~\ref{p3d:thm:v2216-exjet-global-obstruction} proves a strictly
positive same-global-action Ward residual for this nonzero Hessian.
Corollary~\ref{p3d:cor:v2216-raw-euler-nongluing} independently excludes
the raw Euler identification. The complete finite logarithmic germs have
third derivative $-3\chi_7/R$, whereas the integral coefficient has
$\chi_7'>0$, by Proposition~\ref{p3d:prop:v2216-finite-born-limits}.
Consequently even affine corrections cannot establish complete-germ
identity. The complete action residual and its endpoint conditions are
exactly Theorem~\ref{p3d:thm:v2216-complete-action-jets}. This permits a
separately generated effective transport while preventing its inference
from a frozen Hessian alone.
\end{proof}
\section{Capped scalaron and boundary-source action constructors}
\label{paper3e:sec:organized-cap}
\subsection{Scalar Threshold and Terminal Resolution}
\label{p3:sec:scalaron}

\subsubsection{The scalar degree of freedom}

In metric \(\mathfrak f(\scal)\) gravity the nonlinear curvature sector carries
an additional scalar mode through
\[
 \phi=\sqrt{\frac32}\ln\!\bigl(\mathfrak f_{\scal}(\scal)\bigr).
\]
For the minimal response family the required result is its
terminal decoupling.

\begin{theorem}[General metric-\(\mathfrak f\) scalaron decoupling criterion]
\label{p3e:thm:general-scalaron-decoupling}
Let \(\mathfrak f\in C^2((C,U])\cap C^1([C,U])\) be evaluated on a background
where the standard metric-\(\mathfrak f(\scal)\) scalaron formula applies.
If
\[
 \mathfrak f''(\scal)>0,\qquad
 \mathfrak f''(\scal)\longrightarrow0^+,\qquad
 \mathfrak f'(\scal)-\scal\mathfrak f''(\scal)
 \longrightarrow L>0
 \quad(\scal\downarrow C),
\]
then
\[
 m_s^2(\scal)=
 \frac{\mathfrak f'(\scal)-\scal\mathfrak f''(\scal)}
 {3\mathfrak f''(\scal)}
 \longrightarrow+\infty.
\]
This is a local Hessian criterion.  It neither requires the source endpoint
\((F_C,Z_C)=(0,1)\) nor proves the vacuum condition
\(F_C=CZ_C/2\).
\end{theorem}

\begin{proof}
The numerator tends to the strictly positive number \(L\), while the
denominator tends to zero through positive values.
\end{proof}

\begin{theorem}[Scalaron decoupling for the complete minimal response family]
\label{p3:thm:scalar-threshold}
For any response of
Theorem~\ref{p3e:thm:complete-minimal-smooth-response-family} with
\(Z_C>0\),
\[
 \mathfrak f_{\scal\scal}(\scal)
 =w_{\rm nl}(\scal)\,\frac{\gamma_\sigma}{U}\,u(\scal)>0
 \qquad(C<\scal\le U),
\]
while
\[
 \mathfrak f_{\scal\scal}(\scal)\to0^+,
 \qquad
 \mathfrak f_{\scal}(C)=Z_C>0
 \qquad(\scal\downarrow C).
\]
Hence the metric-\(\mathfrak f(\scal)\) scalaron mass parameter
\begin{equation}
\label{p3:eq:spectral-scalaron-mass}
 m_s^2(\scal)
 =\frac{\mathfrak f_{\scal}(\scal)
        -\scal\mathfrak f_{\scal\scal}(\scal)}
       {3\mathfrak f_{\scal\scal}(\scal)}
\end{equation}
defines an algebraic mass function which diverges to \(+\infty\).  Paper~III-F,
Theorem~\ref{p3f:thm:trace-linearized-scalaron-mass}, derives this expression
from the trace equation on each admissible constant-curvature vacuum
background; see also \cite{DeFelice2010}.  Therefore it is a physical
scalaron-decoupling limit only along an independently established family of
such backgrounds approaching \(C\).  Under
\((\mathrm{CAL}_{\rm resp/EH})\), \((\mathrm{CRP4})\), and
\(\mathrm{REG}_{\sigma}\in\mathfrak A_{\rm REG}^{+}\), the terminal response coordinate is
\[
 \boxed{C_{\rm fr}=C=
 \Lambda_{\rm EH}^2\,
 \alpha_{\rm hier}(\sigma_{\rm fin})^{-7}.}
\]
With the legacy alias \(\Rp:=U\), this may be written
\(C=\Rp\alpha_{\rm hier}(\sigma_{\rm fin})^{-7}\).
No geometric-radius identity for this root is asserted.  The physical
decoupling interpretation is conditional on independently admissible
backgrounds and does not prove that \(C\) is a vacuum constant-curvature
solution.  Neither \(F_C\) nor its
source value is used in the decoupling proof.
\end{theorem}

\begin{proof}
Theorem~\ref{p3e:thm:complete-minimal-smooth-response-family} gives positivity of
the Hessian on the open branch, its terminal limit zero, and
\(\mathfrak f_{\scal}(C)=Z_C>0\).  Thus the numerator in
\eqref{p3:eq:spectral-scalaron-mass} tends to \(Z_C\), and
Theorem~\ref{p3e:thm:general-scalaron-decoupling} gives
\(m_s^2\to+\infty\) as an algebraic limit.  The downstream
trace-linearization theorem identifies it with a physical scalaron mass only
on its stated background domain.  The structural factor
\(\Lambda_{\rm EH}^2\) is the generated III-E upper coefficient; the neutral ratio
\(\alpha_{\rm hier}(\sigma_{\rm fin})^{-7}\) is generated, and the physical
identity \(C/U=\alpha_{\rm hier}(\sigma_{\rm fin})^{-7}\) is its unique
inheritance through the certified CRP4 response ancestry.
\end{proof}

\begin{corollary}[The source endpoint is not a terminal vacuum]
\label{p3e:cor:source-endpoint-not-vacuum}
For the source specialization \((F_C,Z_C)=(0,1)\) and \(C>0\),
\[
 C\,\mathfrak f_{\scal}(C)-2\mathfrak f(C)=C\ne0.
\]
Hence the source member decouples the scalaron but does not supply its own
vacuum constant-curvature background.  A matter-supported or otherwise
independently supplied background is required for a physical terminal limit.
\end{corollary}

\begin{proof}
Substitution of \(F_C=0\) and \(Z_C=1\) gives the displayed residual.
\end{proof}

\begin{corollary}[Owner-typed Paper-III terminal handoff]
\label{p3:cor:scalar-threshold-handoff}
On the calibrated, owner-typed \((\mathrm{MRESP}_{\min})\) branch with
\(\mathrm{REG}_{\sigma}\in\mathfrak A_{\rm REG}^{+}\), Paper~III-E exports
\[
 \boxed{\bigl(\mathfrak f_{\sigma,\gamma,F,Z},\
 C,\ U,\ H_{56}^{(\sigma)}\bigr)},
 \qquad
 U=\Lambda_{\rm EH}^2,\quad
 C=\Lambda_{\rm EH}^2
 \alpha_{\rm hier}(\sigma_{\rm fin})^{-7},\quad
 H_{56}^{(\sigma)}=\alpha_{\rm hier}(\sigma_{\rm fin})^{7},
\]
together with convexity on \((C,U]\), terminal Hessian decay, and scalaron
decoupling at \(C\) whenever \(Z_C>0\) and the stated scalaron background
domain holds.  The source response \(\mathfrak f_{\rm sp}\) is the
\((F_C,Z_C)=(0,1)\) specialization.  With the legacy alias \(\Rp:=U\),
replace \(\Lambda_{\rm EH}^2\) by \(\Rp\) only as a notation change.
The smooth branch exports six finite ordered scale crossings
\(\scal_j=U\alpha_{\rm hier}(\sigma_{\rm fin})^{-j}\),
\(j=1,\ldots,6\), and the limiting endpoint
\(C=U\alpha_{\rm hier}(\sigma_{\rm fin})^{-7}\).  If
\((\mathrm{FLAG})\) and its
required realization data are separately supplied, this seven-entry ledger
may additionally index an ordered rank-one refinement.
\end{corollary}

\begin{remark}[Scalar hierarchy and optional projector interpretation]
The six finite integer crossings and the seventh asymptotic endpoint of the
protected rank-trace coordinate are the generated scalar scale hierarchy.
They do not name cosmological processes and are not projector events; the
latter interpretation requires \((\mathrm{FLAG})\).  The corresponding mathematical figure is in
Paper~III-G, Appendix~\ref{p3e:app:hierarchy-bookkeeping}.
\end{remark}

\subsection{A capped action from specified boundary and source ports}
\label{p3e:sec:v223-cap-constructor}
This section supplies a candidate action branch, with exact conditional
reduction.  It is not a replacement of the native heat trace by an
arbitrary scalar profile.  The native constructor of
Corollary~\ref{p3e:cor:native-heat-action-constructor} and this boundary
constructor have distinct typed inputs until an actual operator and
source comparison identifies them.  Prior branch formulae remain
statements under their own hypotheses; they do not veto a replacement
whose full native origin is subsequently proved.

\begin{definition}[Normalized cap data and ports]
\label{p3e:def:v223-cap-data}
Fix $P,R_s>0$, $B_I>B_U>0$, $\Delta=B_I-B_U$, and
$\sigma=B_IR_s$.  The dimensional types are
\[
 [P]=[R]=[R_s]=[\sigma]=\text{curvature},\qquad
 [B_I]=[B_U]=1,\qquad [f]=\text{curvature}^2.
\]
Let
\[
 x=R/R_s,\quad G=\sqrt{1+x^2},\quad t=\frac{x}{1+G},\quad
 c=\sqrt{1-B_U/B_I},\quad b=1-c^2t^2,
\]
\[
 B=B_Ib=B_U+\frac{2\Delta}{1+G},\qquad
 M_\pm(t)=\begin{pmatrix}1&\pm ct\\\pm ct&1\end{pmatrix},
 \quad T_s=M_+(t)\otimes M_-(t).
\]
The reference response fibre is $E_0=\mathbb R^4$ with metric $I$;
the output fibre has metric $T_s$.  In the dimensionful source chart,
$Y:E_0\to E_R$ is a real $4$ by $4$ matrix with entries of curvature
type, and $j_s=\rho I_4$ has the same entry type, $[\rho]=[R_s]$.
Equivalently one may use dimensionless coordinates $y=Y/R_s$, but its
measure must retain the Jacobian $d^{16}Y=R_s^{16}d^{16}y$.
The candidate incoming section is $j_s$ in the specified reference
identification.
Neither response fibre is identified by its dimension with spacetime,
the tangent fibre of $M_8$, or the native momentum carrier.
\end{definition}

\begin{theorem}[Normalized source response and reciprocal transport]
\label{p3e:thm:v223-cap-source}
The matrices in Definition~\ref{p3e:def:v223-cap-data} are positive,
$\det T_s=b^4$, and $\operatorname{Tr}T_s^{-1}=4/b^2$.
Stationary elimination of
\[
 W_s(Y)=\operatorname{Tr}(Y^{\mathsf T}T_sY)
                       +2\operatorname{Tr}(j_s^{\mathsf T}Y)
\]
gives $Y_*=-T_s^{-1}j_s$, $W_{s,*}=-4\rho^2/b^2$.
The reference source density $W_{s,*}(0)/8=-R_s^2/8$ fixes
$\rho=R_s/2$ up to its irrelevant overall sign, and hence
\[
 \frac{W_{s,*}}8=-\frac{\sigma^2}{8B^2}.
\]
For a general source $j$, the response of $W_{s,*}/8$ is
$-T_s^{-1}j/4$ and its source Hessian is $-(I_4\otimes T_s^{-1})/4$.

For $x>0$ let $\widehat t=t(1/x)$ and
\[
 T_r=M_+(t)\otimes M_-(\widehat t),\quad
 P_-=M_-(t),\quad Q_-=M_-(\widehat t),\quad
 A=I_2\otimes(Q_-^{-1/2}P_-^{1/2}).
\]
Then $A^{\mathsf T}T_rA=T_s$.  With
$Y_r=AY_s$, $j_r=A^{-\mathsf T}j_s$, the complete source words,
stationary responses, mixed curvature/source jets and induced Gaussian
measures agree.  In particular
\[
 \operatorname{Tr}(j_r^{\mathsf T}T_r^{-1}j_r)=4\rho^2/b^2,
 \qquad [b(x)b(1/x)]^4|\det A|^4=b(x)^8.
\]
\end{theorem}
\begin{proof}
The port eigenvalues are $1\pm ct>0$, their determinants are $b$,
and each inverse has trace $2/b$.  Tensor multiplication proves the
trace and determinant claims.  Completing the square proves the source
equations and differentiation gives the response.  Since $P_-,Q_-$
are functions of the same symmetric involution, they commute and
$C=Q_-^{-1/2}P_-^{1/2}$ is their positive metric comparison.
The remaining claims follow from
Theorem~\ref{p1b:thm:v223-positive-source-transport}, including its
column Jacobian and the transported linear term.
\end{proof}
Constant coordinate sections $\rho I$ in $T_s$ and $T_r$ are different
physical prescriptions.  The latter has bare return
$4\rho^2/[b(x)b(1/x)]$.  Its logarithmic derivative vanishes at
$x=1$, by reciprocal symmetry, and therefore at $q=7/2$ \emph{only
if} the additional midpoint identification $R_s^2=UC$ is imposed,
with $q(R)=\log(U/R)/\log\alpha$ and $\alpha^7=U/C$.
On this same specialized branch the cap return has
\[
 \left.\partial_q\log(B^{-2})\right|_{7/2}
 =-2L(1-1/\sqrt2)\frac{B(R_s)-B_U}{B(R_s)}\ne0.
\]
For free $R_s$, the reciprocal stationary point occurs at $R=R_s$,
not necessarily at hierarchy depth $7/2$; neither derivative comparison
at $7/2$ is asserted without the stated midpoint guard.
Thus a fixed scalar normalization cannot repair that bare response on
the specialized branch.  Transporting the section repairs the complete
comparison: at $x=1$ the transport has $A=I$ but its derivative in a
nonconstant curvature/depth chart need not vanish, so differentiating
$j_r$ supplies the missing contact.  An isometry transports a supplied section; it does
not prove that the native incoming section is the one just specified.

\begin{proposition}[Running and frozen source returns are distinct marked parents]
\label{p3e:prop:v3028-marked-source-trichotomy}
On the normalized four-port comparison of
Definition~\ref{p3e:def:v223-cap-data}, write
\[
\begin{aligned}
 u&=c\,t(R/R_s),& b&=1-u^2=\frac{B(R)}{B_I},\\
 \rho&=R_s/2,& X&=\begin{pmatrix}0&1\\1&0\end{pmatrix},\\
 T_s&=(I_2+uX)\otimes(I_2-uX).
\end{aligned}
\]
In a \emph{common, specified reference chart} consider the three
positive-port source sections
\[
 j_{\rm fix}=\rho I_4,\qquad
 j_{\rm det}=\rho bI_4,\qquad
 j_{\rm iso}=\rho T_s^{1/2},
\]
where $T_s^{1/2}$ denotes the unique positive square root.  For
\[
 W(Y;j)=\operatorname{Tr}(Y^{\mathsf T}T_sY)
       +2\operatorname{Tr}(j^{\mathsf T}Y),\qquad
 V[j](R)=\min_Y W(Y;j),
\]
the exact stationary \emph{scalar $f$-density} contributions are
\[
 \boxed{V[j_{\rm fix}]=-\frac{R_s^2}{b^2},\qquad
        V[j_{\rm det}]=V[j_{\rm iso}]=-R_s^2.}
\]
The corresponding action-density contribution is $V[j]/8$, not
$V[j]$ itself, in the normalization of
Theorem~\ref{p3e:thm:v223-exact-cap-reduction}.  For $u\ne0$,
$j_{\rm det}$ and $j_{\rm iso}$ are \emph{distinct marked currents}:
\[
 j_{\rm det}^{\mathsf T}T_s^{-1}j_{\rm det}
     =\rho^2b^2T_s^{-1},\qquad
 j_{\rm iso}^{\mathsf T}T_s^{-1}j_{\rm iso}=\rho^2 I_4.
\]
The positive comparison isometry of
Theorem~\ref{p1b:thm:v223-positive-source-transport}
does not select the incoming physical section.  The induced
four-column volume $b^8d^{16}Y$ does not require a determinant
rescaling of that source.  These are conditional alternative
\emph{source constructors}, not coordinate descriptions of the
fixed-reference cap prescribed by
Definition~\ref{p3e:def:v223-cap-data}.
\end{proposition}
\begin{proof}
The four eigenvalues of $T_s$ are $(1+u)^2$, $(1-u)^2$,
$b$, and $b$.  Since $|u|<1$, the form is positive and
\[
 \det T_s=b^4,\qquad \operatorname{Tr}T_s^{-1}=\frac4{b^2}.
\]
Completing the square yields
$Y_*=-T_s^{-1}j$ and
$V[j]=-\operatorname{Tr}(j^{\mathsf T}T_s^{-1}j)$.
Substitute the three currents, using
$T_s^{1/2}T_s^{-1}T_s^{1/2}=I_4$ and $4\rho^2=R_s^2$.
For $u\ne0$, $T_s$ is not the scalar $b^2I_4$, proving the
marked-Gram distinction. A genuine coordinate transport
$Y_r=AY_s$ transforms the kernel, source and measure jointly as
$T_s=A^{\mathsf T}T_rA$ and $j_r=A^{-\mathsf T}j_s$.
Multiplying $j_s$ by $b$ while fixing $T_s$ is instead a change of
the physical source word.  The statement about the induced volume
follows from Proposition~\ref{p3e:prop:v223-source-volume}.
\end{proof}

\begin{corollary}[Reference curvature jets distinguish running from frozen sources]
\label{p3e:cor:v3028-marked-source-jets}
With $d=c^2=1-B_U/B_I$, keep the same $PR+B(R)R^2$ in each
comparison scalar action, and set
\[
 f_{\rm run}=PR+B(R)R^2-R_s^2b^{-2},\qquad
 f_{\rm frozen}=PR+B(R)R^2-R_s^2.
\]
Then
\[
 f_{\rm frozen}-f_{\rm run}=R_s^2(b^{-2}-1)\ge0,
 \qquad b(0)=1,\quad b'(0)=0,\quad
 b''(0)=-\frac{d}{2R_s^2},
\]
\[
 \boxed{f_{\rm run}''(0)=2B_I-d,\qquad
        f_{\rm frozen}''(0)=2B_I.}
\]
The current jets in the common reference chart are
\[
\begin{aligned}
 j_{\rm fix}'(0)&=j_{\rm fix}''(0)=0,\\
 j_{\rm det}'(0)&=0, &
 j_{\rm det}''(0)&=-\frac{\rho d}{2R_s^2}I_4,\\
 j_{\rm iso}'(0)&=\frac{\rho c}{4R_s}
       (X\otimes I_2-I_2\otimes X)\ne0.
\end{aligned}
\]
Equality of eliminated scalar $f(R)$ terms cannot distinguish the
two frozen marked-current parents.  Their native full-source
comparison requires an identified physical source and adequate
source probes, as in
Theorem~\ref{p3f:thm:v289-complete-source-twojet}.
\end{corollary}
\begin{proof}
$t(x)=x/2+O(x^3)$ gives
$b(R)=1-dR^2/(4R_s^2)+O(R^4)$.  Differentiate the
corresponding scalar actions.  Finally
$u'(0)=c/(2R_s)$ and
$\bigl(T_s^{1/2}\bigr)'(0)=T_s'(0)/2$ give the
marked-current jets.
\end{proof}


\begin{proposition}[Finite source volume and multiplicity]
\label{p3e:prop:v223-source-volume}
For the curvature-valued $Y$ chart of
Definition~\ref{p3e:def:v223-cap-data}, let $\hbar_{\rm eff}>0$ have
$[\hbar_{\rm eff}]=\text{curvature}^2$, so all three terms in
$W_s(Y)$ and its exponential denominator have the same dimension.
The matrix-entry metric is $I_4\otimes T_s$ and its volume density is
$b^8d^{16}Y$.  Exactly,
\[
 \int b^8d^{16}Y\,
 e^{-W_s(Y)/(8\hbar_{\rm eff})}
 =(8\pi\hbar_{\rm eff})^8
 e^{\rho^2\operatorname{Tr}T_s^{-1}/(8\hbar_{\rm eff})}.
\]
Both sides have curvature dimension $16$.  On the dimensionless chart
$y=Y/R_s$, the left measure is $b^8R_s^{16}d^{16}y$,
$W_s(R_sy)/R_s^2$ is dimensionless, and
$\bar\hbar_{\rm eff}=\hbar_{\rm eff}/R_s^2$ is dimensionless;
the rescaled integral keeps the full $R_s^{16}$ Jacobian.
A different measure $m(R)d^{16}Y$ with dimensionless
$m(R)/m(0)$ gives the relative effective term
$\hbar_{\rm eff}[8\log b-\log(m(R)/m(0))]$.
For an actual spectator of dimension $n$, $T_n=T_s\otimes I_n$ has
$\operatorname{Tr}T_n^{-1}=4n/b^2$; the same baseline requires
$\rho_n=R_s/(2\sqrt n)$.  On the full matrix space of dimension
$(4n)^2$ the corresponding volume exponent is $8n^2$.
\end{proposition}
\begin{proof}
$\det(I_4\otimes T_s)=b^{16}$.  Translate the completed square and
use the finite positive Gaussian formula; its determinant cancels the
volume.  The spectator assertions follow from tensor trace and
determinant identities with all matrix columns retained.
\end{proof}
The native measure constructor must justify this volume if the source
is quantum integrated.  Covariance does not forbid an additional
positive background weight.  No source reset, vacuum subtraction or
division by a carrier rank is authorized by these equalities.  For
example, independently tensoring two full eight-dimensional lifted
ports $M_\pm\otimes I_4$ retains a spectator of dimension $16$.
The native protected octet, the $15$-dimensional momentum carrier and
this $4$-dimensional response tensor have distinct owner maps.

\begin{theorem}[Exact capped action on a stated constructor domain]
\label{p3e:thm:v223-exact-cap-reduction}
Assume the actual compared constructor has the following specified
data on a smooth four-dimensional Lorentzian patch:
\begin{enumerate}
\item Its scalar curvature channel is $R[g]$, with reference $R_s$,
and its gravitational response is the decaying square-root extension
of $(J+xK)^2$ with the conjugate-flux port
$a+b=2\sqrt\Delta R$ of
Theorem~\ref{p1b:thm:v223-extension-flux}, plus $B_UR^2$ and $PR$.
\item Its source word is the one in
Theorem~\ref{p3e:thm:v223-cap-source}, or its complete transported
presentation, with the stated reference section, multiplicity and
baseline.  Its scalar density weight is $1/8$.
\item Its remaining action has no additional dependence on the
eliminated response variables on this branch; overlaps are removed.
There are no unevaluated derivative or tensor contributions in the
claimed scalar sector.  The canonical smooth non-null endpoint fixes
$h_{ij}$ and the total curvature momentum $\Phi$.
\end{enumerate}
Then classical elimination, with the full coefficient contacts, gives
\[
 \boxed{f_c(R)=PR+R^2B(R)-\sigma^2/B(R)^2,}
\]
\[
 \boxed{\frac{S_c}{\hbar}=
 \frac18\int_M\sqrt{-g}\,f_c(R)\,d^4x
 +\frac14\sum_\partial\epsilon\int_\partial
      \sqrt{|h|}\,f_c'(R)K\,d^3x+\mathcal I_{\rm rest}.}
\]
Here $\mathcal I_{\rm rest}=S_{\rm rest}/\hbar$.
If these finite source variables are also integrated with the stated
geometric volume, their Gaussian reduction adds only a
background-independent normalization to this finite comparison.
\end{theorem}
\begin{proof}
The extension and flux minimum give $\Delta R^2\,2/(1+G)$.
Adding $B_UR^2$ gives $BR^2$.  The normalized source theorem gives
$-\sigma^2/B^2$.  The first-order parent and boundary calculation in
Theorem~\ref{p3f:thm:v223-cap-parent} give the stated complete action,
including the source part of $f_c'$.  The finite Gaussian claim is
Proposition~\ref{p3e:prop:v223-source-volume}.
\end{proof}
The scalar $1/8$ and the source baseline are comparison inputs here.
Paper~III-B's protected constituent provides an antecedent for the
former, but its actual density map into this constructor must still
be identified.  A tangent dimension count, a raw represented trace
factor $96$, or $2I$ representation dimensions do not replace it.
The same applies to the dimensional comparison $P$ versus the native
curvature normalization $U$.  The boundary extension is a
representation of this gravitational response, not a second additive
contribution to an already counted Schur response.

\begin{proposition}[Common-interior incidence distinguishes the cap]
\label{p3e:prop:v223-port-incidence}
Put $u=ct$ and
\[
 \mathcal G=\frac{B_I}{2}\sum_\pm
 (r,z_\pm)M_\pm(t)\binom r{z_\pm}.
\]
Independent stationary elimination gives $Br^2$.  A shared variable
$z_+=z_-=z$ instead gives $B_Ir^2$.  Antisymmetric incidence
$z_+=z$, $z_-=-z$ gives
$\mathcal G=B_I(r^2+2urz+z^2)$ and again $Br^2$.
\end{proposition}
\begin{proof}
Independent elimination gives $z_\pm=\mp ur$.  With a shared symmetric
variable the cross terms cancel and $z=0$; with antisymmetric
incidence completing the square gives $z=-ur$.
\end{proof}
The shared symmetric incidence differs by $B_Iu^2r^2$ and its retained
Hessian by $2B_Iu^2$, the common-variable variance of Paper~III-A.
Co-consumption alone therefore does not select the response: the actual
incidence must be evaluated before separate eliminations are combined.
These boundary and Schur forms are alternative representations of the
same candidate coefficient and must not be added twice.

\begin{remark}[Remaining native identifications are concrete maps]
\label{p3e:rem:v223-native-selection}
The exact theorem above constructs a candidate; unconditional native
selection remains unproved.  It requires the native normal operator
and boundary domain to be the displayed extension, the inherited
curvature-to-flux and scalar-density maps to have its absolute weights,
the actual common-action source to use the displayed comparison and
incoming section with all multiplicities, and the remaining action and
physical measure to reduce without omitted derivative, tensor or
boundary terms.  Shared exponential laws, an available metric isometry
and matching endpoint values do not establish those identities.
The native fixed momentum radius and its carrier are not changed by
this candidate.  No additional ultimate principle or physical field is
postulated to fill these identifications.  A full native evaluation
could select this branch, select a corrected branch, or reject it.
\end{remark}

\begin{remark}[A complete phase-source Gram is not yet the cap source]
\label{p3e:rem:v244-rooted-cap-overlap}
The exact source in Corollary~\ref{p1:cor:v244-quarter-turn-full-gram}
retains a traceless but nonzero signed family cross term.
Proposition~\ref{p3f:prop:v244-rooted-phase-schur-order} shows that
eliminating its family plane before the scalar readout changes the
second source jet and produces higher-order Schur memory, whereas
averaging the linear Gram first removes that term.  Neither operation
is identified here with the actual physical metric/source reduction.
Thus the rooted quarter-turn and its trace cancellation do not
by themselves establish the constant inverse-square cap source,
its absolute \(h^4\) norm, the rank-seven normal defect,
Newton normalization, or the whole-action vacuum allocation.
The source-order and overlap tests of
Proposition~\ref{p3e:prop:v224-source-moments} and
Remark~\ref{p3e:rem:v223-native-selection} remain required; a
shared boundary/Schur contribution must not be counted twice.
\end{remark}


\subsection{Source tests and reference-boundary realization}
\label{p3e:sec:v224-source-matching}
The exact cap reduction of Theorem~\ref{p3e:thm:v223-exact-cap-reduction}
has a second, source-complete realization. Theorem~\ref{p3f:thm:v224-reference-cap-parent}
derives its curvature-squared source from a reference Poisson defect,
and Proposition~\ref{p3f:prop:v224-collar-finite-transport} identifies
the whole form with the existing finite parent before elimination.
These strengthen the constructor without claiming its native selection.

\begin{proposition}[Three-moment source test]
\label{p3e:prop:v224-source-moments}
For the cap kernel set $\mathsf A=X\otimes I$,
$\mathsf B=I\otimes X$, with
$X=\left(\begin{smallmatrix}0&1\\1&0\end{smallmatrix}\right)$.
Let $W=jj^\dagger\succeq0$ be constant in this compared frame. Then
\[
 \operatorname{Tr}(j^\dagger T(u)^{-1}j)
 =\frac{a_0+ua_1-u^2a_2}{(1-u^2)^2},\qquad
 (a_0,a_1,a_2)=
 (\operatorname{Tr}W,\operatorname{Tr}[W(\mathsf B-\mathsf A)],
                              \operatorname{Tr}(W\mathsf A\mathsf B)).
\]
On an open interval, this scalar return equals
$mR_s^2/(1-u^2)^2$ if and only if
\[
 \operatorname{Tr}W=mR_s^2,\qquad
 \operatorname{Tr}[W(\mathsf B-\mathsf A)]=0,\qquad
 \operatorname{Tr}(W\mathsf A\mathsf B)=0.
\]
\end{proposition}
\begin{proof}
The commuting involutions give
$T^{-1}=[I+u(\mathsf B-\mathsf A)-u^2\mathsf A\mathsf B]/(1-u^2)^2$.
Take the source trace; polynomial identity on an open interval equates
its three coefficients.
\end{proof}
The test does not require $W$ to be proportional to identity. Its scalar
return may leave information visible to other interactions. If the source
or frame moves, the general condition is only
$a_0(u)+ua_1(u)-u^2a_2(u)=mR_s^2$, with all its actual derivatives.
The native kernel and the eliminated-source centre from
Proposition~\ref{p1b:prop:v224-canonical-heat-source} remain part of the
same matching problem. The target weight $m=96$ is not yet an evaluated
native source norm or a count of independent species.

\begin{proposition}[Common Einstein/source multiplicity]
\label{p3e:prop:v224-common-multiplicity}
Let $P_0>0$, $R_s^2=P_0C_7$, $\sigma_0=B_IR_s$, and distinguish
the actual Einstein multiplicity $M_E>0$ from source weight $m$.
For the specified complete comparison
\[
 f_{M_E,m}(r)=M_E[P_0r+B(r)r^2]
                         -m\frac{\sigma_0^2}{B(r)^2},
\]
the small-curvature normalization is
\[
 \frac{f_{M_E,m}(r)}{M_EP_0}
             =r-\frac m{M_E}C_7+O(r^2).
\]
In particular equal actual multiplicities preserve the normalized
terminal intercept and the corresponding scaled boundary derivatives.
\end{proposition}
\begin{proof}
At zero curvature $B=B_I$ and $B'=0$. Thus the constant is
$-mR_s^2$ and the linear coefficient is $M_EP_0$. Divide by
the latter; differentiating the complete identity transports its endpoint.
\end{proof}
This is an alternative normalization comparison. It does not simultaneously
identify the raw $P=96U$ and the retained $P=U$ descriptions without
their complete readout map. The marked coefficient result and its gauge,
scalar and source transport are in Proposition~\ref{p3f:prop:v224-normalization-transport}.

\begin{corollary}[Exact source law with remaining physical inputs]
\label{p3e:cor:v224-reference-source-status}
For the supplied normal square $\mathcal A_r=\ell^2[1+(r/R_s)^2]I$
and fixed boundary norm $b_0^2=4\Delta R_s^2/\ell$,
\[
 j_r=(\mathcal A_r-\mathcal A_0)Q_0,\qquad
 B_Ur^2+\min_Q\mathscr E_r-C_0=r^2B(r),\qquad
 C_0=2\Delta R_s^2.
\]
These are whole-action identities on the stated positive Dirichlet
constructor domain, including its reference centre and source derivatives.
\end{corollary}
\begin{proof}
Apply Theorem~\ref{p1b:thm:v224-poisson-defect} and
Theorem~\ref{p3f:thm:v224-reference-cap-parent}.
\end{proof}
The same construction consumes $u_{\rm op}(r)$ as a geometric coefficient;
an independently varied off-shell flux requires a fresh coupled reduction.
The remaining native comparison must fix the normal band, scale and
boundary domain; the actual nonzero datum and susceptibility; complete
curvature, source and normalization maps; tensor, derivative, source-orthogonal
and boundary companions; physical measure and absolute energy allocation.
The route tests in Section~\ref{p3f:subsec:v224-route-tests} reject the
specified finite attenuation, local scalar and intrinsic phase-lift routes.
They do not exclude a different licensed common-action construction.
Neither the reference energy nor the retained normalization cancels the
bare heat vacuum. No prior endpoint-vacuum result is transferred to the
new cap without its own complete action identity.


\subsection{Exact profile-orbit invariance}
\label{p3e:sec:v164-profile-orbit}

\begin{theorem}[Generated gravitational heat orbit]
\label{p3e:thm:v164-generated-gravitational-heat-orbit}
On the premise-minimal scalar heat branch of
Theorem~\ref{p3b:thm:v164-generated-heat-profile}, the profile and cutoff
occur only through
\[
 \Lambda_{\rm eff}^2=\frac{\Lambda^2}{a},
 \qquad F_a(u)=e^{-au}.
\]
The four-dimensional gravitational second moment is
\[
 f_{2,g}(a)=\frac1{(4\pi)^2}\int_0^\infty e^{-au}\,du
 =\frac1{(4\pi)^2a},
\]
so the Einstein--Hilbert coefficient satisfies
\[
 \boxed{
 \Lambda_{\rm EH}^2
 =\frac{f_{2,g}(a)}6\Lambda^2
 =\frac{\Lambda_{\rm eff}^2}{6(4\pi)^2}.}
\]
Every same-action coefficient and response obtained by variation is constant
on \((a,\Lambda)\mapsto(\rho a,\sqrt\rho\Lambda)\).  This proves profile
invariance; it does not minimize \(a\) or choose an SI unit.
\end{theorem}

\begin{corollary}[Newton normalization fixes the orbit]
\label{p3e:cor:v164-newton-fixes-profile-orbit}
On the primitive momentum-curvature realization,
\[
 \boxed{
 \Lambda_{\rm eff}^2=96\pi^2\Lambda_{\rm EH}^2,
 \qquad
 \Lambda_{\rm eff}^2=96\pi^2R_P.}
\]
This fixes the effective-scale orbit, not a preferred representative
\((a,\Lambda)\) within it.
\end{corollary}


\begin{corollary}[Derived-later local Yang--Mills action and phase space]
\label{p3e:cor:v174-yang-mills-action-derived-later}
On the native stationary-product, parity-even, pure-gauge domain of
Paper~III-F,
Theorem~\ref{p3f:thm:v174-generated-joint-gauge-phase-space-action}, the
joint curvature role, its nonlinear local Lorentzian Yang--Mills action, the
regular connection-field cotangent role, Gauss moment map, positive reduced
Hamiltonian, and constrained Legendre equivalence are one generated
co-realized package.  Denote its certificates by
\[
 (\mathrm{ACT}_{{\rm YM},U}^{\rm der})
 \equiv(\mathrm{ACT}_{\rm YM,prod}^{\rm der})
 +(\mathrm{JPHASE}_{\rm YM}^{\rm der})
 +(\mathrm{GAUSS}_{\rm YM}^{\rm der}).
\]
No separate pure-gauge \((\mathrm{LOR}_{\rm dir})\), gauge-phase-space, or
source-free Gauss-law premise remains on the smooth theorem domain.  On the
additional holomorphic first-jet subdomain, the same action also discharges the
pure-gauge restriction of \((\mathrm{CONT}_{\rm cl})\).  This
corollary is sector-local; the Einstein and native odd-scalar action branches
are discharged separately by the joint-Dirac results below.  Coupled
fermionic current, Lorentzian fermionic reality, exact-nonlocal and
curvature-squared continuation, offset/Majorana terms, noncanonical/global
boundaries, finite-regulator and continuum actions, vacuum selection, and
actual solution/history interfaces remain separately typed.
\end{corollary}

\begin{corollary}[Derived-later joint-Dirac Einstein and native odd-scalar actions]
\label{p3e:cor:v175-einstein-scalar-actions-derived-later}
On the clock-adapted local joint-Dirac domain of Paper~III-F,
Theorem~\ref{p3f:thm:v175-generated-joint-einstein-action} derives the
Einstein--Hilbert interior action germ and its positive coefficient ray, while
Theorem~\ref{p3f:thm:v175-native-local-scalar-action} derives the native
parity-even odd-scalar kinetic, quartic, curvature-coupling, and quadratic
terms through canonical dimension four.  Together with
Corollary~\ref{p3e:cor:v174-yang-mills-action-derived-later}, their certificate
is
\[
 \boxed{(\mathrm{ACT}_{{\rm J,EH+YM+\Phi},U_0}^{\rm loc,der}).}
\]
The positive \(+R/6\) coefficient of the unit-preserving scalar Laplacian in
Proposition~\ref{p3:prop:EH_from_a2} remains a scalar-resolution/readout
coefficient.  It is not the action coefficient of the joint Dirac owner,
whose Lichnerowicz endomorphism supplies the negative Euclidean Einstein
coefficient required by the common clock real form.  The local action-role bridge is therefore derived. The later common-heat
carrier evaluation fixes its raw coefficient as
\(\kappa_J=12\Lambda_{\rm EH}^2\), equivalently
\(U_J^{\rm raw}=96U\).  The raw joint coupling is $G_{J,\rm raw}=G_{\rm scal}/96$. A final
physical Newton coefficient uses the complete matched action and source
normalization; the local action-role theorem does not itself normalize that
projection or the separate momentum/EH port. Empirical comparison follows
those evaluated maps.

The corollary does not continue
pure-metric vacuum or curvature-squared terms, exact nonlocal heat channels,
fermionic adjoint/reality data, constant odd/Majorana offsets, or
noncanonical/global boundaries, and it does not prove a broken vacuum.
\end{corollary}


\begin{corollary}[Derived-later local bosonic Noether and constraint packet]
\label{p3e:cor:v176-local-bosonic-noether-constraint-derived-later}
On the same restricted common local action domain as
Corollary~\ref{p3e:cor:v175-einstein-scalar-actions-derived-later},
Paper~III-F now additionally derives the scalar gauge current, the coupled
Yang--Mills--scalar moment map and Gauss constraint, the gauge--scalar
fibre-product constraint locus, the nonminimal gravity--scalar canonical
mixing, the local Hamiltonian/momentum/Gauss constraint packet, and the
Noether--Ward derivative tower.  In particular,
\[
 \boxed{
 \pi-\frac12\operatorname{Tr}_F(\Phi p_\Phi)
 =-2\kappa_J\sqrt h\,K,\qquad
 \Delta_{g\Phi}=\kappa_J>0,}
\]
so no additional conformal primary constraint survives on the regular local
joint-Dirac branch.  These are derived-later descendants of the already
certified common action; they introduce no new dynamics principle or action
owner.

The import is strictly local and owner-typed.  It does not transfer the
restricted broken-Weyl trace identity or ADM constraint equations to the
nonlinear \(f(R)\), curvature-response, scalaron, or RAR branches.  Such a
transfer requires a separately proved common-action ancestry and the relevant
boundary/global data.
\end{corollary}


\subsection{M8 provenance of the CRP4 fourth root}
\label{p3e:sec:v187-m8-crp4-provenance}
The CRP4 uniqueness theorem is unchanged: the response character $H$ is
produced by its stated response owner and hypotheses.  The top-down $M_8$
closure now identifies the geometric rank consumed by that theorem.  The full
relative solder has
\[
 \operatorname{rank}\mathcal E_{M8}=4,
\]
and its local generator descendant carries determinant character
$\det(aI_4)=a^4$.  Therefore the previously derived positive response root
\[
 \boxed{a=H^{1/4}}
\]
has a parent-geometric explanation: the exponent four is the rank of the
$M_8$-descended local coframe.  This statement does not derive $H$ from
$M_8$, nor does it identify any hierarchy energy with an observed mass.


\subsection{Exceptional defect residues versus hierarchy depth}
\label{p3e:sec:v188-defect-residue-hierarchy}

\begin{theorem}[Two independent depth structures]
\label{p3e:thm:v188-two-depth-structures}
The hierarchy depth
\[
 8\to4\to3\to1,
 \qquad q=0,4,5,7,
 \qquad K=8q
\]
is the rank/codimension filtration of one protected Clifford owner.  It is
not the derived series of any exceptional finite group.

Independently, restriction of the $M_8$ order-swap defect gives
\[
 2O\to2T\to Q_8\to\{\pm1\}\to1,
 \qquad
 2T\to Q_8\to\{\pm1\}\to1,
 \qquad
 2I\to2I.
\]
The quotient residues are
\[
 2T/Q_8\cong C_3,
 \qquad 2O/2T\cong C_2,
 \qquad (2I)_{\rm ab}=1.
\]
Thus the trimer clock, normalizer parity and character-free selector are
three different finite dynamical shadows of the same parent defect, whereas
$q$ and $K$ remain Clifford rank data.  No identification of these two depth
notions is licensed by their coexistence.
\end{theorem}

\begin{proof}
The first filtration is defined by the ranks of the protected Clifford
projectors and its depths are codimensions in one eight-dimensional carrier.
The second is defined functorially by repeated group commutators and is
computed from the exceptional derived-subgroup identities.  Since rank loss
and group commutator closure are different constructors with different target
sorts, PCAD retains their common ancestry but supplies no identification.
\end{proof}

\subsection{Amplitude transport, final-octet scalar process, and K16 flux synchronization}
\label{p3e:sec:v194-amplitude-flux-sync}

On the protected complement the relative one-way amplitude and quadratic
transport are distinct descendants:
\[
 \boxed{A_\tau=V_\tau^{1/2}|_W=e^{-\tau/2}I_W,}
 \qquad
 \boxed{\mathcal T_\tau(H)=A_\tau^*HA_\tau=e^{-\tau}H.}
\]
On the post-$K=48$ relative depth $\ell$ this yields the canonical scalar
process coefficient $e^{-\ell}$ and explains the factor of two between a
one-way amplitude and a quadratic/return factor.  A different exponent or an
extra scalar character requires an additional generated response owner once
same-owner ancestry is established; this statement does not itself prove a
cross-owner identification.

At the $K=16$ source-free junction the controlling condition is the generalized
flux law
\[
 \boxed{\chi_+(\nabla_nR)_+=\chi_-(\nabla_nR)_-.}
\]
Since the retained susceptibility is continuous and positive at K16, this
flux law gives $[(\nabla_nR)]_{16}=0$. A $C^1$-curvature branch has the
same derivative-continuity condition; it does not imply
$\nabla_nR=0$. A zero derivative requires a separately supplied turning,
reflection or boundary condition. No independent microscopic smoothing width
is reintroduced by this synchronization.


\subsection{Raw Laurent secant versus physical Euler-response scope}
\label{p3e:sec:v1107-raw-euler-distinction}
\begin{proposition}[The raw Laurent secant cannot be the terminal metric Euler response]
\label{p3e:prop:v1107-raw-not-physical-euler}
Put \(C=C_7>0\).  On \(C<R\le C_2\) the depth-two dressing is
\(D_2=1\) and the raw Laurent secant is
\(\mathfrak P_{2\mid7}=(R^2+UC)/(U+C)\).  The unique
terminal-normalized inverse-Euler primitive of this \emph{raw} function
has
\[
 f_{\rm raw}(R)=R-C+\frac{(R-C)^2}{U+C},\qquad
 (f_{\rm raw})_{RR}(C)=\frac{2}{U+C}>0.
\]
Consequently it cannot be the retained terminally degenerate
\(f_{2\mid7}\) metric scalaron action.  The physical Euler response is
instead
\[
 \boxed{P_{2\mid7}(R)=C+\int_C^R xD_2(x)\chi_7^{\rm jet}(x)\,dx,}
\]
obtained from Paper~III-F's supplied fixed-depth local Hessian on a successfully
evaluated $(\mathrm{EXJET}_7)$ source domain. Without that source
certificate it remains the valid candidate-response reconstruction.
Both response readouts equal \(C\) at the endpoint, but their first
curvature derivatives do not agree there.  No integration constant or
Einstein--Hilbert affine kernel can resolve that derivative mismatch.
\end{proposition}
\begin{proof}
Differentiation of \(Rf_R-f=\mathfrak P\) gives
\(Rf_{RR}=\mathfrak P'\).  On the lower branch
\(\mathfrak P'=2R/(U+C)\), hence \(f_{RR}=2/(U+C)\).
The terminal conditions \(f(C)=0,f_R(C)=1\) uniquely integrate to the
stated quadratic polynomial.  For the physical action
\(f_{RR}=D_2\chi_7^{\rm jet}\), so
\(P'_{2\mid7}=Rf_{RR}\) and \(P_{2\mid7}(C)=C\) give the other
formula, with \(f_{RR}(C)=0\).  The difference already occurs in the
second jet, which is invariant under addition of any affine term.
\end{proof}
\begin{remark}[Logical scope of the corrected handoff]
The raw \(D_2\mathcal H_7^{\rm act}\) is a finite Laurent response
\emph{readout}; the conditionally certified fixed-depth local configuration response
\(D_2\chi_7^{\rm jet}\) is a different differential sort.  The
terminal-normalized metric coefficient is the twofold integral of the
latter.  This corrects the inverse-Euler identification advertised in
older III-E wording, while preserving the K16 finite-third-jet kink,
K48 positive susceptibility, terminal affine representative and the
original raw Laurent hierarchy values.
\end{remark}

\subsection{Fixed momentum heat values and the complete metric-action port}
\label{p3e:sec:v2218-fixed-momentum-action-port}

The momentum carrier and its radius remain fixed throughout this section:
\(\mathcal H_p\cong\mathcal H_1\), \(m_p=1\), and
\(R_{\rm mom}\) is constituted representation data.  The varied
macroscopic metric belongs to the reciprocal-position polarization, as in
Paper~III-F, Theorem~\ref{p3:thm:completed-carrier-dynamical-asymmetry}.
No homothety of the microscopic momentum radius is used.  A finite heat
value, a neutral hierarchy coordinate and a local metric coefficient are
retained as distinct outputs until their actual operator, density and
source comparison is evaluated.

\subsubsection{The raw finite trace and its two direct promotions}

\begin{proposition}[Exact logarithmic raw-heat value and inverse-Euler test]
\label{p3e:prop:v2218-raw-heat-log-port}
Use the fixed finite seed and the Born-half depth of Paper~III-C,
Theorem~\ref{p3c:thm:v2218-born-heat-depth}.  On its specified
physical-operator comparison branch, write the heat trace as
\(Z_H(\theta)\), so that
\[
 A_p(\theta)=\frac12 Z_H(\theta),\qquad
 q_H(\theta)=\frac12(15-Z_H(\theta)),\qquad
 q_H:0\longrightarrow7.
\]
Without that operator comparison, \(A_p\) denotes the Born-weighted
finite comparison value rather than an identified physical action value.
Put \(L=\log\alpha>0\), \(C=C_7=U\alpha^{-7}\), and form the
neutral scalar pullback
\(R_H=U\alpha^{-q_H}\) of
Corollary~\ref{p3c:cor:v2218-born-depth-neutral-character}.
Then eliminating \(\theta\) gives exactly
\begin{equation}
 \label{p3e:eq:v2218-finite-heat-log-value}
 \boxed{A_p=\frac12+\frac{\log(R_H/C)}{L}.}
\end{equation}
This identity is independent of the intermediate heat profile.

If a proposed positive raw-value promotion sets
\(f_{\rm val}(R)=aR+b+\gamma A_p(R)\), \(\gamma>0\), then
\begin{equation}
 \label{p3e:eq:v2218-raw-value-negative-hessian}
 \boxed{f_{\rm val}''(R)=-\frac{\gamma}{LR^2}<0.}
\end{equation}
The affine terms do not repair that sign.  If instead a proposed
positive affine raw-trace \emph{Euler} port is
\(P(R)=C+\lambda\log(R/C)\), \(\lambda>0\), its unique
terminal-normalized inverse-Euler coefficient is
\begin{equation}
 \label{p3e:eq:v2218-log-euler-coefficient}
 \boxed{f_{\log}(R)=R-C+\lambda
 \left[\frac RC-1-\log\frac RC\right],\qquad
 f_{\log}''(R)=\frac{\lambda}{R^2}.}
\end{equation}
Here \(f_{\log}(C)=0\), \(f_{\log}'(C)=1\), and
\(Rf_{\log}'-f_{\log}=P\).  This latter route is convex, but its
Hessian is nonzero at \(C\) and decays at large \(R\), rather than
approaching a positive quadratic susceptibility.
\end{proposition}
\begin{proof}
The two finite trace definitions give
\(A_p=15/2-q_H\), while the character gives
\(\log(R_H/C)/L=7-q_H\).  Subtraction proves
Eq.~\eqref{p3e:eq:v2218-finite-heat-log-value}; differentiating it
twice proves Eq.~\eqref{p3e:eq:v2218-raw-value-negative-hessian}.
For the Euler route, differentiation of \(Rf'-f=P\) gives
\(Rf''=P'=\lambda/R\).  Twice integrating with the displayed
terminal jet gives Eq.~\eqref{p3e:eq:v2218-log-euler-coefficient};
direct substitution verifies its value and Euler response.  Any other
solution differs by a multiple of \(R\), which the terminal value
removes.
\end{proof}

The parameters \(\gamma\) and \(\lambda\) in this proposition have
curvature dimension; the raw finite trace is dimensionless.  Neither
promotion is asserted to be the actual physical coefficient.  The first
fails the positive scalaron-Hessian test, and the second does not supply
the terminally degenerate, quadratically dominated interpolation sought
here.  These exclusions concern the stated direct ports.  They do not
exclude an effective action obtained with its full metric, source,
localizer and stationary-elimination contacts.  In particular the
protected octet and the fifteen-dimensional finite momentum module are
not identified by the equality of their terminal depth values.

\subsubsection{A shared variable proper-time scale in the smooth comparison}

\begin{theorem}[Complete contact residual and forced heat-clock trajectory]
\label{p3e:thm:v2218-factorized-heat-clock}
On a controlled four-dimensional minimal scalar smooth-parent heat
domain, retain the fixed protected process
\[
 V_\tau=P_0+e^{-\tau}P_W,\qquad
 q=7(1-e^{-\tau}),\qquad Z(\tau)=\Tr V_\tau=8-q.
\]
This \(q\) is the protected visibility depth; it is not a new
identification with the finite-seed parameter \(q_H\) above.
Consider the explicitly factorized comparison
\begin{equation}
 \label{p3e:eq:v2218-factorized-heat-comparison}
 S_H(t,\tau;g)=\Tr(e^{-t\Delta_g}\otimes V_\tau)
 =(8-q)\Tr e^{-t\Delta_g},\qquad t>0.
\end{equation}
Its local vacuum and Einstein coefficients, with the same trace and
minimal scalar convention, are
\begin{equation}
 \label{p3e:eq:v2218-factorized-heat-coefficients}
 c_0=\frac{8-q}{(4\pi)^2t^2},\qquad
 c_2=\frac{8-q}{6(4\pi)^2t}.
\end{equation}
These are controlled smooth-parent coefficients, not strict
small-time coefficients of the literal fixed momentum matrix.
For a positive \(C^1\) trajectory \(t=t(q)\), put
\(\mathcal U_H=8c_2\), \(\Phi=Lq\), and \(t_0=t(0)>0\).
The complete scalar action/contact residual on this comparison is
\begin{equation}
 \label{p3e:eq:v2218-factorized-contact-residual}
 \boxed{\eta_H=d\log\mathcal U_H+d\Phi
 =\left[L-\frac1{8-q}-\frac{d\log t}{dq}\right]dq.}
\end{equation}
It vanishes identically if and only if
\begin{equation}
 \label{p3e:eq:v2218-forced-heat-clock}
 \boxed{t(q)=t_0\frac{8-q}{8}\alpha^q.}
\end{equation}
On precisely this root-normalized trajectory,
\begin{align}
 \label{p3e:eq:v2218-heat-clock-coefficient-ratios}
 \frac{\mathcal U_H(q)}{\mathcal U_H(0)}&=\alpha^{-q},&
 \frac{c_0(q)}{c_0(0)}&=\frac8{8-q}\alpha^{-2q},\\
 \label{p3e:eq:v2218-heat-clock-offset-ratio}
 \frac{(c_0/c_2)(q)}{(c_0/c_2)(0)}
 &=\frac8{8-q}\alpha^{-q}.&&
\end{align}
The last ratio reaches \(8\alpha^{-7}\) at the terminal limit.
It is a positive Euclidean coefficient ratio; its value alone does
not determine a Lorentzian vacuum sign or a physical cosmological
constant.
\end{theorem}
\begin{proof}
The factors in Eq.~\eqref{p3e:eq:v2218-factorized-heat-comparison}
commute on the tensor product, and their trace is the product of their
ordinary traces.  The minimal scalar local heat density begins with
\[
 \frac{8-q}{(4\pi)^2}
 \left[t^{-2}+\frac{R_g}{6}t^{-1}+\cdots\right].
\]
This gives Eq.~\eqref{p3e:eq:v2218-factorized-heat-coefficients}
on its stated asymptotic domain.  Taking the logarithmic derivative
of \(\mathcal U_H\propto(8-q)/t(q)\) gives the full residual,
including the cutoff/proper-time derivative.  Its vanishing is the
ordinary differential equation
\[
 d\log t=\left[L-\frac1{8-q}\right]dq.
\]
Integrating from zero to \(q\), with \(t(0)=t_0\), gives
\(\log(t/t_0)=Lq+\log[(8-q)/8]\), proving necessity,
sufficiency and uniqueness of
Eq.~\eqref{p3e:eq:v2218-forced-heat-clock}.  Substituting that
trajectory into the two heat coefficients gives every displayed
ratio.
\end{proof}

\begin{corollary}[Fixed cutoff and distinct-clock tests]
\label{p3e:cor:v2218-distinct-heat-visibility-clocks}
At a fixed cutoff, the comparison residual is
\(\eta_H=[L-1/(8-q)]dq\), which cannot vanish identically on a
nontrivial depth interval.  Internal multiplicity attenuation alone
multiplies \(c_0\) and \(c_2\) by the same factor and leaves
\(c_0/c_2=6/t\) unchanged.  A variable common cutoff can instead
produce the hierarchy in
Eq.~\eqref{p3e:eq:v2218-heat-clock-offset-ratio}, but it must pass
the actual clock trajectory test above.

The successful trajectory cannot have heat proper time proportional
to visibility time over the full branch.  In fact
\begin{equation}
 \label{p3e:eq:v2218-finite-terminal-heat-time}
 t(\tau)=\frac{t_0}{8}(1+7e^{-\tau})
 \alpha^{7(1-e^{-\tau})}
 \longrightarrow\frac{t_0\alpha^7}{8}<\infty
 \quad(\tau\longrightarrow\infty).
\end{equation}
Any positive affine additive identification of heat time with
\(\tau\) is unbounded instead, and would make
\(\mathcal U_H\propto Z/t\) tend to zero.
\end{corollary}
\begin{proof}
At fixed \(t\), the derivative term in the residual vanishes, while
\(1/(8-q)\) is nonconstant.  The coefficient quotient follows
directly from Eq.~\eqref{p3e:eq:v2218-factorized-heat-coefficients}.
Substitution of \(q=7(1-e^{-\tau})\) proves
Eq.~\eqref{p3e:eq:v2218-finite-terminal-heat-time}; the final
statement follows from \(Z\to1\).
\end{proof}

\begin{remark}[Physical origin and the surviving action-port obligation]
\label{p3e:rem:v2218-fixed-origin-status}
The generated gravitational process theorem of this paper explicitly
retains the two-parameter family
\((e^{-t\Delta_g},e^{-\tau P_W})\).  Its common exponential
process-law role does not identify either the generators or their
clocks.  Paper~III-C's visibility calibration relates \(\tau\) to
an additionally calibrated physical time; it does not identify that
time with the heat proper time \(t=\Lambda_g^{-2}\).  Consequently
Eq.~\eqref{p3e:eq:v2218-forced-heat-clock} is an evaluated
necessary/sufficient target for this factorized comparison, rather
than a derived physical scale-selection law.  Setting the trajectory
because it makes the residual zero does not certify its operator
origin or the full action-density transport required by
\((\mathrm{ACD}_H)\).

Likewise the finite norm-input coefficient of Paper~III-D,
Definition~\ref{p3d:def:v2218-fixed-momentum-candidate} and
Theorem~\ref{p3d:thm:v2218-fixed-candidate-endpoints}, is an explicitly
evaluated composition of native data.  Its complete source contacts
are retained there.  Neither its useful two-regime behavior nor the
absence of a fitted intermediate profile identifies that finite word
with the physical heat action.  The actual source incidence,
normalization, density and operator/effective-action port remain to
be established.  These are the same origin obligations recorded by
the finite/smooth interface and by Paper~III-F's certified-response
domain; no variation of \(m_p\) or \(R_{\rm mom}\) is required.
Equation~\eqref{p3e:eq:v2218-factorized-heat-comparison} is not
asserted to rewrite the physical Born sum of exact finite sector
traces as a tensor-product action.
\end{remark}

\subsection{Visible response and complete gravitational ownership}
\label{p3e:sec:v227-gravity-visible}
A source-visible pole/readout can now be compared by the finite energy
and source-jet tests of Theorem~\ref{p1b:thm:v227-finite-response} and
Corollary~\ref{p1b:cor:v227-response-jets}. Once those actual tests pass,
uniqueness of a wholly uncoupled environment is unnecessary for that
readout. The same retirement is unavailable for the complete gravitational
action when it consumes that environment's heat trace or metric response.

In particular Proposition~\ref{p1b:prop:v227-central-boundary} retains
absolute block energies and their actual metric/source derivatives even
when fixed-state even dynamics has forgotten them. The 137 native
coefficient tests still apply only to an admitted native corner; actual
outer operator-valued coefficients, nonreducing defects, boundary terms,
density and trace ownership remain in the complete differentiated action.
This closes an unnecessary visible-sector uniqueness requirement without
asserting the outstanding Einstein/GHY or cosmological normalization.


\subsubsection{Laurent formula and action occurrence}

\begin{remark}[Laurent claim-owner partition]
\label{p3e:rem:laurent-claim-owner-partition}
The continuous positive multiplicative dressing law
\(D(s+t)=D(s)D(t)\), with \(D(1)=Q_\sigma\), gives
\[
 D(s)=Q_\sigma^s,\qquad
 Q_\sigma=8\sqrt{\alpha_\sigma},\qquad
 C_7=U\alpha_\sigma^{-7}.
\]
Taking logarithms gives the continuous additive Cauchy equation;
its value on the rationals and continuity give
\(\log D(s)=s\log Q_\sigma\).  This proves the formula after the
character law has been constituted.  It does not derive the law
from its endpoints.  The corresponding action obligations are:
\par\medskip
\noindent\begin{tabular}{@{}p{.38\linewidth}p{.57\linewidth}@{}}
\(\mathrm{ACT}_{2|7}^{\rm L}\) & conditional on actual EXJET7 and complete-action port\\
\(\mathrm{DCHAR}_2^{\rm coeff}\) & derived later as exponent-space character pullback (III-F)\\
\(\mathrm{JUNC}_{2|7}^{\rm str}\) & derived exact interface; physical crossing is G0/history\\
\(C_7\) scalar-remnant intercept & derived from the Clifford endpoint plus Legendre Euler data\\
\(\mathrm{TRANSPORT}_{J/\mathrm{scal}}^{EH,\rm raw}\)
  & raw coefficient transport \(U_J^{\rm raw}=96U\), derived downstream\\
\(\mathrm{G24\!\!-\!MOMCURV}\)
  & actual comparison \(\varepsilon_{K\to U}=0\) to be evaluated
\end{tabular}\par\medskip
The occurring Hurwitz branch imports
Theorem~\ref{p3c:thm:occurring-hurwitz-reg-origin} and hence
\(\alpha_\sigma=\alpha_0\).  The generic formulas retain their
unanchored scheme domain.  Neither substitution changes the momentum-curvature actualization or any actual G0 history requirement.  Paper~III-F derives the coefficient-character algebra and conditional
reconstruction. The physical placement consumes the actual source/owner
certificates; a supplied susceptibility alone does not evaluate them.
\end{remark}

The Euclidean joint-Dirac represented coefficient and the Lorentzian Laurent
endpoint remain differently typed action coefficients.  Their distinction retains separate raw coefficient owners and coupling
coordinates. The intrinsic momentum curvature is known, while its EH
comparison and the whole-action normalization remain evaluated ports.

Paper~III-C proves that exponent seven is the rank drop
to the infrared scalar remnant, this paper inherits that same hierarchy owner
as \(C_7=U\alpha_\sigma^{-7}\), and Paper~III-F proves that every inverse-Euler
representative with response \(P_{2\mid7}\) has affine terminal form
\[
 Z_7R_{\rm g}-C_7.
\]
Thus the zero-order intercept \(-C_7\) is invariant under the Euler-kernel
freedom.  On the retained slope-one representative,
\[
 f_{\rm term}(R_{\rm g})=R_{\rm g}-C_7
 =R_{\rm g}-2\Lambda_{\rm eff},\qquad
 \Lambda_{\rm eff}=C_7/2,
\]
and the \emph{formally extrapolated affine} Einstein root \(R_{\rm dS}=2C_7\), which is not a source-free root on the lower terminal curvature domain.  The represented-coefficient distinction is not a
second terminal vacuum subtraction; absolute Newton/SI comparison is
metrology/validation.


\begin{theorem}[Terminal endpoint evaluation, no-second-EH firewall, and scalaron status]
\label{p3e:thm:laur17-projected-terminal-freeze}
\label{p3e:thm:laur17-terminal-scalaron-decoupling}
The canonical positive-part projection gives \(n_7^+(C_7)=0\).  The nonlinear scalaron susceptibility reaches its terminal Hessian-degenerate
endpoint there, and the terminal role quotient gives
\begin{equation}
\label{p3e:eq:laur17-postfreeze-raw}
 \boxed{
 \operatorname{ev}_{C_7}\mathcal H_7^{\rm act}=C_7,
 \qquad
 \mathfrak f_{2\mid7}^{\rm term}(\scal_{\rm g})=\scal_{\rm g}-C_7,
 \qquad
 \frac{d^2\mathfrak f_{2\mid7}^{\rm term}}{d\scal_{\rm g}^2}=0.}
\end{equation}
The terminal sector has zero scalaron kinetic Hessian; equivalently, the formal
algebraic scalaron mass is infinite and no finite independent scalaron response
remains on that terminal Einstein quotient.  The
nonlinear action density is the fixed boundary value
\(\Rp C_7=\Rp^2\alpha_\sigma^{-7}\), equal to
\(\Rp^2/\alpha^7\) on the anchored octahedral member.  It has zero derivative with respect
to \(\scal_{\rm g}\) and is therefore not a second Einstein--Hilbert term.

This is an endpoint evaluation, not a negative-index continuation and not the
off-shell substitution
\(\Rp\scal_{\rm g}(\scal_{\rm g}/\Rp)^0=\Rp\scal_{\rm g}\).  On the retained
Legendre branch Paper~III-F gives
\(\mathfrak f_{2\mid7}^{\rm act}(C_7)=0\) and
\((\mathfrak f_{2\mid7}^{\rm act})'(C_7)=1\), so the displayed terminal
Einstein representative is the unique affine continuation with the same
terminal one-jet.  Its zero-order coefficient is therefore the Legendre completion of the
same terminal Clifford scalar-remnant owner.  No independent
\(\mathrm{TERMREN}_7\) subtraction remains on this branch.  The active-side
Hessian is finite and positive on the open branch and tends continuously to zero at the terminal point, so the handoff is \(C^2\) with no action-Hessian atom.
Paper~II-A's exact-interface closure forbids an ungenerated microscopic
smoothing/width/profile datum; a particular physical crossing remains G0
history and is constrained by Paper~III-F's source-free junction condition.
\end{theorem}

\begin{remark}[Smooth MRESP comparison and Laurent action owner]
\label{p3e:rem:laur17-versus-smooth-mresp}
The entropy-Hessian reconstruction of
Theorem~\ref{p3:thm:master-lagrangian} remains a valid smooth
\((\mathrm{MRESP}_{\min})\) realization.  That smooth branch is not used
as the action owner of the Paper~VI-A cosmological branch.  The
minimal quadratic Laurent dressing is the active physical branch; the smooth
graph is retained as an optional comparison and as a possible microscopic
model of a handoff layer only after an explicit matching theorem is supplied.
\end{remark}

\begin{corollary}[Six finite scalar hierarchy hits and one asymptotic endpoint]
\label{p3:thm:active-window-matching-rule}
On the owner-typed ordered bundle of
Theorem~\ref{p3e:thm:uniform-reg-response-bundle}, the smooth rank-trace
coordinate crosses exactly
the first six scalar hierarchy levels at finite parameter:
\[
 q_{\rm sp}(\tau_j)=j,\qquad
 \tau_j=-\log\!\left(1-\frac j7\right)<\infty,\qquad
 \scal_j=U\alpha_{\rm hier}(\sigma_{\rm fin})^{-j},
 \qquad j=1,\ldots,6.
\]
Its seventh ledger entry is reached only as the terminal limit
\[
 q_{\rm sp}(\tau)\longrightarrow7,\qquad
 \scal(\tau)\longrightarrow
 C=U\alpha_{\rm hier}(\sigma_{\rm fin})^{-7},
 \qquad \tau\longrightarrow\infty.
\]
Together these form a seven-entry ordered scale hierarchy for one smooth
coordinate.  They are not seven spectral-projector events and do not require
\((\mathrm{FLAG})\).  Independently, the canonical upper spectral filtration
has the four exact support anchors \(d=0,4,5,7\) proved in
Corollary~\ref{p3e:cor:owner-typed-active-spectral-profile}; it is stepwise,
not a seven-event chronology.  If
\((\mathrm{FLAG})\) and the required realization data are additionally
supplied, the same ledger may index an ordered rank-one refinement; this does
not turn its seventh smooth-semigroup entry into a finite crossing.
No Hessian transport, entropy coupling, terminal datum, or ultraviolet
normalization is used.
\end{corollary}

\begin{proof}
This is the finite-hit and endpoint part of
Theorem~\ref{p3e:thm:uniform-reg-response-bundle}, together with the distinct
spectral-filtration statement cited above.  The listed Hessian, entropy,
terminal, and ultraviolet response data do not occur in its antecedent or
proof.
\end{proof}

\subsubsection{Limit verification}

\begin{proposition}[Ultraviolet and terminal limits]
\label{p3:prop:uv-limit-active-density}
On the calibrated smooth branch,
\[
 \scal(0)=U=\Lambda_{\rm EH}^2,\qquad
 \lim_{\tau\to\infty}\scal(\tau)=C,
\]
and
\[
 \mathfrak f_{\rm sp,\scal\scal}(\scal)>0\text{ on }(C,U],
 \qquad
 \mathfrak f_{\rm sp,\scal\scal}(\scal)\to0^+
 \text{ as }\scal\downarrow C.
\]
In the adopted \(\kappa=2\) normalization representative the exact reduced
response retains \(\scal_{\rm resp}=C+U\) throughout; the EH calibration
does not replace this normalization by \(U\).  Other positive response
normalization representatives give the same calibrated invariant curve after
the compensating conversion-scale change.
\end{proposition}

\begin{corollary}[Terminal frozen nonlinear magnitude readout]
\label{p3:cor:frozen-constant-term-readout}
\label{p3:eq:terminal-correction-value}
Under the source endpoint specialization of
\((\mathrm{MRESP}_{\min})\),
\[
 \boxed{C_{\rm sp}(C)=-C,\qquad
        C=\Lambda_{\rm EH}^2
        \alpha_{\rm hier}(\sigma_{\rm fin})^{-7}.}
\]
This is a generated-owner terminal response endpoint value after the source
endpoint member is specified.  With the alias \(\Rp:=U\), the factor \(\Lambda_{\rm EH}^2\) equals
\(\Rp\) by definition.
\end{corollary}

\label{p3:part:scalar-threshold-handoffs}


\section{Complete conormal and boundary-Schur source realizations}
\label{paper3e:sec:v204-source-realizations}
\subsection{Complete rooted conormal realization of the conditional cap}
\label{p3e:sec:v229-conormal-cap}
The existing complete cap constructor remains conditional on its native
action and source identifications. The following result evaluates a
specific complete conormal source instead of assuming a scalar singlet
forcing. Write $\Delta=B_I-B_U>0$, $R_s=\sqrt{UC}$ and
$G=\sqrt{1+(R/R_s)^2}$. Set $\mu^2=UC/(U+C)$ and
$\gamma_N^2=1/(U+C)$, so $\mu/\gamma_N=R_s$.

\begin{theorem}[Full conormal flux realization and absolute weight]
\label{p3e:thm:v229-conormal-cap}
Use the matched frozen conormal source of
Proposition~\ref{p1b:prop:v229-affine-conormal} and the complete
homothety of Theorem~\ref{p1b:thm:v229-coupled-homothety}, with the
whole operator scaling when more than its constant block is consumed.
Let $\overline K_{0,R}=U_RK_{N,0}U_R^\dagger$ and
$K_{N,R}=G\overline K_{0,R}$ on regular positive inverse domains.
The reference pairing and all source columns are transported together.
For a real inherited boundary amplitude $b$ take
$j_R=b\delta\mathcal J_R=b\gamma_NR J_1$. Define the normalized
complete column trace by $\langle X\rangle_P=\operatorname{Tr}_P X/
\operatorname{Tr}P$ and $d_0=\langle J_0^\dagger K_{N,0}^{-1}J_0\rangle_P$.
Then
\[
 \min_{p_0+p_1=j_R}\frac12\left(
 \langle p_0,\overline K_{0,R}^{-1}p_0\rangle+
 \langle p_1,K_{N,R}^{-1}p_1\rangle\right)
 =\frac{b^2\gamma_N^2d_0R^2}{2(1+G)}.
\]
It equals the required $2\Delta R^2/(1+G)$ exactly when
\begin{equation}
 b^2\gamma_N^2d_0=4\Delta.
 \label{p3e:eq:v229-amplitude-target}
\end{equation}
\end{theorem}
\begin{proof}
The constrained stationary equations give $p_0=j_R/(1+G)$,
$p_1=Gj_R/(1+G)$. Substitution leaves
$\langle j_R,\overline K_{0,R}^{-1}j_R\rangle/[2(1+G)]$.
The complete covariance $\mathcal C=J_1J_1^\dagger=J_0J_0^\dagger$
commutes with $U_R$. Trace cyclicity therefore makes this norm
$b^2\gamma_N^2R^2d_0$, proving the formula and the weight criterion.
All source companions occur in this trace. This positive flux dual
uses the signed boundary-work parent of the existing cap constructor;
it is not the minimum of the opposite-sign Dirichlet source functional.
\end{proof}

\begin{corollary}[Evaluated amplitude target for the rooted block]
\label{p3e:cor:v229-rooted-amplitude}
For the constant rooted two-partner block let
\[
 w_+=\frac{\mu^2-\lambda_-}{\lambda_+-\lambda_-},\quad w_-=1-w_+,
 \qquad \mathcal Q_D=\sum_{\pm}w_\pm\sqrt{\lambda_\pm}
 \coth(L_0\sqrt{\lambda_\pm}).
\]
Then $d_0=\mathcal Q_D/\mu^2$ and
\[
 b_{\rm req}^{,2}=4\Delta R_s^2/\mathcal Q_D.
\]
\end{corollary}
\begin{proof}
The spectral weights are the singlet diagonal weights of $A$.
Theorem~\ref{p1b:thm:v229-conormal-response} evaluates the complete
conormal norm as the singlet $K_D$ return divided by $\mu^2$.
Substitute $\mu^2/\gamma_N^2=R_s^2$ in
Eq.~\eqref{p3e:eq:v229-amplitude-target}.
\end{proof}
The units and value of $b$ belong to the physical incoming boundary
section and density. Unit column normalization does not supply them.
Keeping the existing complementary four-port source and its baseline
then gives the already proved complete candidate
\[
 B_c=B_U+\frac{2\Delta}{1+G},\qquad
 f_c=P_E R+R^2B_c-\frac{\sigma^2}{B_c^2},\qquad \sigma=B_IR_s.
\]
This is an alternative representation of that parent, not an additional
term to be added to it. Its scalar boundary condition remains the one
of Theorem~\ref{p3e:thm:v223-exact-cap-reduction}.

\begin{proposition}[Geometric anchors alone do not select the source]
\label{p3e:prop:v229-source-counterfamily}
Let $x=R/R_s$, $s_0=C/U+U/C$ and
\[
 h(x)=\frac{x^2(x^2-1)^2}{(1+x^2)^4}
 \frac{(x^4-s_0x^2+1)^2}{(s_0+2)^2(1+x^2)^4}.
\]
For $0<\epsilon<1/4$, the changed section
$w_\epsilon=2\sqrt\Delta R\sqrt{1+\epsilon h(x)}$ on the same
common-return positive flux pair gives
\[
 B_\epsilon=B_U+\frac{2\Delta(1+\epsilon h(x))}{1+G}.
\]
It preserves the IR/UV limits, reciprocal midpoint value, and values
and first derivatives at $R=C,U$, but differs away from those anchors.
It obeys $B_U<B_\epsilon\le B_I$, so the existing positive finite
source ports also realize its complete density with the same $\sigma$.
\end{proposition}
\begin{proof}
Both factors of $h$ are nonnegative and reciprocal. The second is at
most one by the triangle inequality. The polynomial
$x^4-s_0x^2+1=(x^2-C/U)(x^2-U/C)$ gives double zeros at the two
finite endpoints; the first factor vanishes at zero, one and infinity.
For $t^2=(G-1)/(G+1)$,
\[
 \frac{[2/(1+G)]h(x)}{t^2}
 \le\frac{2(1+G)(G^2-2)^2}{G^8}\le4.
\]
Thus the extra positive coefficient is at most $4\epsilon\Delta t^2$,
which is smaller than the missing $\Delta t^2$ to $B_I$.
This proves the range and permits the same finite port construction.
The source change is nonzero on an open set, so the responses differ.
\end{proof}
This disproves selection from these anchors alone. It does not assert
that the deformed member occurs. The fixed affine conormal prescription
of Theorem~\ref{p3e:thm:v229-conormal-cap} excludes this deformation
within its own class; its gravitational use still needs its native origin.

\begin{proposition}[The retained inverse-square source cannot hide a
coefficient mismatch]
\label{p3e:prop:v229-complete-cap-injective}
At fixed $R$ and fixed nonzero $\sigma$, the map
$B\mapsto R^2B-\sigma^2/B^2$ is strictly increasing for $B>0$.
Thus equality of complete densities with these same inputs forces
equality of $B$.
\end{proposition}
\begin{proof}
Its derivative is $R^2+2\sigma^2/B^3>0$.
\end{proof}

\begin{remark}[What the conormal theorem establishes]
\label{p3e:rem:v229-native-selection}
The geometric source exists, its full response is evaluated, and the
absolute matching target is explicit. Native selection additionally
identifies the actual gravitational normal operator and its domain,
the actual curvature/source and density maps including $b$, and the
remaining differentiated action and physical measure. These are
identities of existing framework objects, not a new occurrence axiom.
The checked Paper~III-F same-chart cap/UV calculation,
Proposition~\ref{p3f:prop:v229-cap-vs-coherent-uv-hessian}, and the
restricted positive auxiliary-block test,
Remark~\ref{p3f:rem:v229-positive-schur-diagnostic}, respectively exclude
an unchanged same-chart identification and the indicated simple
fixed-coupling positive auxiliary return.  Independently, the terminal
affine extrapolation issue is recorded in
Remark~\ref{p3f:rem:v346-terminal-affine-extrapolation}.
These scoped tests do not exclude every native replacement branch; they
leave a separately proved native replacement admissible, as stated in
Remark~\ref{p3e:rem:v223-native-selection}.
\end{remark}


\begin{proposition}[One-anchor differential certificate for the complete inverse-square source]
\label{p3e:prop:v289-cap-source-ward}
Let $I\subset\mathbb R$ be a connected curvature interval, $B\in C^1(I)$
strictly positive, and let $T(R)$ be a differentiable positive-definite
self-adjoint form on one finite real source port. Let $j(R)$ be a $C^1$
complete incoming source section in that \emph{same} port, with all columns
and multiplicities retained, and put
\[
 s(R)=\langle j(R),T(R)^{-1}j(R)\rangle,\qquad
 X(R)=T(R)^{-1}j(R).
\]
Then $s(R)=\sigma^2/B(R)^2$ at every $R\in I$, for a fixed
$\sigma\ge0$, if and only if for some $R_0\in I$
\[
 B(R_0)^2s(R_0)=\sigma^2
\]
and the following scalar transport identity holds at every $R\in I$:
\begin{equation}
\label{p3e:eq:v289-source-ward}
 \boxed{2\langle j',X\rangle-\langle X,T'X\rangle
          +2\frac{B'}{B}s=0.}
\end{equation}
Under these hypotheses, the eliminated source density $-s(R)$
reproduces the complete inverse-square term and all its curvature jets
on the indicated smooth interval. The identity is an exact finite
comparison test, \emph{not} an assertion that the actual gravitational
source satisfies it.

For the specified $4\times4$ cap source of
Theorem~\ref{p3e:thm:v223-cap-source}, let
$T=I_4\otimes T_s$, $j=\mathrm{vec}(\rho I_4)$,
$\rho=R_s/2$ and $B=B_I(1-c^2t^2)$. Then
\[
 s=\rho^2\operatorname{Tr}T_s^{-1}
  =\frac{R_s^2}{(1-c^2t^2)^2},\qquad
 B^2s=(B_IR_s)^2=\sigma^2,
\]
so its transport identity holds identically. A native identification
additionally requires its actual operator, source, absolute trace and
physical density map to equal these stated ports.
\end{proposition}
\begin{proof}
Differentiate $T^{-1}T=I$ to obtain
$(T^{-1})'=-T^{-1}T'T^{-1}$. Hence
\[
 s'=2\langle j',X\rangle-\langle X,T'X\rangle,
 \qquad (B^2s)'=B^2\left(s'+2\frac{B'}B s\right).
\]
The stated identity is equivalent to $(B^2s)'=0$; connectedness
and the one-point value fix the constant to $\sigma^2$. The converse
is immediate. The explicit source identity follows from the determinant
and inverse-trace calculation in
Theorem~\ref{p3e:thm:v223-cap-source}, with all matrix
columns included. Equality on an interval implies equality of all
existing higher derivatives there by ordinary differentiation.
\end{proof}
\begin{remark}[Scope of the one-anchor test]
The coefficient $B$ and the positive source port must already be the
actual same-owner objects. The equation above tests their \emph{relative
transport}; it does not prove the native transport rule, absolute
normalization, additional tensor contributions, quantum measure, or
boundary action. Matching $B^2s$ does not by itself identify a marked
source current if a physical probe sees more than the scalar norm.
\end{remark}


\subsection{Boundary-Schur realization of the capped source: exact algebra and native scope}
\label{p3e:sec:v353-schur-boundary}
The following is an exact finite response realization of the already declared
cap.  It replaces neither the native common heat action nor its unresolved
physical source incidence.  The channel space here is a \emph{comparison}
complex spinor fibre $E=E_+\oplus E_-$ of dimensions $2+2$;
identification with the actual Dirac boundary domain, source owner, and
normalization remains a separate same-action test.  Put
\[
\begin{aligned}
 x&=R/R_s,& G&=\sqrt{1+x^2},&
 \theta&=\operatorname{arsinh}x,\\
 t&=\tanh(\theta/2)=\frac{x}{1+G},&
 u&=ct,& b&=1-u^2,\\
 c^2&=\frac{B_I-B_U}{B_I}.
\end{aligned}
\]
The positive determinant-one two-channel matrix
\[
 \mathcal B(x)=\begin{pmatrix}G&x\\x&G\end{pmatrix}
 =\exp(\theta\sigma_x)
\]
has positive square root $\exp(\theta\sigma_x/2)$, whose ratio of
its off-diagonal and diagonal entries is $t$.  Thus the existing
cap coefficient is exactly
\begin{equation}
\label{p3e:eq:v353-halfboost-cap}
 B(R)=B_I(1-c^2t^2)
     =B_U+\frac{2(B_I-B_U)}{1+G}.
\end{equation}
This is a half-boost \emph{representation of the prescribed scalar
profile}, not a claim that the ordinary Dirac propagator equals
$\mathcal B^{1/2}$.

\begin{theorem}[Positive two-channel Schur characterization]
\label{p3e:thm:v353-positive-schur}
Let $K_\partial(R)$ be a positive Hermitian source response on
$E_+\oplus E_-$, written in the fixed incoming physical comparison
frame as
\[
 K_\partial=\begin{pmatrix}A&C\\C^\dagger&D\end{pmatrix},
 \qquad D>0,\qquad S_+=A-CD^{-1}C^\dagger>0.
\]
For $P_+$ the projection onto $E_+$, and the rank-two comparison
current $J=(R_s/\sqrt2)P_+$, its normalized stationary quadratic return is
\begin{equation}
\label{p3e:eq:v353-schur-trace}
 F_\partial(R):=-\frac18\operatorname{Tr}
      (JK_\partial^{-1}JK_\partial^{-1})
 =-\frac{R_s^2}{16}\operatorname{Tr}_{E_+}S_+^{-2}.
\end{equation}
In particular $S_+=b(R)I_2$ implies
$F_\partial=-R_s^2/[8b(R)^2]=-\sigma^2/[8B(R)^2]$.
For the normalized diagonal blocks $A=D=I_2$ one has the exact
equivalence
\begin{equation}
\label{p3e:eq:v353-mixing-characterization}
 S_+=bI_2\quad\Longleftrightarrow\quad
 CC^\dagger=u^2 I_2.
\end{equation}
For $u\ne0$ this means $C=uU_\partial$ for some unitary
$U_\partial:E_-\to E_+$; at $u=0$ one has $C=0$.
\end{theorem}
\begin{proof}
The block-inverse identity gives
$P_+K_\partial^{-1}P_+=S_+^{-1}$ on $E_+$, hence the trace formula.
Insert $S_+=bI_2$ and use $\operatorname{rank}P_+=2$.
For $A=D=I_2$, $S_+=I_2-CC^\dagger$ and $b=1-u^2$,
which proves the last equivalence; polar decomposition proves the
unitary statement.  Positivity follows from the positive Schur
criterion $0<b\le1$.
\end{proof}

\begin{proposition}[Exact first curvature-jet target]
\label{p3e:prop:v353-schur-jet}
Assume $C^2$ regularity at $R=0$ and the \emph{actual justified}
reference normalization $A(0)=D(0)=I_2$, $C(0)=0$.
Then
\[
 S_+'(0)=A'(0),\qquad
 S_+''(0)=A''(0)-2C'(0)C'(0)^\dagger.
\]
Consequently native equality $S_+=bI_2$ on a neighbourhood requires
\begin{equation}
\label{p3e:eq:v353-native-jet-target}
 A'(0)=0,\qquad
 A''(0)-2C'(0)C'(0)^\dagger
 =-\frac{c^2}{2R_s^2}I_2.
\end{equation}
If an independently established native identity has $A''(0)=0$,
then $C'(0)C'(0)^\dagger=c^2I_2/(4R_s^2)$.
None of these target equalities is inferred from the cap itself.
\end{proposition}
\begin{proof}
Differentiate $S_+=A-CD^{-1}C^\dagger$ twice at the normalized
reference, using $C(0)=0$.  The half-boost expansion gives
$u(R)=cR/(2R_s)+O(R^3)$ and
$b''(0)=-c^2/(2R_s^2)$.
\end{proof}

\begin{proposition}[Scalar-source non-identifiability]
\label{p3e:prop:v353-schur-orientation}
For any continuous unitary family $U_\partial(R)$, the kernels
\[
 K_U=\begin{pmatrix}I_2&uU_\partial\\
 uU_\partial^\dagger&I_2\end{pmatrix}
\]
have the same $P_+$-supported quadratic return
$-R_s^2/(8b^2)$ and all its scalar curvature derivatives, whereas
$P_+K_U^{-1}P_-=-uU_\partial/b$ depends on the orientation.
Even without scalarity, the family
\[
 S_\varepsilon=b\,\mathrm{diag}\bigl((1+\varepsilon)^{-1/2},
                     (1-\varepsilon)^{-1/2}\bigr),\qquad
 |\varepsilon|<1,
\]
has $\operatorname{Tr}S_\varepsilon^{-2}=2/b^2$.
Hence a scalar inverse-square return cannot establish the actual
operator orientation, source transport or full quadratic response.
\end{proposition}
\begin{proof}
Direct inversion gives $P_+K_U^{-1}P_+=I_2/b$ and
the stated mixed block.  The diagonal example satisfies
$\operatorname{Tr}S_\varepsilon^{-2}
=b^{-2}[(1+\varepsilon)+(1-\varepsilon)]$.
\end{proof}

\begin{remark}[Native-source firewall and dependence]
\label{p3e:rem:v353-schur-native-firewall}
The actual Dirac boundary operator need not have $A=D=I_2$ and
its Dirichlet-to-Neumann map generally depends on the tangential
spectral label, not only on $R$.  The positive kernel above is a
\emph{finite effective comparison}.  The physical incidence test
must derive the common source projector, its transported normalization,
the internal multiplicities, quantum measure, and an actual
mode-reduced $S_+(R)$.  Matching this scalar source does not
remove the independent transverse-traceless and matter-source tests
of Paper~III-F, Proposition~\ref{p3f:prop:v289-flrw-tensor-obstruction}.
\end{remark}

\begin{remark}[Quarter-turn marker and scalar-cap source limitation]
\label{p3e:rem:v3030-quarterturn-cap-firewall}
Corollary~\ref{p1:cor:v3030-rooted-imaginary-coherence}
extracts a nonzero marked imaginary coherence from the complete
rooted source Gram, while Proposition~
\ref{p1:prop:v3030-phase-sample-nonuniqueness} shows that
quarter-turn samples alone do not fix its depth derivative.
Neither result selects the cap's physical source current.
On the positive two-channel comparison of
Proposition~\ref{p3e:prop:v353-schur-orientation}, for each
unitary \(U(R)\) one has
\[
 K_U=\begin{pmatrix}U&0\\0&I_2\end{pmatrix}
      K_{I_2}
      \begin{pmatrix}U^\dagger&0\\0&I_2\end{pmatrix},
 \qquad
 \left[\begin{pmatrix}U&0\\0&I_2\end{pmatrix},P_+\right]=0.
\]
Thus its \(P_+\)-supported scalar return and every admitted scalar
curvature jet remain independent of \(U\), even when \(U\)
varies with curvature; an unmarked cap scalar does not consume a
quarter-turn.  The nonzero marked mixed response, if physically
present, must be determined by the \emph{actual} common-action
source, measure and metric variations, subject to the conditional
signed-response selection rule of Corollary~
\ref{p3f:cor:v3030-signed-heat-twojet}. No extra term is added to
\(f_c(R)\), and no physical cosmological dark matter or RAR/lensing
claim is promoted by this comparison.
\end{remark}

\section{Whole-action Newton matching and quantum reduction}
\label{paper3e:sec:v204-whole-action}
\subsection{Whole-action variation test for Newton matching}
\label{p3e:sec:v356-whole-action-variation}

\begin{proposition}[Equivalent first-jet and action matching]
\label{p3e:prop:v356-whole-action-jet-criterion}
Let $\mathcal U$ be a connected open source/field chart on which
the raw action $S_J$ and a proposed physical projection $S_{\rm phys}$
are continuously differentiable.  For one constant $z>0$, the
following conditions are equivalent:
\begin{enumerate}
\item $dS_{\rm phys}=z\,dS_J$ in \emph{all} retained field,
metric, density, source and boundary directions on $\mathcal U$;
\item $S_{\rm phys}=z S_J+C$ on $\mathcal U$ for some constant $C$.
\end{enumerate}
Consequently all higher jets agree with the same constant factor
where they exist.  Matching only the Einstein--Hilbert coefficient
(or only one directional variation) does not imply either condition.
\end{proposition}
\begin{proof}
The second condition implies the first by differentiation.  Conversely,
$d(S_{\rm phys}-zS_J)=0$ on the connected open chart, so its value is
constant along every piecewise smooth path in that chart.  Repeated
differentiation proves the higher-jet statement on smooth subcharts.
For an explicit counterexample take two independent source coordinates
$x,y$ and $S_J=96x+y^2$, $S_{\rm phys}=x+y^2$.
The coefficient of $x$ agrees with $z=1/96$ but
$\partial_y^2S_{\rm phys}=2\ne 2/96=z\partial_y^2S_J$.
\end{proof}

\begin{remark}[Separate physical matching obligation]
\label{p3e:rem:v356-newton-owner-limits}
This proposition sharpens the inspection of
Proposition~\ref{p3e:prop:v2211-whole-action-newton-certificate}.
It does not generate the physical projection $\mathcal P$, fix its
normalization, or establish $G_{\rm phys}=G_{\rm scal}$.
The complete action, its source jets and the observed Newton
normalization must still be transported and verified on the actual
selected member; a merely formal $1/96$ coefficient is insufficient.
\end{remark}

\begin{remark}[Typed mixed-boundary test for physical Newton projection]
\label{p3e:rem:v3031-newton-boundary-transport}
The source's conditional whole-action normalization
(Proposition~\ref{p3e:prop:v2211-whole-action-newton-certificate})
also requires the actual physical boundary source/domain to be transported.
On a specified Euclidean normal-parallel Dirac collar,
Paper~III-F Proposition~\ref{p3f:prop:v3031-mixed-square-boundary-domain}
provides the scalar-dependent squared-domain Robin increment,
and Proposition~\ref{p3f:prop:v3031-mixed-scalar-boundary-heat}
calculates the ordinary-heat $a_3$ and $a_4$ scalar boundary packets.
Corollary~\ref{p3f:cor:v3031-boundary-scalar-two-jet}
exhibits a nontrivial scalar boundary source Hessian and preserves
an independent spectral/Dirichlet continuation comparison.
These tests neither construct the physical projection $\mathcal P$
nor identify its marked heat, fermion measure, Lorentzian boundary
source and quantum reduction.  They do not alter
$G_{J,\mathrm{raw}}=G_{\mathrm{scal}}/96$ or establish
$G_{\mathrm{phys}}=G_{\mathrm{scal}}$.
\end{remark}


\subsection{Quantum-reduced action normalization and Newton-response scope}
\label{p3e:sec:v358-quantum-newton}
The whole-action first-jet certificate of
Proposition~\ref{p3e:prop:v356-whole-action-jet-criterion} is a
\emph{classical} action test.  The following separate finite quantum
calculation identifies the extra measure information required when a
multiplicative action normalization is transported through an internal
integration at fixed operational $\hbar$.

\begin{proposition}[Exact finite reduction covariance test]
\label{p3e:prop:v358-finite-reduction-covariance}
Let $x$ range over a connected finite-dimensional source chart,
$\mu(dy)$ be a source-independent measure and
$Z_z(x)=\int \exp[-zS(x,y)/\hbar]\,\mu(dy)$ be strictly positive for
$z$ in an interval of positive values containing $1$.  Suppose
$S$ is twice differentiable in $x$ and all displayed differentiation
and integration interchanges are valid.  Put
$\Gamma_z=-\hbar\log Z_z$ and use $\langle\cdot\rangle_z$
for this normalized Gibbs state.  Then for retained source tangents
$u,v$,
\begin{align}
 d_u\Gamma_z&=z\langle d_uS\rangle_z,\label{p3e:eq:v358-quantum-first}\\
 d_vd_u\Gamma_z&=z\langle d_vd_uS\rangle_z
  -\frac{z^2}{\hbar}\operatorname{Cov}_z(d_uS,d_vS),
  \label{p3e:eq:v358-quantum-second}\\
 \frac1z d_u\Gamma_z-d_u\Gamma_1
  &=-\frac1\hbar\int_1^z\operatorname{Cov}_t(d_uS,S)\,dt.
  \label{p3e:eq:v358-covariance-obstruction}
\end{align}
Thus $d\Gamma_z=z\,d\Gamma_1$ on the chart is equivalent to the
vanishing of the indicated \emph{integrated} covariance in all retained
directions.  Covariance vanishing pointwise in $t$ is sufficient but
is not necessary at one chosen value of $z$.
\end{proposition}
\begin{proof}
Differentiate $\log Z_z$ under the integral to obtain the first
identity.  The derivative of a normalized expectation in direction $v$
is $d_v\langle F\rangle_z=\langle d_vF\rangle_z-
(z/\hbar)\operatorname{Cov}_z(F,d_vS)$, giving the second.
Similarly $\partial_z\langle d_uS\rangle_z=
-\operatorname{Cov}_z(d_uS,S)/\hbar$.  Integrate from $1$ to $z$;
the result is oriented also when $0<z<1$.
\end{proof}

\begin{corollary}[Gaussian determinant and measure criterion]
\label{p3e:cor:v358-gaussian-newton-measure}
For a positive $m$-dimensional real Gaussian spectator with
$S(x,y)=A(x)+\frac12y^\mathsf TK(x)y-j(x)^\mathsf Ty$,
$K(x)>0$, and a common source-dependent density $\rho(x)\,dy$
independent of $z$, put
$S_{\rm red}=A-\frac12j^\mathsf TK^{-1}j$.
Then, up to a source-independent constant,
\[
 \Gamma_z=zS_{\rm red}
  +\frac{\hbar}{2}\log\det K-\hbar\log\rho
  +\frac{m\hbar}{2}\log z.
\]
Consequently
\begin{equation}
\label{p3e:eq:v358-gaussian-jet-defect}
 d\Gamma_z-z\,d\Gamma_1
 =\hbar(1-z)\,d\log\frac{\sqrt{\det K}}{\rho}.
\end{equation}
For fixed $z\ne1$, the entire first source jet scales on a
connected chart precisely when $\rho$ is proportional to
$\sqrt{\det K}$ there.  For Lebesgue measure, $K=e^{x^2}$,
$j=0$ and $z=1/96$, first derivatives at $x=0$ coincide after
normalization but the difference of second derivatives is
$\hbar(95/96)$.
\end{corollary}
\begin{proof}
Complete the square and evaluate the positive Gaussian integral.
The $\log z$ term is source-independent; differentiation yields the
claimed defect.  The last example has
$\frac\hbar2(1-z)d^2\log K/dx^2=\hbar(1-z)$ at $x=0$.
\end{proof}

\begin{remark}[Physical projection remains an independent gate]
\label{p3e:rem:v358-newton-quantum-occurrence}
This conditional finite reduction theorem is not a derivation of a
Born--Hopf gravitational Gibbs measure.  The actual joint action,
source-dependent internal Jacobian and boundary measure, retained
metric/clock directions and metrological Newton readout must be
instantiated before Eq.~\eqref{p3e:eq:v358-covariance-obstruction}
can be applied physically.  Neither the raw factor $96$ nor the
Gaussian comparison establishes $G_{\rm phys}=G_{\rm scal}$.
\end{remark}


\section{Finite Petersen, scale-handoff, and exceptional-frame comparisons}
\label{paper3e:sec:v204-finite-comparisons}
\subsection{Petersen geometry and the scope of the gravitational action}
\label{p3e:sec:v319-petersen-gravity}
\begin{proposition}[No automatic metric, Newton, or curvature correction]
\label{p3e:prop:v319-petersen-gravity-firewall}
Integrating the derived five-core incidence packet into Paper I-B and
its two matched-factor counterpart into $M_8$ does not change the
smooth eight-dimensional host, its metric, the existing continuous
structure group, or the Dirac/heat operator that occurs in the
physical gravity action.  The 225-dimensional auxiliary reciprocal
edge tensor decomposes under Born exchange as $120\oplus105$; those
are exchange-module dimensions, not an Einstein--Hilbert heat trace
nor an additional 120-element copy of $2I$.
Therefore the established Newton coefficient, Euler response and
Weyl-tensor terms of this paper are unchanged by the mere availability
of the Petersen packet.  Any proposed curvature correction needs a
same-owner operator and a complete metric two-jet.
\end{proposition}
\begin{proof}
Paper~I-C, Theorem~\ref{p1b:thm:petersen-factor-exchange},
constructs the finite tensor on a \emph{marked} pair of existing
quaternionic factors.  No construction in Paper I supplies a new
metric or a new action term.  The physical variations therefore
remain those of the original action; an insertion would require the
provenance and tensor contacts of Paper~III-F's complete action.
\end{proof}

\begin{remark}[A possible sixfold angular diagnostic, not a gravitational source]
If an independently derived covariant source contains a tangent vector
$v$ and the uniformly weighted ten-axis directional contractions,
its quadratic and quartic scalar moments are rotationally isotropic,
whereas a sextic contraction may contain the nonzero $H_6(v)$ of
Theorem~\ref{p3b:thm:v319-petersen-moment-six}.  This finite
anisotropy does not select a spacetime tensor field, a coupling
coefficient, or a term in $f(R)$; without that sourced incidence no
physical gravitational anisotropy follows.
\end{remark}


\subsection{Independence test for the momentum-to-gravity scale handoff}
\label{p3e:sec:v346-cell-independence}
\begin{theorem}[Exact scale residual and an independent cutoff counterfamily]
\label{p3e:thm:v346-cell-scale-independence}
On the native scalar heat-coefficient branch, let
$U=\Lambda_g^2/(96\pi^2)$ and
$\mathcal K_p=6/R_{\rm mom}^2$; the Hurwitz scalarizer is
$\rho_H=3/\pi$.  Then
\[
 \boxed{\varepsilon_{K\to U}
 =\log\frac{U}{\rho_H\mathcal K_p}
 =2\log\frac{\Lambda_gR_{\rm mom}}{24\sqrt{3\pi}}.}
\]
Thus the source-compatible momentum-cell certificate
$(\mathrm{CELL}_{K\to U})$ is \emph{equivalent} to
$\Lambda_g R_{\rm mom}=24\sqrt{3\pi}$ within this heat branch.
The pre-port facts of round momentum curvature, the finite Hurwitz
representations, the Born-even scalarizer and the normalized exponential
heat-profile \emph{shape} do not force this equality: they admit a
one-parameter mathematical family with arbitrary positive scale ratio.
\end{theorem}
\begin{proof}
Substitute the two coefficient formulae into the logarithmic residual.
For fixed round momentum data, replace $\Lambda_g$ by
$s\Lambda_g$, $s>0$.  This preserves the represented momentum metric,
Hurwitz support, Clifford projectors and normalized semigroup functional
shape $e^{-u}$, while the heat-to-momentum residual changes by
$2\log s$.  The family is a counterexample to inferring the equality
from those \emph{pre-port} data alone.  It does not claim that a freely
varied cutoff is physically licensed after an actual complete action
and source normalization have been certified.  Only such a certificate
can select the required ratio.
\end{proof}
\begin{remark}[Physical owner status]
\label{p3e:rem:v346-cell-owner-status}
This theorem strengthens the independently stated
Eq.~\eqref{p3e:eq:v2211-momentum-eh-port}: a numerical coincidence
between $3/\pi$ and intrinsic $S^3$ curvature cannot substitute for a
common-owner action comparison.  The neutral $U$-owned hierarchy
$C_7/U=\alpha_\sigma^{-7}$ is independent of this additional
momentum-radius identification.  The separate raw-joint-to-scalar
coefficient ratio $U_J^{\rm raw}=96U$ likewise supplies no
whole-action physical Newton normalization.
\end{remark}

\subsection{Exceptional-triad geometric factors and a conditional capped-action branch}
\label{p3e:sec:et-conditional-branch}

\begin{proposition}[Exact exceptional-frame ratios, without physical owner assignment]
\label{p3e:prop:et-frame-factors}
On the marked Hopf horizontal plane with Levi metric $\gamma_0$,
let $G_\Gamma=\gamma_0^{-1}F_\Gamma|_H$ for the unnormalized
nearest-shell tangent frames of Paper~I-B.  Then
\[
 (G_T,G_O,G_I)=\left(\tfrac43I_H,I_H,2I_H\right),\qquad
 \delta_{\rm loc}=\frac1{6\sqrt2},
\]
and the following are exact scalars built from these operators:
\[
 \boxed{D_U=\delta_{\rm loc}^{-1}\frac{G_T}{G_O}=8\sqrt2,\quad
 D_I=\delta_{\rm loc}\frac{G_I}{G_T}=\frac{\sqrt2}{8},\quad
 D_\Lambda=\frac{G_T}{G_T}=1.}
\]
Here the fractions denote ratios of scalar endomorphisms. They obey
$D_UD_I=2$ and $D_U/D_I=64$.
\end{proposition}
\begin{proof}
Theorem~\ref{p1b:thm:v11-three-shell-frames} supplies the three
frame operators.  Substitute $\delta_{\rm loc}$ and simplify.
Equivalently these ratios follow from the protected-edge Gram
Proposition~\ref{p3c:prop:et-shell-edge-gram}.
\end{proof}

\begin{proposition}[Conditional three-regime cap; not a native-selection theorem]
\label{p3e:prop:et-conditional-cap}
\emph{Additional, currently undischarged physical hypothesis
(ET-OWNER):} the complete physical mixed source/gravity action
consumes the exceptional quadratic covariances in precisely the
UV, IR and neutral channels above, carries reciprocal local-determinant
weights and the specified absolute reflected-level normalization,
with its Born $1/8$ measure coefficient counted once.
Then on the same conditional four-port cap chart as
Definition~\ref{p3e:def:v223-cap-data}, the candidate choice
\[
 \boxed{ B_U=D_U\alpha_{\rm hier}^{1/2}=8\sqrt{2\alpha_{\rm hier}},
 \qquad B_I=D_I\alpha_{\rm hier}^{7/2}
 =\frac{\sqrt2}{8}\alpha_{\rm hier}^{7/2} }
\]
with $R_s=P\alpha_{\rm hier}^{-7/2}$,
$\sigma=B_IR_s$, and
\[
 B(R)=B_U+\frac{2(B_I-B_U)}{1+\sqrt{1+(R/R_s)^2}}
\]
produces the already declared (\emph{not additive}) cap form
\[
 \boxed{\begin{aligned}
 \mathcal I_{\rm grav}^{\rm ET}
 &=\frac18\int_M\sqrt{-g}\left[PR+B(R)R^2
 -\frac{\sigma^2}{B(R)^2}\right]d^4x\\
 &\quad+\frac14\sum_{\partial M}\epsilon\int_{\partial M}
 \sqrt{|h|}\,f_{\rm ET}'(R)K\,d^3x\\
 &\quad+\mathcal I_{\rm rest}.
 \end{aligned}}
\]
Its algebraic endpoint identities are
\[
 \boxed{B_UB_I=2\alpha_{\rm hier}^{4},\quad
 B_U/B_I=64\alpha_{\rm hier}^{-3},\quad
 \sigma/P=\sqrt2/8,\quad
 -\sigma^2/B(0)^2=-R_s^2.}
\]
Its asymptotic curvature-linear coefficient is
$P_{\rm UV}=P+2(B_I-B_U)R_s$, which is not equal to $P$.
For $\alpha_{\rm hier}>1$ of the selected REG$_0$ branch,
$B_I>B_U>0$ and $2B_U>1-B_U/B_I$; the regular-Legendre
hypothesis of Proposition~\ref{p3f:prop:v223-cap-legendre}
holds on the original declared cap domain.
\end{proposition}
\begin{proof}
Use the exact source stationary elimination of
Theorem~\ref{p3e:thm:v223-cap-source} and the same boundary term;
substitution yields the asserted conditional action with no second
Schur contribution.  The endpoint ratios follow by multiplication.
At $R=0$, $B(0)=B_I$ and $\sigma=B_IR_s$, hence the source
vacuum term is exactly $-R_s^2$.  The square-root interpolation gives
$R^2B(R)=B_UR^2+2(B_I-B_U)R_s^2
(\sqrt{1+(R/R_s)^2}-1)$, whose UV leading linear term is
$2(B_I-B_U)R_s R$.  The inequality follows directly from the
positive hierarchy and the displayed coefficients.
\end{proof}

\begin{remark}[Three independent native-owner gates; observational claims remain separate]
\label{p3e:rem:et-owner-gates}
The physical hypothesis ET-OWNER decomposes into
\emph{ET-INCIDENCE}, the full unnormalized \emph{quadratic}
exceptional-shell covariance on occurring source lanes (rather than
an antipodally cancelling coherent first moment);
\emph{ET-TRANSPORT}, reciprocal $\delta_{\rm loc}$ determinant weights
and marked metric/source projection from the actual native action;
and \emph{ET-LEVEL/ABSOLUTE}, transfer of the protected $E_{36}/E_{48}$
weights to the reflected $E_8/E_{56}$ baseline with actual source
measure, absolute Newton normalization and no double-counting.
The raw edge-square ratio gives $48\alpha_{\rm hier}^{3/4}$,
not the candidate $64\alpha_{\rm hier}^{-3}$.
A necessary \emph{noncircular} reference check would be to calculate
\[
 H_{RY}^{\rm native}(H_{YY}^{\rm native})^{-1}H_{YR}^{\rm native}
 \overset{?}{=}\frac18(1-64\alpha_{\rm hier}^{-3})
\]
\emph{from the native source action before substituting either
endpoint coefficient}; it must additionally fix $B_I$ absolutely.
The mathematical isometry of Proposition~\ref{p3f:prop:et-marked-port-isometry}
and the cap-chart Schur equality of
Proposition~\ref{p3f:prop:et-cap-source-schur}
are not an occurrence theorem.
The neutral $D_\Lambda=1$ identifies an algebraic cancellation,
not an established $2T/2T$ physical dark-energy source.
These statements retain the earlier nonselection firewall of
Remark~\ref{p3e:rem:v223-native-selection}; neither the physical
inflationary pivot clock nor a dark-matter abundance has been derived.
\end{remark}

\section{Conclusions and actual physical action obligations}
\label{paper3e:sec:organized-conclusion}

\label{p3e:sec:conclusion}
Paper~III-E has proved the choice-free gravitational heat role, its closed
trace, and the smooth-parent fixed scalar heat-kernel channels on the protected
scalar-corner boundary
\(\mathsf C_{\rm IR}^{\rm der}+({\rm GRAV\!-\!ASYM})\).
The established \(({\rm GRAV}_{\circ})\) and
\(({\rm GRAV\!-\!HEAT})\) labels are theorem-status certifications and no
longer add process occurrences.
The tags \(({\rm CORN})\) and \(({\rm BD}_{\rm EH})\) record, respectively,
the derived ambient rank-one constituent and Born-orbit diagonal pullback;
they are no longer additional boundary assumptions.  The generated law fixes
\(f_{2,g}=(4\pi)^{-2}\) and the exact Euclidean relation
\(\Lambda_g^2=96\pi^2\Lambda_{\rm EH}^2\).  Downstream closure of the native
common-heat carrier gives the exact raw transport
\(\UEHjoint^{\rm raw}=96\Lambda_{\rm EH}^2\), which is distinct from the equality
target \(\UEHjoint=\Lambda_{\rm EH}^2\).  The raw relation alone does not supply the final physical action
projection. The separate cell comparison $\varepsilon_{K\to U}=0$ is
an actual operator/action test; neutral geometry alone does not fix it.

The optional
moment--scale equivalence class survives
only on an optional enlarged-profile process root.  The construction imports the fixed
momentum seed \(\mathcal H_p\cong\mathcal H_1,\ m_p=1\) from Paper~II-B and
the premise-minimal Paper-II-C/GP2 reciprocal-role/Born closure completed in
Paper~III-A; no metric-curry or momentum selector is reopened here.  The
\((\mathrm{BR}_{\rm ind})\) and \((\mathrm{IND}_{xp})\) labels are therefore
certification records, not assumptions of the present paper.
The reciprocal group isomorphism transports normalized Haar measure exactly,
its Born exchange preserves the Haar product, and the normalized constant
scalar ground induces the same Haar probability.  These are compact geometric
measure theorems, not a prepared-vacuum assertion.  They imply no relation
for independently parameterized local effective determinant densities; no
such extension is claimed or consumed.  On the generated Born-diagonal
pure-metric image, determinant-density naturality derives
\((\mathrm{BAL}_B^{\rm eff})\) exactly and supplies every active use in this
paper.

Ambient provenance and the natural scalar-density Born pullback give
\[
 w_{\rm vac}=w_{\rm EH}=\frac18,\qquad w_{\rm nl}(U)=1.
\]
The first equality is a support/normalization theorem, not a calculation of
the observed dark-energy density.  The last equality is a full-octet
ultraviolet nonlinear weight, not by itself an inflationary solution.  At the
infrared endpoint \(w_{\rm nl}(C)=1/8\) but the minimal-response Hessian contains
\(u(C)=0\), so its nonlinear contribution vanishes there.

The retained Hopf-to-succession map, the typed Hopf inclusion, and the
faithful reciprocal isometry action generate a nonstationary clocked seed.
Its least physical doctrine is closed under isometry conjugation, time
translation, restriction, and compatible sequential composition; convex
time-sharing and local-uniform closure describe only the mathematical causal
envelope.
With the explicit calibration \((\mathrm{CAL}_t):T_t>0\), it has the sharp
ceiling \(c_{\rm op}=2\pi R_x/T_t\) and supplies the operational cone
without a separate propagation-family occurrence, with
future direction inherited through the Paper-I-A to
Paper-II-A to Paper-III-A scalar-succession chain rather than chosen again in
this paper.  The local field analysis derives the Lorentzian conformal sign.
It keeps the operational cone in \(TM\), the principal symbol in \(T^*M\),
and derives the tangent front cone by the Hamilton map;
on the primitive unlabelled undeformed principal branch it also derives
\(c_{\rm field}=c_{\rm op}\).  Later photon principal ancestry gives
\(c_{\rm light}^{\rm th}=c_{\rm op}\), while empirical channel identification
remains a separate bridge.
\((\mathrm{CAL}_{xp})\) and \((\mathrm{CAL}_t)\) assign spatial and temporal
units but are unnecessary for the scale-free cone.  Euclidean real-time conclusions retain their claim-specific Wick chronology:
finite cyclic, same-carrier, and finite-OS tags are \texttt{DERIVED-LATER}
only on the corresponding Paper-IV theorem domains, while interacting
infinite-volume reconstruction remains unresolved.  The constructor
firewall classifies these as analytic certifications: neither G0 actuality nor
GP2 same-owner invariance can manufacture the missing continuation owner or
its analytic hypotheses.

For the nonlinear branch, Paper~III-E imports Paper~III-D's exact abstract
response and its normalization equivalence.  The invariant curve is
\[
 0<C<U,\qquad
 g_{\rm eff}(\scal)=\frac{\scal^2+CU}{(C+U)\scal},
\]
and Paper~III-E adopts the convenient \(\kappa=2\) representative
\(\scal_{\rm resp}=C+U\).  This normalization is not an additional
selector.  The neutral visibility family is imported separately.  The
occurring minimal-response composite co-realizes the hierarchy-active state
\(\sigma_{\rm act}(\scal)=\rho_{\tau_{\rm H}(\scal)}\); its entropy defines
\(n_{\rm act}\), which is distinct from the geometric secant exponent
\(n_{\rm sec}\).  The active nonlinear fraction is
\((1/8)8^{n_{\rm act}}\).  Ordinary covariant tensor pullback uniquely fixes
the Hessian attenuation, and the continuous entropy-sufficient,
tensor-multiplicative octet normalization uniquely fixes this active fraction;
neither is a residual functional selector in
\((\mathrm{MRESP}_{\min})\).  The Einstein--Hilbert constituent remains fixed
at \(1/8\).  The scalar-action coefficient compiler fixes
\(U=\Lambda_{\rm EH}^2\).  Any physical Planck/Newton rewrite is separately
owner-typed: \(U=\Rp\) on a primitive curvature realization and
\(U=\Rp/96\) on the raw represented coefficient.

These statements continue to classify the optional smooth
$(\mathrm{MRESP}_{\min})$ branch.  The active Laurent input exported here is
instead the fixed-depth local susceptibility
\[
 \chi_{2\mid7}=Q_\sigma^{n_2^+}\chi_7^{\rm jet},
 \qquad \chi_7^{\rm jet}
 =\frac{(1+n_7^+)(\scal^2-C_7^2)}{(U+C_7)\scal^2},
\]
with $Q_\sigma=8\sqrt{\alpha_\sigma}$.  Paper~III-F derives its unique
terminal-normalized same-action reconstruction.  The structural Laurent
junction is also derived: the action is $C^2$ through the coefficient freeze
and has only a finite third-jet kink there, with no generated microscopic
smoothing profile or intrinsic width.  An actual interface
encounter/crossing belongs to G0 history and is constrained by the same-action
field equation on the source-free branch.  The terminal
Einstein lane \(R_g-C_7\) has zero scalaron Hessian and is the degenerate
Einstein limit of the retained active scalaron branch;
the retained endpoint magnitude \(\Rp C_7\) is not a second EH constituent.
At the depth-two coefficient freeze the full response is
\(\mathcal H_7^{\rm act}(C_2)=(C_4+C_7)/(1+\alpha_\sigma^{-7})\).  The protected-octet
reflection pairs its label with \(j=6\) only at trace/scale level.

The ordered response theorems hold uniformly over
\(\mathfrak A_{\rm REG}^{+}\), so no REG point is needed for family-level
existence, ordering, endpoint convergence, convexity, or hit counting.  A
numerical prediction still requires an evaluation
\(\mathrm{REG}_{\sigma}\); the generated octahedral interface fixes that
evaluation to \(\mathrm{REG}_0\), while unanchored branches retain the affine
family.  The certified ancestry morphism in \((\mathrm{CRP4})\) places the
occurring
response-root owner in the determinant-dual owner,
\(C=U\alpha_{\rm hier}(\sigma_{\rm fin})^{-7}\) and the intermediate
ratios are inherited from generated neutral dual-character outputs, not
matching data.  The same logarithmic depth pulls back Paper~III-C's unique
upper spectral resolution to the physical response owner.  This yields the
four exact support anchors
\[
 (U,\Pi_8),\quad
 (U\alpha_\sigma^{-4},\Pi_4),\quad
 (U\alpha_\sigma^{-5},\Pi_3(u)),\quad
 (C,P_0)
\]
without \((\mathrm{ACT}_{\rm Cl})\).  The complete orbit family is generated;
a particular \(u\), an actual transition, or any alternative profile remains
outside that conclusion.
The complete convex response is a theorem for arbitrary typed endpoint pair
\((F_C,Z_C)\) and arbitrary positive dimensionless susceptibility calibration
\(\gamma_\sigma\) on the unrestricted MRESP atlas.  The kernel of \(D^2\) is exactly two-dimensional, so two
independent raw endpoint data are mathematically necessary unless a new
physical relation or larger-action quotient is supplied.  Scalaron
decoupling needs only \(Z_C>0\), not the source value \(F_C=0\).  The source
endpoint branch uses \(\mathrm{TERM}_f^0=(0,1)\), which is not selected by GP
closure.  On the separately stated PQTI branch, the half-weight ancestry and
relative amplitude are no longer free: the III-A/III-C tensor-square
half-character, the III-B trace ratio, and the III-D unit primitive response
give \(\widehat\gamma_\sigma=16\sqrt{\alpha_\sigma}\).  Once the complete
typed input tuple jointly occurs,
the unique reconstructed response graph is co-realized by witness transport;
the action-role and solution-history sorts remain separate.  Paper~III-F later
completes the same-action inverse-Euler descent and proves that the terminal
zero-order intercept is \(-C_7\), invariant under the remaining Euler-kernel
linear normalization.  That downstream completion does not rewrite the raw
III-E reconstruction.
Its rank-trace coordinate has six
finite scale crossings and one asymptotic terminal endpoint.  These ordered
scalar levels need no \((\mathrm{FLAG})\); only an ordered rank-one projector
semantics requires that refinement.  Scalaron decoupling at
that endpoint is conditional on an independently admissible background and
does not establish a vacuum constant-curvature solution at \(C\).

Matter sourcing and common-action backreaction are addressed only after
Paper~III-F states its own common-action assumptions.
The scalar action remains the action root. Its actual coefficient and
derivatives are generated descendants; contact ancestry and the physical
quadratic response consume the evaluated $(\mathrm{ACD}_H)$ and PQTI
comparisons. The local Laurent susceptibility additionally consumes
$(\mathrm{EXJET}_7)$.  G0 alone
does not select an arbitrary gravitational profile or actual history; the
exponential law follows only after the
choice-free generated-role theorem within the displayed minimal scalar
Laplace-type signature above.  The response compiler does not follow merely
from the M8 grammar and does not replace any claim-specific
Wick/reconstruction hypothesis, unit map, or solution theorem.


\subparagraph{Exact structural/cross-owner remainder through Paper~VI-B.}
The initial action coefficient obeys $U=\Lambda_{\rm EH}^2$ on its own
frozen pairing. The actual contact/action transport must pass
$(\mathrm{ACD}_H)$; only then do CRP4 and the positive root give the
physical weight-one hierarchy. The proposed primitive susceptibility must
pass the actual normalized PQTI second-jet comparison. The shifted Laurent
excess must pass $(\mathrm{EXJET}_7)$, or the stronger interval Ward test
$(\mathrm{EXCHAR}_7)$, before its local Hessian is a physical prediction.
These tests retain actual action, density, localizer and source derivatives.
The candidate functions and their terminal-normalized integral reconstruction
remain exact conditional mathematics.


The positive exponent-space character $D_2$ and its positive-part freeze are
still algebraically derived. Given the supplied continuous susceptibility,
the reconstructed action is $C^2$ across K16 with a finite third-jet step.
An actual source-free crossing, boundary condition or microscopic transition
remains a separate complete-action/history claim. The neutral rank hierarchy,
visibility family, contact residue and anchored formal determinant value are
retained. No ungenerated smoothing profile is inserted.

The raw coefficient relation $U_J^{\rm raw}=96U$ is retained exactly, with
$G_{J,\rm raw}=G_{\rm scal}/96$. A physical equality of the two couplings
requires an actual whole-action/source projection, not gravity-only trace
renormalization. The independent momentum-to-EH comparison must evaluate
$\varepsilon_{K\to U}=0$; the intrinsic curvature
$\mathcal K_p=6/\Rmom^2$ and scalarizer $3/\pi$ alone do not evaluate it.
Absolute metrology follows the finally matched physical coefficient.


\printbibliography
\end{document}
